\RequirePackage{fix-cm}
\documentclass[11pt,a4paper]{ltjsarticle}
\usepackage[margin=24mm]{geometry}
\usepackage{amsmath,amssymb,amsthm,mathtools,booktabs,longtable,array}
\usepackage{tikz-cd}
\usepackage[unicode,hidelinks]{hyperref}
\usepackage{bookmark}
\setlength{\emergencystretch}{3em}
\setlength{\parskip}{0.22em}
\newtheorem{theorem}{定理}[section]
\newtheorem{proposition}[theorem]{命題}
\newtheorem{lemma}[theorem]{補題}
\newtheorem{corollary}[theorem]{系}
\theoremstyle{definition}
\newtheorem{definition}[theorem]{定義}
\newtheorem{assumption}[theorem]{仮定}
\newtheorem{remark}[theorem]{注意}
\newtheorem{example}[theorem]{例}
\renewcommand{\proofname}{証明}
\newcommand{\cA}{\mathcal A}\newcommand{\cB}{\mathcal B}
\newcommand{\cC}{\mathcal C}\newcommand{\cD}{\mathcal D}\newcommand{\cH}{\mathcal H}
\newcommand{\cT}{\mathcal T}\newcommand{\cF}{\mathcal F}
\newcommand{\cK}{\mathcal K}\newcommand{\cZ}{\mathcal Z}\newcommand{\OO}{\mathcal O}
\newcommand{\Db}{D^b}\newcommand{\Sfun}{\mathbb S}
\newcommand{\AR}{\mathcal A_{\mathrm R}}
\newcommand{\AV}{\mathcal A_{\mathrm V}}
\newcommand{\AAS}{\mathcal A_{\mathrm{AS}}}
\newcommand{\AAC}{\mathcal A_{\mathrm{ac}}}
\DeclareMathOperator{\Hom}{Hom}\DeclareMathOperator{\Ext}{Ext}
\DeclareMathOperator{\End}{End}\DeclareMathOperator{\RHom}{RHom}
\DeclareMathOperator{\coh}{coh}\DeclareMathOperator{\qgr}{qgr}
\DeclareMathOperator{\QGr}{QGr}\DeclareMathOperator{\gr}{gr}
\DeclareMathOperator{\Perf}{Perf}\DeclareMathOperator{\Pic}{Pic}
\DeclareMathOperator{\Ind}{Ind}\DeclareMathOperator{\thick}{thick}
\DeclareMathOperator{\rk}{rk}\DeclareMathOperator{\add}{add}
\DeclareMathOperator{\Ann}{Ann}
\title{非可換del Pezzo曲面の反標準heartと根圏\\
\large 幾何学的認識・Euler選択・アフィン有限商}
\author{}
\date{2026年11月17日\qquad 公表準備稿 hearts-v28}
\hypersetup{pdftitle={非可換del Pezzo曲面の反標準heartと根圏},pdfsubject={導来比較、HRS tilt、共鳴、反標準corner}}
\begin{document}
\maketitle
\begin{abstract}
非可換del Pezzo曲面のVdB，Rains，反標準，Artin--Schelterという四つのheartを，
構造対象と自然な比較関手を指定して比較する．genuineな滑らかな楕円gluingでは
反標準導来切断が導来同値を与え，heartの差は境界非交叉部分のHRS tiltに集中する．
VdBの右加群規約で初期平面twistを運ぶと，Rainsの辞書は$L=L_1$，$q=L_0L_1^{-1}$となる．
規格化したmutation列から得る楕円decupleについて，全27関係係数をtheta関数で表示し，
構造頂点を除いたcornerが$L_0$に依存せず，$L_1$による元の六点の可換cornerに一致することを示す．
これにより，係数比の十一条件を元の六点の共線・conic条件へ引き戻す．
さらにampleな構造対象を含むvery strong helixからnoetherianなorbit代数を再構成する．
その結果，規格化した共通領域ではAS heartと反標準heartが一致し，
noetherian性と十一因子の非消滅は独立の追加条件でなく結論となる．
一方，順序付きRains heartは実在単純根のEuler標数で選ばれ，
整数倍点族では整列による単純tiltで同じ反標準heartへ移る．
次数1--6では全根の同時共鳴点がZariski稠密であり，指定heartの一致点と不一致点は
同じ楕円chartでともに稠密となるため，一致条件は構成可能でない．
標数0・非torsionの範囲では，実在するADE根の拡大閉包と標的の有限次元heartを
それぞれDynkin preprojective加群圏と同定する．標的単純のEuler次元を$m_i$とすると，
元のアフィン代数の全有限次元表現は
$\End_{\Pi(\Gamma)}(\bigoplus_iP_i^{m_i})$という一つの有限次元商を経由する．
この商は非被約基底上の古典的な表現族も支配し，そのFrobenius性は$m_i$のNakayama不変性で決まる．
AS比較には指定されたfull very strong helixとrolled-up環の同定を入力とする．
零指数根壁については古典的な反標準収縮を記載し，標準弱楕円gluingでは
全反標準orbitに直交するdg部分圏が古典的収縮の核に一致し，gluingパラメータに依存しないことを示す．
この核を除くdg商はorbit環を保存する．coherent反標準模型との商圏比較，
壁上のheart・局所特異点・相対収縮は，証明した核の同定から分離した課題として残す．
\end{abstract}
\tableofcontents

\section{序論：反標準heartと二種類の根現象}\label{sec:intro}
滑らかな可換del Pezzo曲面では，点のblow-up，反標準環，例外列から作るhelixは
同じcoherent sheaf圏を回復する．非可換の場合には，同じ導来圏に異なる標準heartが現れ，
導来同値の存在だけでは構造対象や反標準切断の同定は決まらない．
本稿の問いは，\emph{指定された構造対象と反Serre作用がどのheartを選ぶか}である．
その答えには，非零指数の共鳴に伴うheartのtiltと，零指数壁に伴う対象の収縮という
二つの異なる現象が現れる．

構造対象を$e$，Serre関手を$\Sfun$，反Serre作用を$s=\Sfun^{-1}[2]$とし，
$R(e,s)=\bigoplus_{n\ge0}\Hom(e,s^ne)$を反標準orbit環と呼ぶ．
名称は構成の由来を表し，\emph{VdB heart}はVan den Berghのraw blow-up，
\emph{Rains heart}はgluingとblowdown表示で指定されるt構造，
\emph{AS heart}はArtin--Schelter正則$\mathbb Z$-代数，
\emph{反標準heart}は指定された$R(e,s)$の$qgr$から得る．
四名称が四つの相異なるabel圏の存在を意味するわけではない．
圏の構成，導来同値，そのt-exactness，構造対象・境界・標識の保存を分けて扱う．

\subsection{genuine領域：heartの選択と有限次元表現}
滑らかな射影有理曲面$S$の$H=-K_S$がampleである場合をgenuineと呼ぶ．
滑らかな楕円境界$E\in|H|$と$q\in\Pic^0(E)$による標準Rains gluingを固定すると，
比較の基本図式は次のようになる．
\begin{equation}\label{eq:overview}
\begin{tikzcd}[column sep=large,row sep=large]
\AV \arrow[r,"\text{同じblow-up表示}","\simeq"'] & \AR
\arrow[d,phantom,"\text{HRS tilt}" description] \\
\AAS \arrow[r,"\text{指定AS実現}","\simeq"'] & \AAC
\end{tikzcd}
\end{equation}
右側の縦の表示はabel同値ではなく，指定導来同値の下でのheartの関係である．
下側にはfull very strong helixとrolled-up同定などの入力があり，
任意のpotentialに対して存在する矢とは解釈しない．

\begin{theorem}[genuine領域の比較と表現論]\label{thm:summary}
以下はそれぞれに記した仮定の下で成り立つ．
\begin{enumerate}
\item 滑らかなgenuine $S$，滑らかな楕円境界$E$，任意の$q\in\Pic^0(E)$からの
標準gluingでは，自然反標準導来切断$\Phi$はdg同値であり，Rains heart上の振幅は$[0,1]$である．
従って
\[
 \Phi^{-1}(\AAC)=\langle\cF[1],\cT\rangle,\qquad
 \cT=\ker H^1\Phi,\quad\cF=\ker H^0\Phi
\]
となる（定理\ref{thm:derived}，\ref{thm:globaltilt}）．
\item さらに$q$が非torsionなら，$\cF$は境界非交叉の有限長Serre圏$\cZ$に含まれる．
その実在単純根は有限ADE図式$\Gamma$をなし，
\[
 \Phi\text{ がt-exact}\iff(e,s)\text{ がAZ-ample}
 \iff\chi(e,S_i)>0\quad(\text{全ての実在単純 }S_i\in\cZ)
\]
である．負指数単純の有限回のtiltが標的heartを回復する
（定理\ref{thm:globaldefect}，\ref{thm:Eulerchamber}，\ref{thm:globalADEtilts}）．
\item $\cZ$と標的の有限次元heartは，それぞれ$\operatorname{mod}_{\rm fd}\Pi(\Gamma)$とabel同値である．
$A=R/(g-1)$，標的単純の次元$m_i=\chi(e,\beta_i)>0$に対して，
全有限次元$A$加群は一つの有限次元商
\[
 A/I\simeq\End_{\Pi(\Gamma)}\left(\bigoplus_iP_i^{\oplus m_i}\right)
\]
を経由する（定理\ref{thm:ADE-abelian}，系\ref{cor:ADE-affine-quotient}）．
これは162次元の頂点除去cornerとは別の代数である．
\item $D=A/I$について，非被約基底を含む古典的な表現スキームは
$\operatorname{Rep}_n(A)\simeq\operatorname{Rep}_n(D)$となる．
表現stackは$\sum_i m_id_i=n$を満たすpreprojective表現stackの和である
（定理\ref{thm:finite-rep-stack}）．$D$は自己入射的で，
Frobeniusとなる条件は$m_i=m_{\nu(i)}$である（命題\ref{prop:quotient-Frobenius}）．
\item 同じ初期heart，中心，gluing bimodule，blow-down関手を同定したVdB表示とRains表示では，
誘導される同値はt-exactである（定理\ref{thm:RV}）．同じ最後のraw blow-upについては
section-model関手を自然反標準関手へ同定できる
（定理\ref{thm:VdBcanonical}，系\ref{cor:fourfunctors}）．
\item 三次曲面のAS helixのfoldingと反標準orbitの同定，命題\ref{prop:corner}の有限性入力の下では，
自然corner restrictionがabel同値となる必要十分条件は$\dim B/BpB<\infty$である．
規格化decupleでは全27theta係数から入力の$L_1$六点を復元でき，頂点除去cornerは$L_0$に依存しない．
指定されたfull very strong AS実現とrolled-up同定の下で，AS--反標準比較そのものがt-exactである
（系\ref{cor:ASac}，定理\ref{thm:theta-potential}，\ref{thm:original-six-points}，\ref{thm:normalized-abelian}）．
\end{enumerate}
\end{theorem}
実在する負指数根sheafは，導来切断で消えるのでなく次数1へ移る．
それが反標準AZ amplenessの失敗を検出する（系\ref{cor:negative}）．
一般の対象の総Euler値だけではtorsion pairは判定できず，
全ての非零商または部分対象のEuler値を用いる（命題\ref{prop:numericaltorsion}）．

整数倍点$p_i=N_iq$の実在配置では，単純tiltは整数列の整列として記述される．
一つの負指数根だけを持つ配置と$E_6$の混合配置を第\ref{sec:single-root}節までに具体化する．
$E_6$の普遍有限商は例によって970次元となり，そのbasic代数の156次元とも区別される．
生成元ごとの作用行列を計算したことと，忘却関手が指定する商の同型型を決めたことは異なる．

\subsection{最大共鳴と零指数壁}
固定した$E,q$について，相異なる正整数倍点は次数1--6の全根を非零指数で共鳴させる．
任意の厳密な整数順序を課してもその配置はZariski稠密である
（定理\ref{thm:maximal-density-main}）．整列定理と合わせると，
指定heartが反標準heartと一致する点も一致しない点も稠密であり，
その一致条件は構成可能でない（系\ref{cor:nonconstructible-texact}）．
これは点ごとに指定したheartの結論で，代数的な相対t構造の存在や全曲面モジュライの成分分類ではない．

一方，$H=-K_S$がnefかつbigでampleでない弱del Pezzoでは，
反標準収縮$f:S\to\overline S$の核が存在する．
古典的な$(-2)$曲線$C$の$i_*\OO_C(-1)$は全反標準orbitから見えず，
零指数での対象の消失は，genuine領域のtiltと異なる現象である．

\begin{theorem}[弱gluingのcoherent収縮と商圏の記述]\label{thm:weak-summary}
$S$を滑らかな射影有理弱del Pezzo曲面，$E\in|-K_S|$を滑らかな楕円曲線とし，
任意の$q\in\Pic^0(E)$による標準gluing $\cD_q$を取る．
\begin{enumerate}
\item $R_q$は左右noetherianで，各$X\in\cD_q$の次数付きExt
$\bigoplus_{n\ge0}\Hom(e,s_q^nX[j])$は有限生成で，非零の$j$は有限個である．
従ってcoherentな自然dg関手$\Phi_q:\cD_q\to\Db(\qgr R_q)$が存在する．
\item その全核は
\[
 \ker\Phi_q=\mathcal N_q
 =\iota_q\{G\in\Perf(S):Rf_*G=0\}
\]
である．ここで$\mathcal N_q$は全反標準orbitに直交するdg部分圏で，
$\iota_q$は共通gluing成分$\OO_S^\perp$の包含である．
\item $(\cD_q/\mathcal N_q)^\natural$は，$\OO_{
\overline S}^{\perp}\subset\Db(\coh\overline S)$にexceptional対象を
$\RHom_E(L\bar i^*(-),q)$でglueした特異gluing $\overline{\cD}_q$と同値である．
$q=\OO_E$ではこれは$\Db(\coh\overline S)$である．
\end{enumerate}
証明は定理\ref{thm:weak-orbit-kernel}，\ref{thm:weak-coherent-sections}，
\ref{thm:singular-gluing-target}に与える．
\end{theorem}
この定理は，coherent有限性を追加仮定として課すことなく標準弱gluingへ適用できる．
ただし一般の$q$について$\overline{\cD}_q\simeq\Db(\qgr R_q)$までは結論しない．
その比較は自然関手$\overline\Phi_q$の完全忠実性と本質的全射性に帰着する．
非genuineの場合，$\overline{\cD}_q$には無限個の高次数に自己Extを持つ対象があり，
source $\cD_q$のpropernessを商に移してはいけない（系\ref{cor:singular-gluing-unbounded}）．
さらに古典的特異性圏$D_{\mathrm{sg}}(\overline S)$は共通部分圏$\overline{\cB}$の商として残り
（命題\ref{prop:duval-singularity-quotient}），Du Val点$x_i$のideal sheafの像は
全$n\ge3$で高次Ext次元$4\delta_{ij}$を持ち，
gluingパラメータに依存しない（定理\ref{thm:duval-high-ext}）．
さらに特異点idealとgluingのexceptional対象の間には
$\RHom(J_{x,q},\bar e)\simeq R\Gamma(E,q)$が成立し，
非自明$q$では両側直交する一方，$q=\OO_E$では次数$0,1$に各一次元のExtが現れる
（命題\ref{prop:duval-ideal-boundary-gluing}）．これは古典的局所自己Extの$q$不変性と
例外成分へのgluingの$q$依存性を区別するものである．
また，相異なるDu Val点のideal間の導来Homは$k[-1]$である
（命題\ref{prop:duval-offdiagonal-ext}）．
非自明$q$では，これらのidealと$\bar e$が生成する有限部分圏は直和に分かれるが，
$\overline{\cD}_q$全体の分解は主張しない
（系\ref{cor:duval-local-splitting}）．
重要なことに，$J_{x,q}$の反標準$qgr$への像は実際にcoherent対象であり，
その切断加群は$H_R(t)-H_M(t)=1/(1-t)$を満たす
（定理\ref{thm:duval-point-hilbert}）．これは$qgr$側で非perfect対象を
調べるための具体的な出発点である．
全比較が完全忠実なら非自明$q$で
$\RHom_{\qgr R_q}(\mathcal I_{x,q},\pi R_q)$の全次数の消滅が必要となる
（系\ref{cor:duval-target-hom-obstruction}）．
さらに一つの$A_1$壁では$uv-w^2$の明示的な一周期matrix factorizationから
局所特異性圏のidealの安定自己準同型代数を
$M_2(k[\eta,\eta^{-1}])$（$\deg\eta=1$）と同定する
（命題\ref{prop:A1-stable-end}）．
これは普通のExtの低次数の積でも，非自明$q$の反標準$qgr$の計算でもない．
一方，反標準orbitから生成される完全部分圏
$\mathcal T_q=\thick\{s_q^ne:n\in\mathbb Z\}$上では，
弱領域でも自然関手が$\thick\{\pi R_q(n)\}$への同値を与える
（定理\ref{thm:orbit-core-equivalence}）．
非genuineならこの完全部分圏は商$\mathcal Q_q$全体を生成せず，
残余商$\mathfrak S_q$は非零である．
$q=\OO_E$では$\mathfrak S_q$は古典的な$D_{\mathrm{sg}}(\overline S)$に一致する
（定理\ref{thm:residual-singularity}）．

\subsection{先行研究と寄与の範囲}
Van den Bergh\cite{VdB}は非可換blow-upと，一般位置の下でのsection algebraへの導来比較を構成した．
Rains\cite{Rains,Filtered}はそのgluingとt構造，effective roots，数値的なnef/ample coneを解析した．
数値的amplenessと指定pairのAZ amplenessは定義から同じではない．
本稿のVdB--Rains一致は既知のpseudo-canonicalな認識の適用であり，
追加的な問題は\emph{自然な反標準関手が選ぶheart}を実在対象で決定することである．
Rainsの2025年monograph\cite{RainsBook}では，旧稿のLemma 4.9，Proposition 4.17，
Theorem 4.30，Theorem 9.10，Corollary 10.16はそれぞれ
Lemma 11.10，Proposition 11.19，Theorem 11.31，Theorem 16.10，Corollary 17.17に対応する．
単純根sheafの存在にはProposition 18.27直前の記述を用い，nef仮定付きの同命題自体を根へ誤適用しない．

Abdelgadir--Okawa--Ueda\cite{AOU}の三次曲面型構成とOkawa--Ueda\cite{OU,OUHelices}の
helix理論はAS比較の入力である．本稿では実際の右加群shiftから境界辞書
$L=L_1$，$q=L_0L_1^{-1}$を運び，構造対象を保存する規格化mutationを固定する
（第\ref{sec:boundarydata}節，定理\ref{thm:AOU-pointed}）．
AOUの印刷座標の異なるshift規約を混在させず，未変更の全座標への翻訳を確定したとも主張しない．
指定$\Psi_+$のt-exactness，actual theta係数と入力六点の一致，$L_0$非依存を
この規格化の下で証明する．elliptic Cauchy恒等式自体は既知\cite{EllipticCauchy}である．
係数トーラスの十一条件と一致領域は定理\ref{thm:AOUexactlocus}，系\ref{cor:coefficientopen}に記す．

HRS理論\cite{CHZ}，Kleinian圏のWeyl反射\cite{BridgelandKleinian}，
length heartの単純tilt\cite{BCPW}，ADE zigzagの形式性\cite{LiuWangZigzag}と
拡大閉包の復元\cite{KellerAinf}，Dynkin torsion classの分類\cite{Mizuno}は既知の機構である．
本稿ではこれらを二つの実在heartへ適用し，Euler次元から元のアフィン代数の
\emph{特定の非basicな有限商}を取り出す．自己入射性・Nakayama置換\cite{AiharaMizuno}と
matrix Hilbert公式\cite{EtingofRains}も既知入力として区別する．
GDAHAとmultiplicative preprojectiveを通じたquantized del Pezzoの先行研究\cite{EOR}は，
境界や表示を確認せずに本稿の滑らかな楕円gluingへ同一視しない．

弱領域ではRainsの一般有理曲面のgluingを出発点とし，
境界の正次数性とKeelerのnoetherian性\cite{KeelerAmple}からcoherent導来切断を構成する．
古典的な反標準収縮\cite{DolgachevWalls}に対する
Bondal--Kapranov--Schechtmanのfiber次元高々1の局所化定理\cite[Theorem 2.14]{BKSWalls}を用い，
dg商\cite{DrinfeldWalls}を特異モデル上のgluingに同定する．
これらの一般定理自体を新しい結果とはせず，固定された$(S,E,q,e,s)$の
全orbit核・coherent関手・収縮圏を接続することを本稿の結論とする．
与えられた収縮から始まる非可換収縮環\cite{DonovanWemyssWalls}は，
任意の非可換曲面上の収縮射の存在を代替するものではない．

\subsection{仮定・未解決点と第2論文の役割}
自然導来\emph{同値}とHRS比較はgenuine性を使う．実在ADE根によるheart選択と
アフィン有限商には非torsion性を加える．AS比較には指定helixとrolled-up同定を加える．
これに対しcoherent反標準\emph{関手}と全核・特異gluing表示は弱del Pezzoまで成り立ち，
$q$の位数制限を必要としない．任意の非可換曲面，特異楕円境界，相対t構造への拡張は含まない．

第2論文\texttt{dp-hearts2.tex}は，符号別Ext・混合拡大の詳細，固定dg圏のframed形式変形，
三次曲面cornerの完全Hilbert関数と162次元平坦族・厳密計算を扱う．
良いgeometric helixのcorner有限性自体は既知一般論\cite{dTdVV}でも説明できる．
本稿は自然比較に必要な有限性と六点復元を自立して証明し，細かいHilbert計算を重複させない．
零指数壁での一般$q$の$qgr$同定，指定heartでの収縮対象の位置，局所完成代数，
曲面族の底変更・降下は第\ref{subsec:coherent-future}節の課題として明記する．
これらの未解決点と，本稿で証明するpointwiseな結論を混同しない．

\section{記法，realization，反標準環}\label{sec:setup}
基礎体は$k=\mathbb C$とし，圏は$k$-線形で適合するdg enhancementを用いる．
sourceの滑らかなgluing $\cD_q$は冪等完備でproper，すなわち各RHomの全cohomologyが有限次元である．
零指数壁の商$\mathcal Q_q$や特異gluing $\overline{\cD}_q$にはpropernessを仮定しない．
有界heart$\cH\subset\cD$からのrealization
$\operatorname{real}_{\cH}:\Db(\cH)\to\cD$が同値となる場合を\emph{derived type}と呼ぶ．
heartが有界であるだけではこの同値は従わない．

Serre関手を$\Sfun$，$\theta=\Sfun[-2]$，$s=\theta^{-1}=\Sfun^{-1}[2]$と書く．
$s$を単にAS代数の次数1 shiftと同一視しない．一般の$\mathbb Z$-代数には
一周期的なshiftの指定自体がなく，helixの一周期は複数頂点を通ることがある．
derived typeのheartを$s$が両方向で保つなら
\[
\Ext^i_{\cH}(M,N)=\Hom_{\cD}(N,s^{-1}M[2-i])^*=0\quad(i>2)
\]
なので，$\cH$のhomological dimensionは高々2である．
$\OO$がexceptionalなら$\Ext^2(\OO,s^{-1}\OO)=k$であり，ちょうど2となる．

\begin{definition}[構造対象つき比較]
$(\cH,\OO,s)$の間の\emph{pointed同値}とは，構造対象の同型と$s$との自然な両立同型を
指定した同値をいう．境界restrictionやmarkingの保存は，必要な箇所で別に指定する．
\end{definition}
構造対象を$e=\OO$と書き，右加群の規約で
\[
R=R(e,s)=\bigoplus_{n\ge0}\Hom(e,s^ne),
\qquad ab=s^{\deg b}(a)\circ b
\]
と置く．局所有限右noetherian性を確認した範囲で
$\AAC=\qgr R=\gr R/(\text{有限次元加群})$と定める．
$\AR$と$\AAC$を同じ三角圏のheartと呼ぶには，まず$\Phi$を構成してtransportする必要がある．

\begin{definition}[AZ ample pair]
noetherian abel圏のpair $(e,s)$がAZ-ampleであるとは，任意の対象が
有限個の$s^{-n}e$（$n>0$）の直和の商であり，任意の全射$M\twoheadrightarrow N$に対して
$\Hom(e,s^nM)\to\Hom(e,s^nN)$が$n\gg0$で全射となることをいう\cite{AZ}．
\end{definition}
高い全Extの消滅をこの二条件だけから無条件に推論しない．
またhelixの全頂点による生成性と，一つの頂点の$s$-orbitによる生成性は別である．

\section{比較の一般論：振幅1，ample orbit，corner}\label{sec:formal}
\subsection{振幅1の同値とHRS tilt}
\begin{proposition}\label{prop:hrs}
$F:\Db(\cA)\simeq\Db(\cB)$をenhancementと両立する同値，$G=F^{-1}$とし，
$F(\cA)\subset D^{[0,1]}(\cB)$とする．このとき
\[
\cT=\ker H^1F,\qquad\cF=\ker H^0F
\]
は$\cA$上のtorsion pairであり，
\[
G(\cB)=\{X:H^{-1}_{\cA}X\in\cF,\ H^0_{\cA}X\in\cT,
\ H^i_{\cA}X=0\ (i\ne-1,0)\}=\langle\cF[1],\cT\rangle.
\]
これはderived typeのheartである．$F$がt-exactであることと$\cF=0$は同値である．
\end{proposition}
\begin{proof}
有界t構造のaisleを有限拡大とshiftで記述すると，振幅仮定と直交性から
$G(\cB)\subset D^{[-1,0]}(\cA)$を得る．
$M\in\cA$，$B_i=H^i(FM)$に対しtruncation triangleへ$G$を施すと
$GB_0\to M\to GB_1[-1]$となる．長完全列から
$H^{-1}_{\cA}(GB_0)=H^0_{\cA}(GB_1)=0$．よって
\[
0\longrightarrow GB_0\longrightarrow M\longrightarrow
H^{-1}_{\cA}(GB_1)\longrightarrow0
\]
は$\cT$と$\cF$への分解である．$T\in\cT,U\in\cF$では
$FT\in\cB$，$FU\in\cB[-1]$なので$\Hom(T,U)=0$．
逆に二つのcohomologyが表示条件を満たす$X$は$FX\in\cB$となり，heartの式を得る．
このheartのrealizationは$G$を通して標準の$\Db(\cB)$と同定されるため同値である．
$\cF=0$ならtorsion分解から$\cT=\cA$で，両heartは一致する．逆は明らか．
\end{proof}
\begin{remark}
$\cF\ne0$の対象は，$F$で消えるのでなく$H^1(FM)[-1]\ne0$になる．
抽象的なtorsion pairからHRS heartを作るだけではderived type性は従わない\cite{CHZ}．
\end{remark}

\subsection{反標準orbitによるheartの一意性}
\begin{proposition}\label{prop:unique}
$\cH\subset\cD$はderived typeの有界heart，$s$は$\cH$を両方向で保ち，
$(e,s)$はAZ-ampleとする．さらに各$M\in\cH$について，すべての$p>0$に共通の閾値で
$\Ext^p(e,s^nM)=0$ $(n\gg0)$を仮定する．このとき
\[
\cH=\{X\in\cD:\Hom(e,s^nX[i])=0\quad(i\ne0,\ n\gg0)\}.
\]
従ってこの条件を満たすheartは固定した$(\cD,e,s)$に対し高々一つである．
\end{proposition}
\begin{proof}
有限個の$H^j_{\cH}X$に高次消滅を同時に適用し，hyper-Ext spectral sequenceを退化させると
$\Hom(e,s^nX[i])=\Hom(e,s^nH^i_{\cH}X)$となる．
非零$M$について右辺がすべての十分大きい$n$で零となることは，任意に大きい$N$で
$s^NM$を負twistの有限直和から生成するAZ条件に反する．従って指定した消滅は
$H^i_{\cH}X=0$ $(i\ne0)$と同値である．
\end{proof}
\begin{corollary}\label{cor:ampleexact}
二つのheartが命題\ref{prop:unique}の条件を満たすなら，pointed導来同値はt-exactである．
同じ条件のheartの非零$s$-固定対象$M$は$\chi(e,M)>0$を満たす．
$s^rM\simeq M$なら少なくとも$\sum_{a=0}^{r-1}\chi(e,s^aM)>0$である．
\end{corollary}
\begin{proof}
第一の主張はheartの認識式による．固定対象では高次Extが消え，
$\chi(e,M)=\dim\Hom(e,M)$となる．この値が零なら全orbitからのHomが零となってAZ生成性に反する．
周期対象には$\bigoplus_{a=0}^{r-1}s^aM$を用いる．
\end{proof}

\subsection{idempotent cornerの判定}
\begin{proposition}\label{prop:corner}
$B$は局所有限右noetherian非負次数付き代数，$p\in B_0$はidempotentとする．
$R=pBp$が右noetherian，$Bp$が右$R$有限生成，$pB$が左$R$有限生成であると仮定する．
自然な$(-)p:\qgr B\to\qgr R$が同値である必要十分条件は$\dim_k B/BpB<\infty$である．
\end{proposition}
\begin{proof}
$Bp$の有限生成性から$Mp$は有限生成であり，$(-)p$はexactかつ有限次元加群を保つ．
左随伴$L(N)=N\otimes_RpB$は有限生成加群を有限生成加群に送り，
$pB$の左有限生成性から有限次元加群も有限次元に送る．
$B/BpB$が有限次元なら，counit $Mp\otimes_RpB\to M$の核・余核は$p$で消される．
それらは右noetherian性により有限生成の$B/BpB$加群なので有限次元である．
$L$を短完全列へ施した際の左端の核も同じ議論で有限次元となり，$L$は商圏上exactとなる．
$(N\otimes_RpB)p=N$とcounitが逆同値を与える．
逆に$(-)p$が同値なら$\pi(B/BpB)$は零に送られるので零であり，有限生成性から有限次元である．
\end{proof}

\section{VdB heartとRains heart：blow-up表示の同定}\label{sec:VdBRains}
\subsection{blow-downからのt構造の認識}
Rains \cite[Lemma 4.9, Proposition 4.17]{Rains}はVdB blow-up
$\alpha:\widetilde X\to X$について，$X$とblow-up中心を含む可換divisorがGorensteinなら
$R\alpha_*$が\emph{pseudo-canonical}であることを示す．本稿で必要な内容は
\begin{equation}\label{eq:pseudocanonical}
M\in D^{\ge0}(\widetilde X)
\quad\Longleftrightarrow\quad
R\alpha_*\theta_{\widetilde X}^{\ell}M\in D^{\ge0}(X)
\quad(\ell\in\mathbb Z)
\end{equation}
という非負部分の認識である．初期heartとblow-down関手およびSerre作用を固定すれば，
右辺はt構造を一意に決める．

\begin{theorem}[同じblow-upデータに対する一致]\label{thm:RV}
VdBとRainsの二つの表示に共通の初期heartを取り，有限回の各blow-upについて，
中心，gluing bimodule，blow-down関手を同定する．各段階で
\eqref{eq:pseudocanonical}の適用条件が成り立つとする．
これらの同定から得るdg同値はt-exactで，$\AV\simeq\AR$を与える．
構造対象も同じラベルの対象として同定している場合，この同値はpointedである．
\end{theorem}
\begin{proof}
一段階について，基底のt-exact同値を$E$，上のdg同値を$F$とし，
$R\alpha_{\mathrm R*}F\simeq E R\alpha_{\mathrm V*}$を用いる．
Serre関手の一意性から$F\theta_{\mathrm V}\simeq\theta_{\mathrm R}F$である．
従って\eqref{eq:pseudocanonical}の右辺は$F$で保存され，非負部分が同定される．
直交性から非正部分も同定されるので$F$はt-exact．初期heartから帰納する．
構造対象の保存はt-exactnessからでなく，指定したgluingラベルの同定から従う．
\end{proof}
\begin{remark}
Rains \cite[Theorem 4.10, Proposition 4.13, Theorem 4.30]{Rains}は
VdB blow-upのsemiorthogonal decompositionとgluingを用いる．従って同じ構成を
この二通りに記述した場合には上の仮定を満たす．しかし，独立に選んだ
blowdown構造の比較，異なるパラメータ表示の照合，途中でsection modelに置き換える操作は，
同じデータという言葉だけで済ませられない．
\end{remark}

\subsection{section modelへの移行は別の操作}
VdB \cite[\S11.1]{VdB}の図式を
\[
X_j\xleftarrow{\alpha_j}\widetilde X_j\xrightarrow{\delta_j}X_{j+1}
\]
とする．$\alpha_j$はblow-up，$\delta_j$はsection-algebra modelへの比較である．
Definition 11.1.2の一般位置条件を，同論文で固定した平面埋込みに関する
$(\sigma p_i)$へ課すと，Theorem 11.1.3より$R\delta_{j*}$と$L\delta_j^*$は逆同値で，
heart上の振幅は$[0,1]$と$[-1,0]$．coherent有界部分への制限にはLemma 9.1.5も用いる．
命題\ref{prop:hrs}から
\[
\cT_j=\ker R^1\delta_{j*},\qquad \cF_j=\ker\delta_{j*},\qquad
(L\delta_j^*)(\coh X_{j+1})=\langle\cF_j[1],\cT_j\rangle.
\]
$\cF_j=0$のとき，かつそのときに限り，この比較はabel同値である．

三次曲面では$X_7=\operatorname{Proj}F$である．VdBのChapter 12の表示$F=P/(C_3)$は，
Theorem 11.1.3の一般位置条件を仮定せずにも構成される．
従って超曲面表示の存在からその点での導来同値を結論してはいけない．
同書Remark 12.0.5も，raw blow-upとsection modelが非可換の場合に一般には一致しないことを区別する．
六段の比較の合成を，対象・関手の照合なしに一回のtiltと呼ぶこともしない．

\section{Rains heartと反標準heartの導来比較}\label{sec:Rac}
\subsection{genuine楕円gluingと既知入力}
$S$を滑らかな可換del Pezzo曲面で$-K_S$がampleなものとし，
$E\in|-K_S|$を滑らかな楕円曲線，$q\in\Pic^0(E)$とする．
$\cB=\OO_S^\perp\subset\Perf(S)$にexceptional object $e$を
\begin{equation}\label{eq:glue}
\RHom(e,M)=0,\qquad \RHom(M,e)=\RHom_E(Li^*M,q)
\quad(M\in\cB)
\end{equation}
でglueしたdg圏を$\cD_q$とする．Rainsの構成から得るheartを$\AR$とし，
$\cD_q=\Db(\AR)$とする．これは任意の抽象的gluingにheartの存在を仮定するものではなく，
\cite[Theorem 4.30, \S5]{Rains}と\cite[\S4]{Filtered}の構成を指定したものである．
その境界関手$F_q$と右随伴$i_{q*}$は
\[
F_q(e)=\OO_E,\quad F_q|_{\cB}=Li^*(-)\otimes q^{-1},\quad
\theta X\longrightarrow X\longrightarrow i_{q*}F_qX\longrightarrow\theta X[1]
\]
を満たす．$s=\theta^{-1}$はheartを保ち，$F_qs=\Psi F_q$で，
$\Psi$は境界の平行移動と次数$d=(-K_S)^2$の線束によるtensor積の合成である．

Rains \cite[Proposition 5.1]{Filtered}を$e$へ適用すると
\begin{equation}\label{eq:ring}
R=\bigoplus_{n\ge0}\Hom(e,s^ne),\quad g\in R_1\text{ は中心正則},\quad
R/gR=B(E,L,\sigma),\quad\deg L=d>0.
\end{equation}
ここで$B(E,L,\sigma)$はtwisted homogeneous coordinate ringである．
$R$は$A=R/(g-1)$のRees環であり，$\operatorname{gr}_{(g)}R=B(E,L,\sigma)[g]$．
固定次数での$g$進filtrationは有限なので，整域性と左右noetherian性は商から持ち上がる．
$\sigma$が無限位数なら\cite[Proposition 2.4]{RSS}により$\chi$条件と有限cohomological dimensionを得る．
この命題は$\deg L=1$も許す．有限位数なら$B(E,L,\sigma)$はskew Laurent環の部分環としてPIであり，
\cite[Theorems 5.13, 5.1(1), 4.5]{YZ}を順に用いると同じ有限性を得る．
ここで必要なのは$\pi R(a)$からのExtの有界性で，標的の有限大域次元を先取りしてはいない．

\begin{lemma}[orbit間のExt]\label{lem:orbit}
すべての整数$a,b$に対し，sourceの$s^ae,s^be$とtargetの$\pi R(a),\pi R(b)$の双方で
\begin{equation}\label{eq:orbit}
\Ext^j(-,-)=
\begin{cases}
R_{b-a}&j=0,\ b\ge a,\\
R_{a-b-1}^*&j=2,\ b<a,\\
0&\text{その他}
\end{cases}
\end{equation}
となる．targetの次数2のpairingは境界$H^1(E,\OO_E)$へのtraceで定まる．
\end{lemma}
\begin{proof}
sourceでは$e$のexceptionalityと境界triangleから始め，
正次数の楕円線束の$H^1$消滅で帰納する．$g$の単射性とSerre双対性が負のtwistを処理する．
targetの消滅とtrace pairingは\cite[Corollary 3.8とその証明]{Filtered}である．
Homの同定が自然な$R_n\to\Hom(\pi R,\pi R(n))$であることは，
$g$の短完全列に対する切断の列と$n=0$での$\End(\pi R)=k$から帰納して確認する．
\end{proof}

\subsection{orbit生成，trace，飽和性}
\begin{theorem}[指定した自然導来切断の同値]\label{thm:derived}
上のgenuine楕円gluingに対し
\[
\Phi=\pi\mathbf R\Gamma_*:\cD_q\xrightarrow{\sim}\Db(\qgr R),
\qquad (\mathbf R\Gamma_*X)_n=\RHom(e,s^nX)
\]
はdg同値である．$\Phi(e)=\pi R$，$\Phi s\simeq(1)\Phi$を満たし，
$\qgr R$のhomological dimensionは2である．$q$と$\sigma$の位数に制限はない．
\end{theorem}
\begin{proof}
\textit{第一段階：Ind圏での生成．}
$X\in\Ind(\cD_q)$がすべての$s^ne$に右直交するとする．
$\theta X$も同じ条件を満たすので，境界triangleと随伴性から
\[
\RHom_E(\Psi^n\OO_E,F_qX)=0\quad(n\in\mathbb Z)
\]
となる．正次数のtwisted line bundlesは$D_{\mathrm{qc}}(E)$をcompact生成するから$F_qX=0$．
さらに$\RHom(e,X)=0$より$X\in\Ind(\cB)$．対応する$G\in D_{\mathrm{qc}}(S)$は
$R\Gamma(S,G)=0$，$Li^*G=0$を満たす．$U=S\setminus E$はaffineである．
$j:U\hookrightarrow S$に対しCartier divisorの局所化から
\[
G\simeq\operatorname*{hocolim}_{n}G(nE)\simeq Rj_*j^*G.
\]
affine上の大域切断の保守性より$G=0$，従って$X=0$．これでorbitはInd圏でcompact生成し，
compact objectsは$\cD_q=\thick\{s^ne:n\in\mathbb Z\}$となる．

\textit{第二段階：関手と完全忠実性．}
Rains heartのInd圏のinjective resolutionを取り，exactな$s$の作用をstrictifyして
右$R$作用つき$\mathbf R\Gamma_*$を構成する．まず$D(\QGr R)$への関手である．
\eqref{eq:orbit}より$\Phi(s^ae)=\pi R(a)$で，高次cohomologyは有限内部次数にのみ残って$\pi$で消える．
thick生成により全対象の像は$\Db(\qgr R)$へ入る．
境界の$V\in\coh E$については$R\Gamma(E,\Psi^nV)$となり，Serre消滅により
$\Phi(i_{q*}V)$は通常の境界包含$\coh E\simeq\qgr(R/gR)\to\qgr R$の像である．
この同定は短完全列を保存する．

非零の$u\in H^1(E,\OO_E)$を$\OO_E$の自己拡大で表す．境界へ押し出し，
$e\to i_{q*}\OO_E$で引き戻し，境界triangleのconnecting mapを取ると
$\Ext^2(e,\theta e)$の非零元を得る．$e$のexceptionalityから
$H^1(E,\OO_E)\simeq\Ext^2(e,\theta e)$である．targetでも同じ議論が成り立ち，
$\Phi$は$u$を同じ$u$へ送る．従ってtraceの写像は非零．
全shiftとYoneda積に対する自然性から，orbit間の次数2の写像は次数0の同型の双対となる．
残る次数は\eqref{eq:orbit}で消えるので，orbit間のRHomはquasi-isomorphism．
両変数でthick closureを取れば$\Phi$の完全忠実性が従う．

\textit{第三段階：本質的全射性．}
$\cB$は$\Perf(S)$のadmissible部分圏でsmooth properである．gluing moduleは
pseudo-perfectであり，smoothnessからperfectとなる．従って\cite[Theorems 3.18, 3.25]{Orlov}より
$\cD_q$はsmooth properかつ飽和的である．任意の$Y\in\Db(\qgr R)$に対し
$X\mapsto\RHom(\Phi X,Y)$は，環の有限性とorbit生成からpseudo-perfect moduleとなり，
sourceの対象で表現される．こうして右随伴$\Phi^!$を得る．
counitのcone $K$はすべての$\pi R(a)$に右直交する．
標的の標準ample pair，$\chi$と有限cohomological dimensionを使い，有限個の$H^iK$へ
同時にSerre消滅を適用すると，非零cohomologyがあれば高いtwistの非零大域切断を得て矛盾する．
よって$K=0$であり$\Phi$は本質的全射．

同値でSerre関手を運ぶとtarget上で$(-1)[2]$となる．$(-1)$はheartを保つので
Extは次数2を超えて消える．$\Ext^2(\pi R,\pi R(-1))=k$からhomological dimensionはちょうど2である．
\end{proof}
\begin{remark}
この証明は共鳴により元のheartのAZ amplenessが失敗しても働く．
ただし$S\setminus E$がaffineであることを本質的に使う．
weak del Pezzo，特異楕円境界，非局所自由gluingにそのまま適用はしない．
また全$q$で同値という表現は，各$q$のsourceと対応するtargetの比較を意味する．
\end{remark}

\subsection{全heartの一回のtilt}
\begin{theorem}\label{thm:globaltilt}
定理\ref{thm:derived}の設定で$H=-K_X$がnefなRains heartを取る．
全$M\in\AR$に対し$\Phi(M)\in D^{[0,1]}(\AAC)$であり，
\[
\cT=\ker H^1\Phi,\quad\cF=\ker H^0\Phi,\quad
\Phi^{-1}(\AAC)=\langle\cF[1],\cT\rangle.
\]
\end{theorem}
\begin{proof}
Rains \cite[Theorem 9.10]{Rains}によるline bundlesの生成性を用い，
$\bigoplus_j L_j\twoheadrightarrow M$を取る．Serre双対性から
\[
\Ext^2(e,s^nM)=\Hom(M,\theta^{n+1}e)^*.
\]
このHomは$\bigoplus_j\Hom(L_j,\theta^{n+1}e)$へ単射する．
非零射があれば\cite[Corollary 10.12]{Rains}により
$(n+1)K_X-c_1(L_j)$が有効だが，その$H$との交点は
$-(n+1)d-Hc_1(L_j)<0$ $(n\gg0)$でありnef性に反する．
従って次数2のcohomologyは有限内部次数にしかなく，$\pi$で消える．
負次数と次数$>2$はheart性とhomological dimensionで消える．命題\ref{prop:hrs}を適用する．
\end{proof}
\begin{remark}
$\cT,\cF$は$H^j\Phi$によって定義される．境界外でない一般の対象について，
これを生の$\Ext^j(e,M)$の消滅と同一視してはならない．
$\Phi$のt-exactnessは正確に$\cF=0$と同値である．
t-exactならpointed性により$(e,s)$はAZ-ampleとなる．逆向きの認識には，
例えば命題\ref{prop:unique}の高次消滅条件も確認すれば十分である．
\end{remark}

\section{共鳴，数値的ampleness，AZ ampleness}\label{sec:resonance}
\subsection{共鳴の向きと実在条件}
古典的avatarの根格子$K_S^\perp$の根を$\alpha^2=-2$，$\alpha K_S=0$で定める．
$\rho(\alpha)=\OO_E(\alpha)$と書き，rank-zeroの類について
\begin{equation}\label{eq:resonance}
Li^*M=0\quad\Longrightarrow\quad \rho(\alpha)\otimes q^m\simeq\OO_E,
\qquad\alpha=c_1(M),\quad m=\chi(e,M)
\end{equation}
となる向きに$q$を取る．$q$が非torsionなら$m$は一意である．
これはrestrictionのdeterminantの条件であり，特定のheartのsheafの存在とは別である．
特に$\alpha$と$-\alpha$を入れ替えるだけで，実在sheafの指数を自由に反転することはできない．

Rainsの数値的ample coneは有効divisor錐の双対の内部である\cite[\S4]{Filtered}．
有効根$\alpha$があれば$H\alpha=0$なので$H$はこの意味ではampleでない．
以下では，そのことと指定pair $(e,s)$のAZ amplenessの失敗を別に証明する．

\begin{theorem}[既約根sheafの符号別Ext]\label{thm:sign}
境界非交叉の既約純一次元root sheaf $M$が実在するとする．$m=\chi(e,M)$なら全$n\in\mathbb Z$で
\[
\dim\Hom(e,s^nM)=\max(m,0),\quad
\dim\Ext^1(e,s^nM)=\max(-m,0),\quad\Ext^2(e,s^nM)=0.
\]
$M$は2-sphericalであり，自然導来切断の像は
\[
\Phi(M)\simeq
\begin{cases}
V_m & m>0,\\ 0&m=0,\\V_{-m}[-1]&m<0,
\end{cases}
\]
ただし$V_{|m|}$は境界外の$|m|$次元$A=R/(g-1)$加群に対応する対象である．
\end{theorem}
\begin{proof}
境界triangleから$\theta M\to M$は同型なので$s^nM\simeq M$．
line bundleから一次元以下のsheafへのExt$^2$は\cite[Corollary 7.23]{Rains}で消える．
$m\le0$で非零$e\to M$があれば，像$I$は$c_1(I)=\alpha$を持ち，余核$T$は零次元．
$I$は$e$の商なので\cite[Corollary 10.16]{Rains}より
$\chi(e,I)\ge-\alpha(\alpha+K_X)/2=1$．零次元$T$のEuler値は非負なので$m\ge1$となり矛盾する．
従ってこの場合$h^0=0$，$h^1=-m$である．

$m>0$ではRainsの反変導来双対を使う．純一次元sheafの双対は
$\mathbf R\mathrm{ad}(M)=N[-1]$の形で，$N$も双対曲面の既約純一次元境界外根sheafである
\cite[Corollaries 5.9, 7.26, 7.30, 7.31]{Rains}．双対構造line bundleを$e^{\mathrm{ad}}$とすると
\[
\Ext^j(e^{\mathrm{ad}},N)\simeq\Ext^{j+1}(M,e)
\simeq\Ext^{1-j}(e,M)^*\quad(j=0,1).
\]
従って$\chi(e^{\mathrm{ad}},N)=-m$．前段の議論を双対側の指定line bundleへ適用して
$h^1(e,M)=0$，$h^0(e,M)=m$を得る．

既約性と純粋性から非零自己射は単射で，同じ数値類の間の単射は同型である
\cite[Corollary 7.7]{Rains}．従って$\End(M)=k$．Serre双対性と
$\chi(M,M)=-\alpha^2=2$から自己Extは次数0,2で$k$，その他で零．
最後に$g$は$M$上で可逆に作用するから，非零の次数付きExt加群は
$k[g]$上有限自由で，局所化の次数0部分が所要の有限次元$A$加群となる．
\end{proof}

\begin{corollary}\label{cor:negative}
$m<0$の上記$M$が存在すれば$(e,s)$はRains heartでAZ-ampleでなく，$\Phi$はt-exactでない．
$\Phi$でtransportした反標準heartはRains heartと一致せず，両者の標準pointed abel圏の同値もない．
定理\ref{thm:derived}の範囲では，それでも導来同値は成立する．
\end{corollary}
\begin{proof}
$\Hom(e,s^nM)=0$が全$n$で成り立つためAZ生成性に反する．
$\Phi(M)=V[-1]$は非零であり，標準heartに属さない．pointed abel同値はample pairを保存する．
最後の主張は定理\ref{thm:derived}による．
\end{proof}
\begin{remark}[三つの非同値な条件]
数値的共鳴，実在する根sheaf，AZ amplenessの失敗は同義語でない．
正指数の個々の根sheafは$\cT$に入り，それだけではAZ失敗を検出しない．
負指数の根sheafが見つからないことだけから$\cF=0$とも結論しない．
またheartの非一致は指定した同値の下での主張であり，構造対象を忘れた任意のabel同値の不存在ではない．
\end{remark}

\subsection{稠密性と零指数の境界}
\begin{proposition}\label{prop:density}
滑らかな平面三次$E$，非torsion $q$と，一般の五点$p_1,p_2,p_4,p_5,p_6\in E$を固定する．
$L=\OO_E(1)$とし，$p_3^{\varepsilon}(n)$を
\[
\OO_E(p_3^{\varepsilon}(n))=L(-p_1-p_2)\otimes q^{-\varepsilon n},
\qquad n\ge1,\quad\varepsilon=\pm1
\]
で定める．各符号で有限個を除き六点は一般位置であり，
$h-e_1-e_2-e_3$の指数は$m=-\varepsilon n$となる．
各符号の点列は$E$にZariski稠密である．既約sheaf存在結果が適用される範囲で，
負指数の列はAZ失敗の稠密な例を与える．
\end{proposition}
\begin{proof}
$\Pic^1(E)\simeq E$より点は一意で，$q$の無限位数性から相異なる．
固定五点に対し，三点共線または六点conicの条件を生む追加点は有限個である．
残る無限集合は曲線上稠密で，restrictionの計算が指数を与える．
\end{proof}
\begin{proposition}\label{prop:zeroroot}
$S$がgenuineで$E\in|-K_S|$が滑らかなら，非零根$\alpha$について
$\OO_E(\alpha)\not\simeq\OO_E$である．
\end{proposition}
\begin{proof}
$H=-K_S$のamplenessから$H^0(\pm\alpha)=0$．
$\chi(\alpha-H)=0$，$H^0(\alpha-H)=H^2(\alpha-H)=0$なので$H^1(\alpha-H)=0$．
$0\to\OO_S(\alpha-H)\to\OO_S(\alpha)\to\OO_E(\alpha)\to0$から
$H^0(E,\OO_E(\alpha))=0$となる．
\end{proof}
弱del Pezzoの$(-2)$曲線$D$では，対照的に$i_*\OO_D(-1)$が全反標準orbitに直交し，
反標準収縮で導来的に消える．これは$m=0$の現象であり，非零負指数のshiftと区別する．
実際，$H|_D=\OO_D$なので全$n$で$R\Gamma(D,\OO_D(-1)\otimes\OO_D(nH))=0$である．


\section{楕円境界データの復元と規約の整合性}\label{sec:boundarydata}
独立な構成のパラメータを同定するには，pointedな導来圏だけでなく，
境界関手の\emph{実際の像}とその上の反Serre作用を揃える必要がある．
Rains \cite[(4.11)--(4.13)]{Filtered}の復元を境界作用まで明示し，
VdBの右加群規約から運んだ実際のcollectionに適用する．
一般的復元原理自体は既知であり，本節の要点は，平面twist・blow-up・section modelを
一つの境界同定で追うことで具体的な辞書を確定することにある．

\subsection{rank・Chern類・Euler標数によるrestrictionの式}
$S=\operatorname{Bl}_{p_1,\ldots,p_r}\mathbb P^2$，$d=9-r>0$，
$H=-K_S=3h-\sum_i e_i$とし，滑らかな境界$E$への古典的restrictionを
$\rho:\Pic S\to\Pic E$と書く．$L=\rho(h)$，$q\in\Pic^0(E)$とし，
Rainsのgluingにおける境界関手を$F_q$とする．原点$o\in E$を選び，
$a\in E$と$P_a=\OO_E(a-o)\in\Pic^0(E)$を同一視する．
平行移動を$t_a(x)=x+a$と書くので，線束$M$について
\begin{equation}\label{eq:translation-convention}
 t_a^*M\simeq M\otimes P_a^{-\deg M},\qquad
 t_{a*}M\simeq M\otimes P_a^{\deg M}
\end{equation}
である．この二つを入れ替えてはならない．

\begin{proposition}[restrictionと反Serre境界作用]\label{prop:boundarynormal}
$x\in K_0(S)\simeq K_0(\cD_q)$のrank，第一Chern類，Euler標数を
$u,D,c=\chi(e,x)$と書く．このとき
\begin{equation}\label{eq:detrestriction}
 \operatorname{rank}F_q(x)=u,\qquad
 \det F_q(x)=\rho(D)\otimes q^{c-u}.
\end{equation}
境界上で$F_qs=\Psi_EF_q$と書き，$\Psi_E$が平行移動と線束のtensor積からなる
Rainsの自己同値であるとする．$q=P_a$なら
\begin{equation}\label{eq:boundarynormal}
 \boxed{\Psi_E=t_{a*}(-\otimes\rho(H)),\qquad
 \Psi_E(\OO_E)=\rho(H)\otimes q^d.}
\end{equation}
これらは線束・自己同値の同型類の式であり，すべての線束同型の標準的な選択を主張しない．
\end{proposition}
\begin{proof}
$K_0(S)$の基底を$[\OO_{e_i}(-1)]$，$[\OO(-2h)]$，$[\OO(-h)]$，$[\OO]$とする．
それらのrankとEuler標数は順に$(0,0),(1,0),(1,0),(1,1)$である．
Rainsの境界像は，順に$k_{p_i}$，$L^{-2}q^{-1}$，$L^{-1}q^{-1}$，$\OO_E$であるから，
全基底について\eqref{eq:detrestriction}が成り立ち，加法性により一般の類にも従う．

Serre作用の数値式は可換側と同じなので，$s$は
\[
 (u,D,c)\longmapsto(u,D+uH,c+D\cdot H+ud)
\]
と作用する．$\delta=D\cdot H$と置けば，$\det F_q(sx)$は
$\det F_q(x)\otimes\rho(H)^u\otimes q^{\delta+ud}$となる．
これは\eqref{eq:translation-convention}による\eqref{eq:boundarynormal}の作用と一致する．
数値式だけで任意のdg関手を同定しているわけではない．ここでは既知の
平行移動・tensor積型の中で作用を決めている．実際，
$s[\OO_{e_i}(-1)]=[\OO_{e_i}(-1)]+[\mathrm{pt}]$から
$\Psi_E(k_{p_i})=k_{p_i+a}$となり，平行移動が$a$と決まる．
また$\det F_q(se)=\rho(H)q^d$がtensor因子を決める．
\end{proof}

\subsection{実際のrestriction tupleによる辞書}
\begin{proposition}[境界像からのpointed復元]\label{prop:actualdictionary}
境界関手を備えたdg圏にfull exceptional collection
$(T_1,\ldots,T_r,B_2,B_1,e)$があるとする．
その\emph{実際の}境界像が$(k_{p_1},\ldots,k_{p_r},U_2,U_1,\OO_E)$であるとする．
$\deg U_2=-6$，$\deg U_1=-3$とし，forward pairのRHomを境界restrictionが
quasi-isomorphismに写すと仮定する．このとき候補パラメータは一意に
\begin{equation}\label{eq:actualdictionary}
 q=U_2U_1^{-2},\qquad L=U_1U_2^{-1}
\end{equation}
である．$L$による埋込みと$p_i$のblow-upにRainsの構成を適用できれば，
そのordered directed dg模型は元の圏とpointedに同値となる．
この主張は指定heartのt-exactな一致を含まない．
\end{proposition}
\begin{proof}
$U_2=L^{-2}q^{-1}$，$U_1=L^{-1}q^{-1}$を解けば\eqref{eq:actualdictionary}を得る．
各forward RHomとその合成を境界のdg模型へ写すと，後向き射を零とし対角を$k$にした
同じtriangular dg代数が得られる．これは\cite[(4.11)]{Filtered}の模型である．
full exceptional collectionは各圏をそのperfect module圏に同定するので結論を得る．
個々の$K_0$類の一致だけを用いず，実際のrestriction関手が保存する合成を使っている．
\end{proof}

$c_1(T_i)=e_i$，$D_E=\sum_i p_i$とすると，この辞書には次の必要条件が伴う：
\begin{equation}\label{eq:dictionarycompatibility}
 \Psi_E(k_{p_i})=k_{p_i+a},\qquad
 \Psi_E(\OO_E)=U_2^{d-3}U_1^{3-2d}\otimes\OO_E(-D_E),
 \qquad P_a=U_2U_1^{-2}.
\end{equation}
特に三次曲面では第二の式は$U_1^{-3}(-D_E)$である．
これはEuler行列からも独立に検査できる．Orlov順序のEuler行列を$C$とすれば，
列ベクトルへの反Serre作用は$(C^{\mathsf t})^{-1}C$であり，三次曲面では
\begin{equation}\label{eq:point-class-audit}
 [\mathrm{pt}]=[B_2]-2[B_1]+[e],\quad
 [sT_i]=[T_i]+[\mathrm{pt}],\quad
 [se]=4[e]-3[B_1]-\sum_i[T_i].
\end{equation}
従ってこの検査は，gluingの公式の単なる再代入ではない．

\subsection{文献の表示を共通座標へ移す際の障害}\label{subsec:boundaryaudit}
ここではAOU v1の(9.2)，(9.3)，(9.25)，(9.47)を，\emph{同じ曲線上の同じ記号と同じrestriction同定}
で同時に読むことができるかを検査する．原典の全構成や主定理を否定するものではなく，
本稿で独立な表示を接続する際の，明確な整合性条件を与える．
$\sigma=t_b$，$L_1=\sigma^*L_0$，$\deg L_0=3$，$Q=L_1L_0^{-1}=P_{-3b}$とする．
印刷されたrestriction tupleは
\begin{equation}\label{eq:printedtuple}
 U_2=L_0^{-5}L_1^3=L_0^{-2}Q^3,\qquad
 U_1=L_0^{-2}L_1=L_0^{-1}Q
\end{equation}
である．これから得る$q_{\rm res}=Q$と$L_{\rm res}=L_0^3L_1^{-2}$に対し，
\eqref{eq:dictionarycompatibility}は境界の平行移動を$t_{-3b}$とし，
構造対象の像を$L_0^3Q^{-3}(-D_E)$とする．
一方，印刷された境界shiftを文字通り
\begin{equation}\label{eq:printedshift}
 \Psi_{\rm lit}=\sigma_*^3(-\otimes N),\qquad
 N=L_0\otimes\sigma^*L_0\otimes(\sigma^2)^*L_0\otimes\OO_E(-D_E)
\end{equation}
と読むと，$\Psi_{\rm lit}(k_p)=k_{p+3b}$，
$\Psi_{\rm lit}(\OO_E)=L_0^3(-D_E)$となる．
従って二つの表示の同時成立には
\begin{equation}\label{eq:normalization-obstruction}
 Q^2\simeq\OO_E,\qquad Q^3\simeq\OO_E
\end{equation}
が必要であり，結局$Q\simeq\OO_E$となる．$\sigma$が無限位数ならこれは不可能である．
この不整合は形式的な整数指数でなく，\eqref{eq:point-class-audit}による
点支持と構造対象の像の二つの独立な検査で検出される．

\begin{remark}[比較する規約を固定する]
この計算は二つの印刷式を同じ境界同定で併用することを排除するのであって，
AOUの主定理全体を否定するものではない．以下ではVdBの右加群規約を定義から固定し，
初期の$\OO_X(-1),\OO_X(-2)$を実際のfunctorで運んでrestrictionを再計算する．
この規格化した構成から作るAS代数と同値には下添字$+$を付ける．
印刷されたdecupleを同じパラメータでそのまま使ったAS代数や同値との同一視は主張しない．
したがって，以下の結論を著者により確認された訂正表として引用してはならない．
\end{remark}

\subsection{境界作用を先に固定した復元と九つの候補}
\begin{proposition}[boundary-first復元]\label{prop:boundaryfirst}
実際の境界作用を$\Psi_E=t_{a*}(-\otimes N)$と固定し，
$T_i$のrestrictionを$k_{p_i}$，$e$のrestrictionを$\OO_E$とする．
上のRains型gluingと両立する候補は
\[
 q=P_a,\qquad L^{\otimes3}\simeq N(D_E)
\]
である．$N$の次数が$9-r$なら，$L$は次数3の線束であり，その選択は$E[3]$の九通り．
$L$を選ぶと$U_1=L^{-1}q^{-1}$，$U_2=L^{-2}q^{-1}$が決まる．
これらから実際のblow-upとgluingを作れる場合には，復元された境界作用は指定した$\Psi_E$に一致する．
\end{proposition}
\begin{proof}
命題\ref{prop:boundarynormal}から$q=P_a$，$N=\rho(H)=L^3(-D_E)$が必要である．
逆にこの式を満たす$L$を選び，対応するblow-up上のRains gluingを作れば同命題が十分性を与える．
$\Pic^3(E)\to\Pic^9(E)$，$L\mapsto L^3$は複素楕円曲線上の$[3]$のtorsorであり，
各fiberは$E[3]\simeq(\mathbb Z/3)^2$である．同じ境界作用だけで九つのgluing圏やheartが
互いに同値になるとは主張しない．
\end{proof}

たとえば\eqref{eq:printedshift}を境界作用として固定し，点のラベルも固定するなら，
$N=L_1^3(-D_E)$であるから，これと両立する候補は必ず
\begin{equation}\label{eq:boundaryfirst-AOU}
 \begin{split}
 q&=L_0L_1^{-1},\qquad L=L_1\otimes\varepsilon,\qquad\varepsilon\in E[3],\\
 U_1&=L_0^{-1}\varepsilon^{-1},\qquad
 U_2=(L_0L_1)^{-1}\varepsilon^{-2}
 \end{split}
\end{equation}
となる．境界作用だけからは九候補を区別できない．
次の計算では初期の平面twistをさらに固定することによって$\varepsilon=\OO_E$を選ぶ．


\subsection{右加群規約で運んだOrlov collection}
VdBの量子平面$X=\operatorname{Proj}A(E,\sigma,L_0)$を取る．
右加群を使い，境界を$E$へ同定する規約は
\cite[(5.4), \S10冒頭]{VdB}に固定する．すなわち境界上の次数1 shiftは
\begin{equation}\label{eq:right-plane-shift}
 S_+(M)=\sigma_*(M\otimes L_0),\qquad
 S_+^{-1}(M)=\sigma^*M\otimes L_0^{-1}.
\end{equation}
特に$S_+$を$M\mapsto L_0\otimes\sigma^*M$へ置き換えてはならない．

\begin{lemma}[負の平面twistの実際の境界像]\label{lem:plane-negative-twists}
上の規約で$L_1=\sigma^*L_0$とすると，
\[
 Li^*\OO_X(-1)=L_0^{-1},\qquad
 Li^*\OO_X(-2)=(L_0L_1)^{-1}.
\]
\end{lemma}
\begin{proof}
正のshiftはinvertible bimodule $(L_0)_\sigma$によるtensorである．
その右加群への作用は\eqref{eq:right-plane-shift}であるから，
$Li^*\OO_X(-k)=S_+^{-k}(\OO_E)$を$k=1,2$について計算すればよい．
これは$K_0$だけでなく線束そのものの同型である．
\end{proof}

\begin{proposition}[反復比較を通るrestrictionの保存]\label{prop:transported-tuple}
VdB \S11.1の各段階$X_j\xleftarrow{\alpha_j}\widetilde X_j
\xrightarrow{\delta_j}X_{j+1}$が同論文の導来比較の仮定を満たすとする．
境界の同定をblow-up側の$\beta_j$とsection側の$\delta_{j,E}$で運び，
$G_j=R\delta_{j*}L\alpha_j^*$と置く．
平面の標準三対象を$G_{j}$の合成で運び，例外対象をその後の同じ合成で運んだ
Orlov collectionを$(T_1,\ldots,T_r,B_2,B_1,e_F)$と書く．
このとき，その実際の境界像は
\begin{equation}\label{eq:right-actual-tuple}
 (k_{p_1},\ldots,k_{p_r},(L_0L_1)^{-1},L_0^{-1},\OO_E)
\end{equation}
である．例外対象の順序は同じ表示で固定し，対応する点を同じ順序で記す．
\end{proposition}
\begin{proof}
VdB Theorem 6.3.1(1),(4)より$\beta_j$は滑らかな曲線の一点でのblow-up，従って同型で，
$R\alpha_{j*}j_*=i_*R\beta_{j*}$が成り立つ．左随伴の一意性から
$Lj^*L\alpha_j^*\simeq L\beta_j^*Li^*$を得る．
Proposition 9.1.8，式(9.11)は
$Li_{j+1}^*R\delta_{j*}\simeq R\delta_{j,E*}Lj^*$を与える．
これらの実際の境界同型を使えば$G_j$はrestrictionを保存する．
構造対象の保存は$L\alpha_j^*\OO_{X_j}=\OO_{\widetilde X_j}$と
Proposition 9.2.3(1)による．
例外対象のrestrictionが$k_{p_j}$であることは，同じ境界同定に関する
VdBの例外対象の列，または\cite[(9.24)--(9.27)]{AOU}による．
全てを後の段階へ運び，補題\ref{lem:plane-negative-twists}を使う．
\end{proof}

\begin{theorem}[共通規約での具体的なgluing辞書]\label{thm:right-dictionary}
$r=6$とし，$\sigma p_1,\ldots,\sigma p_6$が$L_0$による平面埋込みで一般位置にあるとする．
命題\ref{prop:transported-tuple}のordered collectionを備えた
$\Db(\qgr F)$に対応するRainsのパラメータは
\begin{equation}\label{eq:right-dictionary}
 \boxed{L=L_1,\qquad q=L_0L_1^{-1},\qquad
 S=\operatorname{Bl}_{p_1,\ldots,p_6}\mathbb P(H^0(E,L_1)^*) .}
\end{equation}
$S$はgenuineな滑らかな三次曲面であり，Rains gluingの圏$\cD_q$と
$\Db(\qgr F)$は，構造対象と境界関手を保つdg同値で結ばれる．
境界の反Serre作用は双方で
\[
 \sigma_*^3\bigl(-\otimes L_1^3(-D_E)\bigr)
\]
となる．境界作用から残っていた$E[3]$の九候補のうち，
この指定collectionに対応するのは$\varepsilon=\OO_E$だけである．
この同値のt-exactnessは本定理には含めない．
\end{theorem}
\begin{proof}
$U_2=(L_0L_1)^{-1},U_1=L_0^{-1}$を\eqref{eq:actualdictionary}へ代入すると，
$q=L_0L_1^{-1}$，$L=L_1$を得る．
$\sigma$は$(E,L_1)$から$(E,L_0)$への埋込みつき同型なので，
$\sigma p_i$の$L_0$一般位置条件は$p_i$の$L_1$一般位置条件と等しい．
従って対応する可換blow-upはgenuineである．

初期collectionはVdBのsemiorthogonal decompositionと各$R\delta_*$の同値からfullである．
標的のSerre関手は$(-1)[2]$であり，境界triangleから，任意のforward exceptional pairの
RHomは境界でのRHomに自然にquasi-isomorphicとなる\cite[Proposition 9.3]{AOU}．
この同定はrestriction関手に由来するため合成を保存する．
よって命題\ref{prop:actualdictionary}を適用して，Rainsの三角dg模型と同定する．
これは全RHomと合成の同定であり，数値類だけの同定ではない．
$\sigma=t_b$なら$q=P_{3b}$で，$\rho(H)=L_1^3(-D_E)$である．
命題\ref{prop:boundarynormal}が表示した境界作用を与える．
最後に\eqref{eq:boundaryfirst-AOU}の$U_1$が実際の$L_0^{-1}$と等しいことから
$\varepsilon=\OO_E$となる．
\end{proof}

\begin{remark}[印刷tupleが現れる別のshift]
$S_-(M)=L_0\otimes\sigma^*M$と定義すれば，
$S_-^{-1}(\OO_E)=L_0^{-2}L_1$，
$S_-^{-2}(\OO_E)=L_0^{-5}L_1^3$となり，\eqref{eq:printedtuple}を再現する．
この計算は，二種類のshiftを混在させたときに不整合が起こる理由を説明する．
しかし$S_-$は\eqref{eq:right-plane-shift}の$S_+$とは異なり，
点対象と正の境界shiftの規約も同時に変更しなければならない．
本稿は$S_+$と明示したtransportを採用する．原典の全記号をどの変更で同時に読み替えるか，
また著者による正式な訂正の有無については別の確認事項とする．
\end{remark}

\subsection{正規化したdecupleと明示的な許容開集合}
$D_{123}=p_1+p_2+p_3$，$D_{456}=p_4+p_5+p_6$と置く．
mutationにはcohomological shiftの規格化も必要である．本節では
\begin{equation}\label{eq:mutation-shift-normalization}
 L_EF=\operatorname{Cone}(\RHom(E,F)\otimes E\to F),\qquad
 R_F^-E=\operatorname{Cone}(E\to\RHom(E,F)^*\otimes F)[-1]
\end{equation}
とし，これらが互いに逆のcollection操作になる規約を使う．
AOUの(9.29)--(9.43)と同じ隣接交換の順序でこの操作を行うが，
$R_6^-R_7^-R_8^-$の段階の後には，末尾の三つの変異対象を$[1]$でshiftする．
これはその境界像を次数2の\emph{線束}にするための指定である．
個々の対象のshiftはexceptionalityとfullnessを保ち，この三対象は以後の操作で動かない．
この手順を\emph{sheaf-normalized mutation列}と呼ぶ．
原典のmutation記号やline bundleの略記だけから，このshiftを暗黙に推論しない．
\begin{proposition}[actual restrictionに対する三条件]\label{prop:normalized-admissible}
相異なる六点について，命題\ref{prop:transported-tuple}の初期collectionへ
上記sheaf-normalized mutation列を施す．次の三条件を仮定する：
\begin{equation}\label{eq:normalized-open}
 \OO_E(D_{123})\not\simeq L_0,\qquad
 \OO_E(D_{456})\not\simeq L_1,\qquad
 \OO_E(D_E)\not\simeq L_0L_1.
 \tag{$U_+$}
\end{equation}
このとき最後のcollectionのrestrictionは
\begin{equation}\label{eq:normalized-decuple}
\begin{gathered}
 \OO_E,\quad (L_0L_1)^{-1}(D_E),\quad L_0^{-1}(D_{123}),\\
 \OO_E(p_1),\ \OO_E(p_2),\ \OO_E(p_3),\\
 L_0(-p_4),\ L_0(-p_5),\ L_0(-p_6)
\end{gathered}
\end{equation}
であり，これはadmissible elliptic decupleである．
この式の九線束のadmissibilityは，六点が相異なる範囲で正確に\eqref{eq:normalized-open}と同値である．
固定した$(E,\sigma,L_0)$に対して，この三条件とVdBの一般位置条件の共通部分は
$E^6$の非空Zariski開集合である．
\end{proposition}
\begin{proof}
境界でのmutationはexceptional pairのRHom比較を通じて評価triangleと両立する．
点と線束についての局所計算は$0\to M\to M(p)\to k_p\to0$による．
さらに次数$-2$の線束$M$を$\OO_E$の右にmutateする箇所では，$M^{-1}$が
次数2の大域生成線束なので$0\to M\to\OO_E^{\oplus2}\to M^{-1}\to0$を使う．
従って$R^-_{\OO_E}M=M^{-1}[-1]$であり，指定した$[1]$により$M^{-1}$になる．
点対象を線束の右へ動かす他の箇所では，$R^-_M k_p=M(p)$で追加shiftは不要である．
この二つの列を各mutationへ順に適用すると，最後の二回の直前のrestrictionは
\[
 (U_2(D_E),U_1(D_{123}),\OO_E,\OO_E(p_1),\OO_E(p_2),\OO_E(p_3),
 U_1^{-1}(-p_4),U_1^{-1}(-p_5),U_1^{-1}(-p_6))
\]
となる．最初の三対象の次数は零である．\eqref{eq:normalized-open}の第一・第三条件は
$U_1(D_{123})$と$U_2(D_E)$が非自明であることを意味するため，
最後の二回のmutationは$\OO_E$を変えず，\eqref{eq:normalized-decuple}を得る．

同じdegree blockで相異なることがadmissibilityである\cite[Definition 7.1]{AOU}．
次数1，2のblockでは点の相異性だけでよい．次数0の三線束の比較は，
第一・第三条件と，その比$L_1^{-1}(D_{456})$の非自明性，すなわち第二条件である．
各違反条件は和写像$E^k\to\Pic^k E$の一つのfiberであり真の閉集合である．
VdBの六点一般位置の違反も対角線，三点和，六点和の有限個の真の閉集合からなる．
$E^6$は既約なのでその有限和の補集合は非空開となる．
この議論では$\sigma$を固定しており，無限位数性を族におけるZariski開条件とはしていない．
\end{proof}

\begin{corollary}[規格化したAOU型同値への適用]\label{cor:normalized-pointed}
前命題の開集合上で，\eqref{eq:normalized-decuple}からAS代数$\mathsf A_+$を作り，
運ばれたcollectionのsheaf-normalized mutationからAOUの方法で同値を構成する．すると
\[
 \Psi_+:\Db(\qgr F)\xrightarrow{\sim}\Db(\qgr\mathsf A_+),\qquad
 \Psi_+(\pi F)=\OO_{\mathrm{AS},+}(0,0)
\]
はpointedであり，$F=R(e_{\mathrm{AS},+},s_{\mathrm{AS},+})$となる．
標準周期表示に対する右noetherian性を確認すれば，定理\ref{thm:AOUexactlocus}の十一条件は
この\emph{実際のdecupleから作ったpotential}に対する$\Psi_+$のt-exactnessを判定する．
\end{corollary}
\begin{proof}
前命題の最後のcollectionはfullであり，境界上の線束の次数と同次数での相異性から
forwardの高次Extが消える．Proposition 9.3を使えば圏内でもstrongで，End代数は
\eqref{eq:normalized-decuple}のEnd代数と同じである．
AOU \S7，Lemma 7.4とその直後のAS構成および\S9のtilting比較を使い，同じEnd代数を
持つAS標準collectionを得る．両方のfull strong collectionの比較が$\Psi_+$を与える．
前命題の第一・第三条件は条件\eqref{eq:two-nontrivial}を実際のrestrictionについて保証する．
定理\ref{thm:AOU-pointed}，系\ref{cor:AOU-ring}と定理\ref{thm:AOUexactlocus}を適用する．
これは印刷式(9.53)を変更せず代入したAS代数との同一視ではない．
この構成での右noetherian性と十一条件の非消滅は，
後の定理\ref{thm:normalized-abelian}で独立の追加仮定なしに導く．
\end{proof}

\begin{proposition}[decupleからの逆辞書]\label{prop:inverse-decuple}
原点を選んだ滑らかな楕円曲線$E$上のordered admissible decupleを
\[
 (\OO_E,A_0,B_0;
 K_1,K_2,K_3;
 C_1,C_2,C_3)
\]
と書く．$A_0,B_0$は次数0，$K_i$は次数1，$C_j$は次数2である．
\eqref{eq:normalized-decuple}の形に書くための線束と点は一意に
\begin{equation}\label{eq:inverse-decuple}
\begin{split}
 \OO_E(p_i)&=K_i\quad(1\le i\le3),\\
 L_0&=K_1K_2K_3B_0^{-1},\\
 \OO_E(p_{3+j})&=L_0C_j^{-1}\quad(1\le j\le3),\\
 L_1&=\OO_E(D_E)L_0^{-1}A_0^{-1},\\
 q&=L_0L_1^{-1}
\end{split}
\end{equation}
で復元される．$\sigma^*L_0=L_1$を満たす平行移動$\sigma$は九つ存在し，
それらは$E[3]$のtorsorをなす．全六点の相異性と$L_1$一般位置を追加すれば，
各選択に対するVdB構成と命題\ref{prop:normalized-admissible}によって，元のdecupleへ戻る．
無限位数の平行移動を要求するなら，さらに$q$の非torsion性を課す．
\end{proposition}
\begin{proof}
次数1の線束は一意の点の線束である．$B_0=L_0^{-1}(D_{123})$から$L_0$，
$C_j=L_0(-p_{3+j})$から残る点，$A_0=(L_0L_1)^{-1}(D_E)$から$L_1$が順に一意に決まる．
$\sigma=t_b$について$\sigma^*L_0=L_0P_{-3b}$なので，$[3]b$を指定する方程式となる．
複素楕円曲線の$[3]$は全射で核が$E[3]$，従って九つの解を持つ．
元のadmissibilityは三つの条件\eqref{eq:normalized-open}に戻るため，
一般位置の仮定の下で前の構成と互いに逆になる．$3b$が非torsionであることと$b$が非torsionであることは同値である．
\end{proof}
\begin{remark}
ここでの九つの$\sigma$は，\emph{decupleだけを入力として平面translationを復元する}際の曖昧さである．
既に$\sigma$と初期平面twistを固定した定理\ref{thm:right-dictionary}で，
$\varepsilon=\OO_E$と決まった九候補とは入力が異なる．
原点とline bundleの同型類を固定したPicardパラメータ上では，この九つの選択は$[3]$の引戻しによる有限étale被覆である．
自己同型で商を取った全モジュライstackが九対一で同型になる，あるいは九つの平面が
同じ標準heartを与えるという主張ではない．
\end{remark}

\begin{corollary}[整数倍点の明示的な共通領域]\label{cor:normalized-cyclic-family}
flex $o\in E$，非torsion $q=P_a$，相異なる正整数$N_1,\ldots,N_6$を取り，
$p_i=N_i a$とする．$L_1=\OO_E(3o)$，$L_0=L_1q$とし，$3b=a$を満たす$b$を選んで$\sigma=t_b$と置く．
このデータは$\sigma^*L_0=L_1$とVdBの一般位置条件を満たし，\eqref{eq:normalized-open}も満たす．
従って本節の規格化したpointed比較を適用できる．
\end{corollary}
\begin{proof}
$\sigma^*L_0=L_0P_{-3b}=L_1$である．$L_1$埋込みの三点共線は三つの$N_i$の和，
六点conic条件は全$N_i$の和が$q$の位数で零になることを意味するため起こらない．
従って$\sigma p_i$は$L_0$一般位置である．
三条件の左辺と右辺の比は順に
\[
 q^{N_1+N_2+N_3-1},\qquad
 q^{N_4+N_5+N_6},\qquad
 q^{\sum_iN_i-1}
\]
であり，全指数は正なので非自明である．
\end{proof}
\begin{remark}
この系は，整数倍点の共鳴例を実際の規格化したdecupleと接続する．
この三指数の計算だけでnoetherian性や十一因子$\Delta$の非零性を結論するわけではない．
それらは後の第\ref{sec:ellipticcoefficients}節で，実際のpotentialとvery strong helixを用いて証明する．
\end{remark}

\section{AS heartと反標準heart：helixからcornerへ}\label{sec:AS}
\subsection{AS構成が指定するもの}
AS heartは，適切なnoetherian性を満たす三次元Artin--Schelter正則
$\mathbb Z$-代数の$\qgr$として定義する．三次曲面の場合には
$\{0,1,2\}\times\mathbb Z$を添字とする代数を使い，標準的な九対象のfoundationがある
\cite[\S3--4]{AOU}．より一般のdel Pezzo型の標識とAS条件は
\cite[Definitions 1.1, 3.1]{OU}に従う．
ここでのAS正則次元3は曲面のhelix代数の次元であり，
非可換射影三空間を与える四次元AS正則ambientと同じものではない．
また非周期的な$\mathbb Z$-代数を，根拠なく一周期の次数付き代数に置き換えない．

very strong helixは例外性，特定の周期範囲でのstrong性，Serre作用との周期関係を与える．
そのhelixがheartを再構成するためには，生成性と高次Ext消滅を伴う豊富性も用いる．
しかし，九対象の全orbitの豊富性から，そのうちの一つ$e$のorbitの豊富性は自動的には従わない．
本稿で反標準環を作るのは後者であり，ここにcornerを導入する理由がある．

\subsection{実際のrestrictionを用いた構造対象の同定}
AOU\cite[\S9]{AOU}の設定を固定する．楕円曲線$Y$は滑らかで，
$L_0=\OO_Y(H)$，$L_1=\sigma^*L_0$は次数3，$\sigma$は無限位数の平行移動とし，
六点は同節の比較が成り立つ一般位置にあるとする．
同節の$\coh\widetilde X$は\emph{最初から}$\qgr F$であり，
構造対象は$e_F=\pi F$である．Lemma 9.2はSerre関手を$(-1)[2]$と同定し，
Theorem 9.4はAS代数$\mathsf A$に対するdg導来同値
\begin{equation}\label{eq:AOU}
\Psi:\Db(\qgr F)\xrightarrow{\sim}\Db(\qgr\mathsf A)
\end{equation}
を与える．右側の構造対象を$e_{\mathrm{AS}}=\OO(0,0)$と書く．
原典の構成では，mutation後のfoundationの先頭$E_0$が$e_{\mathrm{AS}}$へ移る．
以下では$E_0$を実際の対象として同定する．これは$K_0$や境界restrictionだけによる同定ではない．
本稿で新たに構成する$\Psi_+$については，第\ref{eq:mutation-shift-normalization}式の
mutation規約と明示したshiftを使う．以下の先頭対象の議論は，後方三対象のshiftを変更しても影響を受けない．

ここで境界の同定を一つ固定し，初期collectionの最後の三対象の実際のrestrictionを
$(\mathcal U_2,\mathcal U_1,\OO_Y)$と書く．$\mathcal U_2,\mathcal U_1$はそれぞれ
次数$-6,-3$の線束とする．$D_6=p_1+\cdots+p_6$，$D_3=p_1+p_2+p_3$と置き，
\begin{equation}\label{eq:two-degree-zero}
\mathcal L_A=\mathcal U_2\otimes\OO_Y(D_6),\qquad
\mathcal L_B=\mathcal U_1\otimes\OO_Y(D_3)
\end{equation}
を定める．両者の次数は零である．以下のpointed比較には，\emph{実際のrestriction}に対する条件
\begin{equation}\label{eq:two-nontrivial}
\mathcal L_A\not\simeq\OO_Y,\qquad\mathcal L_B\not\simeq\OO_Y
\tag{R}
\end{equation}
を明示的に仮定する．実際のdegree-zero blockがadmissibleならこの条件を満たす．
AOU v1の(9.47)に印刷された$\mathcal U_2=L_0^{-5}L_1^3$，$\mathcal U_1=L_0^{-2}L_1$を
他の境界式と同じ座標で無条件に使えるわけではない．第\ref{subsec:boundaryaudit}節の
検査がその障害を示すため，ここでは印刷式の代入でなく，条件\eqref{eq:two-nontrivial}を入力とする．
この形なら，本節の論証は座標式の規約に依存しない．
第\ref{thm:right-dictionary}定理の右加群規格化では実際の$\mathcal U_i$を計算し，
命題\ref{prop:normalized-admissible}の三条件の下で(R)を確認する．
そのdecupleから構成した同値は$\Psi_+$と書き，元の印刷式でラベル付けした同値とは区別する．

\begin{theorem}[AOUの構造対象の同定]\label{thm:AOU-pointed}
AOU導来比較の設定と実際のrestriction条件\eqref{eq:two-nontrivial}の下で，最後のmutationで得る対象は
$E_0\simeq e_F$である．従って\eqref{eq:AOU}は
\[
\Psi(e_F)\simeq e_{\mathrm{AS}},\qquad
\Psi\circ(1)\simeq s_{\mathrm{AS}}\circ\Psi,
\quad s_{\mathrm{AS}}=\Sfun_{\mathrm{AS}}^{-1}[2]
\]
を満たすpointed導来同値として規格化できる．
この結論だけでは$\Psi$のt-exactnessを主張しない．
\end{theorem}
\begin{proof}
AOUの(9.28)--(9.43)のmutation列を$\tau_1,\ldots,\tau_7$とする．
$\tau_1$の末尾は$e_F$である．$\tau_2,\tau_3$への操作は末尾に触れず，
$\tau_3\to\tau_4$の$R_6R_7R_8$はこの対象を変えずに第六位置へ動かす．
右mutationの組$(C,D)\mapsto(D,R_D C)$では右側の$D$自体は変わらないためである．
$\tau_4\to\tau_5$はその位置に触れず，$R_3R_4R_5$で第六位置から第三位置へ動く．
従って$\tau_6$の先頭三対象は$(A,B,e_F)$である．
ここまでの$e_F$の位置は$9,9,9,6,6,3$であり，いずれも対象そのものは保存されている．

点対象のspherical twistが線束を対応する点でtwistすることを用い，
AOU (9.48)--(9.52)のmutation計算を実際の$\mathcal U_i$で行うと，境界の導来restrictionは
$Li^*A=\mathcal L_A$，$Li^*B=\mathcal L_B$，$Li^*e_F=\OO_Y$である．
$(A,e_F)$と$(B,e_F)$は例外対なので，AOU Proposition 9.3により
\[
\RHom(A,e_F)\simeq R\Gamma(Y,\mathcal L_A^{-1}),\qquad
\RHom(B,e_F)\simeq R\Gamma(Y,\mathcal L_B^{-1}).
\]
楕円曲線の非自明な次数零線束の全cohomologyは零である．
実際，非零切断があれば次数零の有効divisorが零となって線束が自明になる．
従って$H^0=0$，Riemann--Rochから$H^1=0$である．
\eqref{eq:two-nontrivial}より上の二つのRHomは零となる．

左mutationの定義triangle
\[
\RHom(B,e_F)\otimes B\longrightarrow e_F
\longrightarrow L_Be_F\longrightarrow
\]
から，実際の自然射$e_F\to L_Be_F$は同型となる．
さらに$\RHom(A,L_Be_F)=0$なので$L_A L_Be_F\simeq e_F$．
これが最後の$L_1L_2$で得る先頭$E_0$である．
AOU (9.55)--(9.57)の同値は$E_0$を$e_{\mathrm{AS}}$へ送るので構造対象を保つ．
反Serre作用への適合は，dg同値によるSerre関手の一意性とLemma 9.2から従う．
\end{proof}

\begin{corollary}[同じ反標準環の回復]\label{cor:AOU-ring}
上のpointed同定と反Serre作用の整合的な反復を選ぶと，次数付き環の同型
\[
F\simeq R(e_{\mathrm{AS}},s_{\mathrm{AS}})
=\bigoplus_{n\ge0}\Hom(e_{\mathrm{AS}},s_{\mathrm{AS}}^n e_{\mathrm{AS}})
\]
を得る．従ってAS側から作る反標準heartは，AOUのsection modelの$\qgr F$と同じである．
ただし，AS標準heartそのものがこれと一致するという結論ではない．
\end{corollary}
\begin{proof}
まず$F_n\to\Hom_{\qgr F}(e_F,e_F(n))$が全$n\ge0$で同型であることを確認する．
AOU (9.1)の中心正則元$t$による短完全列を使う．境界の次数$n$の線束は
$n>0$で次数$3n$を持つため，$H^1$が零となる．
$e_F$のexceptionalityから始めて帰納すると
$\Ext^1(e_F,e_F(n))=0$ $(n\ge0)$である．
この短完全列の大域切断と
$0\to F_{n-1}\xrightarrow{t}F_n\to(F/tF)_n\to0$を比較し，
$F_0=\End(e_F)=k$を使えば自然な飽和写像は各次数で同型となる．
これらは乗法に適合する．定理\ref{thm:AOU-pointed}の同定を施すと主張を得る．
反Serre関手との自然同型は一回分を選んで整合的に反復するので，積の整合性も保たれる．
\end{proof}

\begin{proposition}[構造対象の同定の相対版]\label{prop:AOU-relative-object}
基礎$T$上で上記mutation列と$e_F$が相対的に定義されているとする．
$R\!\operatorname{Hom}_T(A,e_F)$，$R\!\operatorname{Hom}_T(B,e_F)$がperfectで，
導来基底変更に適合し，各幾何学的fiberで\eqref{eq:two-nontrivial}と
上のrestriction比較が成り立つと仮定する．
このとき評価triangleの自然射から$E_0\simeq e_F$が族として得られ，
その同型は任意の導来基底変更に適合する．
\end{proposition}
\begin{proof}
二つの相対Hom複体は全幾何学的fiberで零である．perfect complexに対する
導来Nakayama補題から双方は零となる．局所的に有限自由複体を取り，
任意の点のfiberがacyclicならその点の近傍で可縮にできることからもこの消滅が分かる．
従って最後の二つの評価triangleの自然射が相対的にも同型である．
評価射とconeの形成は指定したdg模型で基底変更に適合するので，得た同型も適合する．
\end{proof}
\begin{remark}
相対命題はmutation列・perfectness・基底変更を仮定している．
相対圏の構成，AS admissibilityの全退化点への延長，heart同値やモジュライstackの降下を
この命題だけから結論しない．ただし，その仮定下では$E_0$の同定に新たな
fiberごとのスカラー選択を要求せず，評価triangleの自然射が同型を与える．
\end{remark}

\begin{corollary}[VdB--ASの導来比較の範囲]\label{cor:VdBAS}
VdBの最後のsection-model比較とAOUの比較\eqref{eq:AOU}の両方が適用でき，
中間の$\qgr F$を同じデータで同定しているなら，それらの合成は
VdBの当該blow-up heartとAS heartの導来同値を与える．
そのt-exactnessと構造対象の保存は，これだけでは従わない．
\end{corollary}
\begin{proof}
二つの指定dg同値を合成する．VdB側では既に非自明なHRS tiltがあり得るため，
合成がt-exactであることは同値性だけの帰結ではない．構造対象の保存には各関手での像を照合する必要がある．
AOU側の保存は定理\ref{thm:AOU-pointed}で確定したが，VdB側の関手との
同じデータにおける適合性は別に指定する．
\end{proof}

\subsection{一対象の反標準orbitを判定する}
AS側の標識されたfoundationを$E_1=e,E_2,\ldots,E_r$とし，
$T_0=\bigoplus_iE_i$，$B_n=\Hom(T_0,s^nT_0)$と置く．
$E_1$への射影を$p\in B_0$とすると，積の規約を合わせれば
\begin{equation}\label{eq:foldedcorner}
pBp=\bigoplus_{n\ge0}\Hom(e,s^ne)=R(e,s).
\end{equation}
以下の条件は，任意のAS退化点で自動的に成立するものではない．

\begin{corollary}[AS--反標準のabel同値判定]\label{cor:ASac}
上のfoldingが$\AAS\simeq\qgr B$を与え，$e$が$\pi(pB)$に対応するとする．
$B$と$R=pBp$が局所有限右noetherianで，$Bp$が右$R$有限生成，
$pB$が左$R$有限生成であると仮定する．このとき自然なcorner関手
\[
K:\AAS\simeq\qgr B\longrightarrow\qgr R(e,s)
\]
は，$\dim_k B/BpB<\infty$と同値にabel同値となる．
その場合$K(e)=\pi R$であり，確認したfoldingの下で$s$と$(1)$にも適合する．
\end{corollary}
\begin{proof}
式\eqref{eq:foldedcorner}と命題\ref{prop:corner}を使う．
$(pB)p=pBp$が構造対象の保存を与える．反標準次数のshiftと$s$の適合性は
foldingの定義に含めたものである．
\end{proof}

この系が同定するのは，\emph{AS側の指定対象から作った}反標準heartである．
実際のrestriction条件\eqref{eq:two-nontrivial}を確認したAOU領域では，
系\ref{cor:AOU-ring}がこれを$F$と同定する．
独立に指定したRains側の環との比較には，そのパラメータと構造対象の同定をなお要する．
$B/BpB$が無限次元でも，それはこの自然なabelian corner関手が同値でないことを意味するだけで，
別の導来同値の不存在ではない．

\begin{corollary}[指定されたAOU同値のt-exactnessへの帰結]\label{cor:AOU-texact}
実際のrestriction条件\eqref{eq:two-nontrivial}を満たすAOU領域で，系\ref{cor:ASac}のfoldingと有限性条件が成立し，
$\dim_k B/BpB<\infty$であるとする．このときAOUの指定同値$\Psi$自体がt-exactである．
\end{corollary}
\begin{proof}
corner同値$K$と系\ref{cor:AOU-ring}から，AS標準heartは
$(e_{\mathrm{AS}},s_{\mathrm{AS}})$に関してAZ生成性と全高次Extのtail消滅を満たす．
これらは$\qgr F$の標準heartから$K$で運ばれる性質である．
定理\ref{thm:AOU-pointed}により$\Psi$はpointedなので，
命題\ref{prop:unique}の認識式を両側に適用するとt-exactnessが従う．
\end{proof}
この系は，指定したAOU関手と抽象的なcorner同値とを接続する．
規格化した実際の共通領域では，追加の有限性条件を後の定理\ref{thm:normalized-abelian}で満たすことを示す．
全AOU領域で有限corner条件を証明したという意味ではない．

\begin{lemma}[有限次数による判定が有効になる条件]\label{lem:cutoff}
$C=B/BpB$が非負次数付きで，有限次元$C_0$上，次数1で生成されるとする．
このとき$C_N=0$なら全$n\ge N$で$C_n=0$である．
従って$C$が有限次元であることと，ある$N$で$C_N=0$であることは同値である．
\end{lemma}
\begin{proof}
任意の次数$n\ge N$の生成単項式は，長さ$N$の先頭部分を含む．
$C_N=0$ならそれが零となる．逆は次数ごとの局所有限性による．
\end{proof}
三次曲面のJacobi表示では道長次数と反標準cut次数は異なる．
補題を道長次数に適用した結果を反標準cornerへ移すには，この二つのgradingの関係を確認する．
単に有限次数の計算が良好だったという理由だけで，全次数の有限性を結論しない．

\begin{remark}[三次曲面の具体的corner計算との関係]
同プロジェクトの別稿\nolinkurl{cubic_heart_corner.tex}は，指定された九頂点表示について，
四交比の十一のdeterminant条件，一様な道長7での消滅，および
$8+21t+33t^2+38t^3+33t^4+21t^5+8t^6$という有限cornerのHilbert関数を扱う．
一様な道長境界と全Hilbert関数の詳細は別稿に保持する．本稿では次節で，
有限性に必要十分な十一条件の証明を自立した形で収め，指定AOU同値の
t-exactnessへ接続する．この係数判定から全potentialのAS正則性を結論しない．
\end{remark}


\section{三次曲面型の一致領域：四交比による完全判定}\label{sec:explicitcorner}
本節では一般のcorner条件を，AOUの三block表示で係数の条件に置き換える．
有限cornerの幾何学的判定は同プロジェクトの別稿で扱っているが，
ここでは必要な証明を全て収め，前節で確定したpointed性と接続する．
純粋なJacobi環の命題と，AS heartへ適用する際のnoetherian仮定を区別する．

\subsection{二つのgradingと自動的に得られる有限性}
頂点を$A_a,B_b,C_c$ $(0\le a,b,c\le2)$，矢を
$x_{ab}:A_a\to B_b$，$y_{bc}:B_b\to C_c$，$z_{ca}:C_c\to A_a$とする．
\[
 W(w)=\sum_{a,b,c}w_{abc}[z_{ca}y_{bc}x_{ab}],\quad
 B(w)=k\widetilde Q/(\partial W),\quad p=e_{A_0},\quad C(w)=B(w)/B(w)pB(w).
\]
全矢の次数を1とする道長次数を$\ell$，$x,y$の次数を0，$z$の次数を1とする
次数をcut次数と呼ぶ．block番号を$\lambda(A)=0,\lambda(B)=1,\lambda(C)=2$とすると，
始点$u$，終点$v$の道について
\begin{equation}\label{eq:cutpath}
 \ell=3\deg_{\mathrm{cut}}+\lambda(v)-\lambda(u).
\end{equation}
従って両次数での右有界性は同値であり，cutの各次数も有限次元である．
cut次数のcoverの$(A_a,n),(B_b,n),(C_c,n)$をそれぞれ
$(a,3n),(b,3n+1),(c,3n+2)$に対応させると，AOUの無限quiverそのものになる．
このreindexingは加群圏とtorsion部分圏を同定する．標準の周期的なAS実現では
\begin{equation}\label{eq:explicitfold}
 \AAS\simeq\qgr B^{\mathrm{cut}},\quad
 e_{\mathrm{AS}}=\pi(pB),\quad s_{\mathrm{AS}}=(1)_{\mathrm{cut}},\quad
 pB^{\mathrm{cut}}p=R(e_{\mathrm{AS}},s_{\mathrm{AS}}).
\end{equation}
最後のHomの飽和と周期性にはAS projective helixの既知の性質
\cite[\S4]{AOU}\cite[(3.9)--(3.10), Remark 5.4]{OU}を用いる．
道長gradedの$\qgr B$を，そのまま一つのAS heartと同一視してはいけない．

\begin{lemma}[有限cornerから両bimoduleの有限性へ]\label{lem:automaticcorner}
有限quiverの道長斉次関係による商$B$と頂点の和$p$について，$B/BpB$が有限次元なら，
$Bp$は右$pBp$有限生成，$pB$は左$pBp$有限生成である．さらに$B$が右noetherianなら$pBp$も右noetherianである．
\end{lemma}
\begin{proof}
$C_N=0$なる$N$を取る．任意の長さ$N$の道は，$p$の頂点を通る同じ長さの道の線形結合となる．
$p$から出る長い道の終端側$N$個をこの形で書き換え，最後に$p$を通る位置で切れば，
長さ$\le N$の$Bp$の道と$pBp$の道の積になる．従って前者の有限個の道が右加群を生成する．
逆に$p$へ入る道の始端側を用いれば左有限生成性も従う．この操作はcut次数にも斉次的である．
$R=pBp$の右イデアル$I$に対して$(IB)p=I$なので，右イデアルの増大列は
$B$のnoetherian性により停止する．
\end{proof}
従って標準的な右noetherian AS表示では，有限cornerを証明した後に
両bimoduleの有限性を別の未解決仮定として残す必要はない．

\subsection{係数による有限cornerの必要十分条件}
$C(w)$は$w_{1bc},w_{2bc}$の十八係数だけに依存する．これらが非零なら
$q_{bc}=w_{2bc}/w_{1bc}$，$Q=(q_{bc})$と置く．$Q^{\langle-1\rangle}$は成分ごとの逆数行列であり，
通常の逆行列ではない．

\begin{theorem}[十一条件による判定]\label{thm:eleven}
任意の$w\in\mathbb A^{27}$について，$\dim_k C(w)<\infty$となる必要十分条件は，
十八係数$w_{1bc},w_{2bc}$が全て非零で，次の十一個の値が全て非零となることである：
\begin{equation}\label{eq:elevenconditions}
 q_{bc}q_{b'c'}-q_{bc'}q_{b'c}\quad(b<b',\ c<c'),\qquad
 \det Q,\qquad \det Q^{\langle-1\rangle}.
\end{equation}
この命題は$B(w)$のAS正則性やnoetherian性を仮定しない．
\end{theorem}
\begin{proof}
\emph{必要性．}
一つの$w_{abc}$ $(a=1,2)$が零なら，三角形$x_{ab},y_{bc},z_{ca}$だけを残す．
各矢を$t$倍の行列単位へ送る$k[t]$上の表現が全Jacobi関係を満たし，
三角閉路の像は$t^3$倍の冪等元となる．その全ての冪が非零なので$C(w)$は無限次元である．

十八係数が非零なら$y$のrescalingで$w_{1bc}=1$とできる．
二行二列の小行列式が零なら，その長方形上で$q_{bc}=r_bs_c$と書ける．
行和・列和が零の$T=\left(\begin{smallmatrix}1&-1\\-1&1\end{smallmatrix}\right)$を置き，
その支持で$x_{1b}=z_{c1}=1$，$x_{2b}=-r_b^{-1}$，$z_{c2}=s_c^{-1}$，$y_{bc}=T_{bc}$，
他の矢を零とする．全関係は行和・列和または$1-q_{bc}/(r_bs_c)$に帰着する．
再び各矢に$t$を掛けると非零の三角閉路の全冪が残る．

$\det Q=0$なら$Q=UV$，$U\in M_{3,2}(k)$，$V\in M_{2,3}(k)$と書き，
$0\ne h\in\ker U^{\mathsf t}$，$0\ne k\in\ker V$を取る．
$A_1$をrank 2，他の残る頂点をrank 1とする$k[t]$表現を
\[
 x_{1b}=tU_{b\bullet},\quad z_{c1}=-tV_{\bullet c},\quad
 x_{2b}=z_{c2}=t,\quad y_{bc}=th_bk_c
\]
で定める．$\partial y$の関係は$-UV+Q=0$，他の関係は
$Vk=0,h^{\mathsf t}U=0,Qk=0,h^{\mathsf t}Q=0$である．
$h_bk_c\ne0$の閉路が$t^3h_bk_c$に写るため無限次元となる．
$\det Q^{\langle-1\rangle}=0$の場合は$A_1,A_2$を交換して同じ議論を使う．

\emph{十分性．}
$P=(Q^{\langle-1\rangle})^{\mathsf t}$の列を$p_1,p_2,p_3$とし，
$p_4,p_5,p_6$を座標点とする．各成分と九つの小行列式と$\det P$の非零性は，
六点のどの三点も共線でないことに等しい．座標三点を通るconicは
$aX_0X_1+bX_1X_2+cX_2X_0=0$であり，残る三点の評価行列は，各行の非零因子を除き
$P$の成分逆数の転置である．従って六点がconic上にない条件は$\det Q\ne0$．
ゆえに$S=\operatorname{Bl}_{p_1,\ldots,p_6}\mathbb P^2$は滑らかな可換三次曲面である．

この六点からAOUの九線束helixと幾何学的potential $w^{\rm geo}$を作ると，矢の基底を選んで
\begin{equation}\label{eq:twoseries}
 w^{\rm geo}_{1bc}=P_{cb},\qquad w^{\rm geo}_{2bc}=1
\end{equation}
となる\cite[(8.4)--(8.7)]{AOU}．この恒等式は次のようにも確認できる．
添字$b$を3で読み，$u_b=(P_{0b}P_{1b},P_{1b}P_{2b},P_{2b}P_{0b})$，
$D_b=(u_{b+1}\times u_{b+2})\cdot(X_0X_1,X_1X_2,X_2X_0)$，
$M_b=(P_{\bullet,b+1}\times P_{\bullet,b+2})\cdot X$，$L_{bc}=\det[p_{b+1},p_{c+4},X]$とする．
$D_b$は$p_{b+1}$以外の五点を通るconicであり，Cramerの恒等式から
\[
 \sum_bP_{cb}D_bL_{bc}=0,\qquad \sum_bM_bL_{bc}=0.
\]
例えば$c=0$の前者は
$X_1\sum_bP_{0b}P_{2b}D_b-X_2\sum_bP_{0b}P_{1b}D_b=0$
で，二つの和は共通の行列式を掛けた$X_2X_0$と$X_0X_1$に等しい．
他の$c$は巡回対称性による．各関係空間は一次元なので\eqref{eq:twoseries}が従う．
従って二系列の比は元の$Q$と等しく，$C(w)\simeq C(w^{\rm geo})$となる．
ここでは$w_{0bc}$を幾何学的な値へ変えてよい．それらは商$C$に影響しないからである．

$H=-K_S$とすると幾何学的rolled-up環は
$\bigoplus_{v,v',n\ge0}\Hom(E_v,E_{v'}(nH))$である\cite[(4.9)]{AOU}．
各$v$に対し$E_v^\vee(mH)$が大域生成される$m$を取り，評価の核を$K_v$とする．
評価列を$E_{v'}((n-m)H)$でtensorしてSerre消滅を使うと，$n\gg0$で
\[
 H^0(E_v^\vee(mH))\otimes H^0(E_{v'}((n-m)H))
 \twoheadrightarrow\Hom(E_v,E_{v'}(nH))
\]
となる．左辺は全て$\OO_S(mH)$を経由するため，構造頂点を通る積である．
有限個の$v,v'$に共通の閾値を取れば全高cut次数で$B_n\subset BpB$となり，
式\eqref{eq:cutpath}から商は有限次元である．
\end{proof}

\subsection{四交比と指定AOU同値のt-exactness}
行・列rescalingを除き
\[
 Q\sim\begin{pmatrix}1&1&1\\1&a&b\\1&c&d\end{pmatrix},
\qquad a,b,c,d\in k^\times
\]
とできる．四交比は$q_{ij}q_{00}/(q_{i0}q_{0j})$ $(i,j=1,2)$である．
\begin{equation}\label{eq:Deltaheart}
\begin{split}
\Delta={}&(a-1)(b-1)(c-1)(d-1)(a-b)(c-d)(a-c)(b-d)(ad-bc)\\
&\cdot\bigl((a-1)(d-1)-(b-1)(c-1)\bigr)\\
&\cdot\bigl(bc(1-a)(1-d)-ad(1-b)(1-c)\bigr).
\end{split}
\end{equation}
最後の二因子は$\det Q$と$abcd\det Q^{\langle-1\rangle}$である．
従って有限corner条件は正確に$\Delta\ne0$である．

\begin{lemma}[標準AS helixからの高次消滅]\label{lem:ASfulltail}
標準noetherian AS実現で，foundationを$E_v$とする．任意の$M\in\AAS$に対して
$\Ext^j(E_v,s^nM)=0$ $(j>0,n\gg0)$である．
\end{lemma}
\begin{proof}
$M$をnoetherianなAS加群の像として取り，AS正則性とnoetherian性から有限生成projectiveの
有限解消を選ぶ．$\pi$を施すと，各項が有限個のhelix対象からなる解消を得る．
標準projective helixはvery strongであり，十分大きい$s$-shiftの後では各項に対する
$\Ext^{>0}(E_v,-)$が零となる\cite[Definition 3.1と(3.10)の直前]{OU}．
有限個の項に共通の$n$を取り，この解消へ$\RHom(E_v,s^n-)$を適用する．
項は非正cohomological次数に置かれるので，正次数のcohomologyはない．
\end{proof}

\begin{theorem}[指定AOU比較の明示的な必要十分条件]\label{thm:AOUexactlocus}
AOUの導来比較の領域で，実際のrestriction条件\eqref{eq:two-nontrivial}を満たす
定理\ref{thm:AOU-pointed}の同値$\Psi$を取り，
その標準周期AS表示$B(w)$が右noetherianであるとする．以下は同値である．
\begin{enumerate}
\item $\Psi$は標準heartについてt-exactである．
\item 自然なcorner $(-)p:\AAS\to\qgr F$はabel同値である．
\item $C(w)$は有限次元である．
\item $w_{1bc},w_{2bc}$は全て非零で，四交比が$\Delta\ne0$を満たす．
\end{enumerate}
従ってこの表示では，両corner bimoduleの有限性は独立の追加仮定ではない．
\end{theorem}
\begin{proof}
(iii)$\Leftrightarrow$(iv)は定理\ref{thm:eleven}．
(iii)$\Rightarrow$(ii)は補題\ref{lem:automaticcorner}と命題\ref{prop:corner}，
(ii)$\Rightarrow$(iii)は$\pi C$が零へ送られることによる．
pointed性により$R(e_{\mathrm{AS}},s_{\mathrm{AS}})=F$である．
(ii)ならAS側の$(e,s)$は標準反標準pairのAZ豊富性を継承する．
補題\ref{lem:ASfulltail}および標的のSerre消滅を用い，命題\ref{prop:unique}を
pointedな$\Psi$へ適用すると(i)を得る．

逆に(i)なら$(e,s)$はAS heartでAZ-ampleとなる．各$E_{v'}$に対し
$\bigoplus_l s^{-m_l}e\twoheadrightarrow E_{v'}$を取り，核を$K$とする．
補題\ref{lem:ASfulltail}から，十分大きい$n$で
\[
 \bigoplus_l\Hom(E_v,s^{n-m_l}e)\twoheadrightarrow\Hom(E_v,s^nE_{v'})
\]
となる．各像は$E_v\to s^{n-m_l}e\to s^nE_{v'}$と分解され，$n\ge m_l$なら
両因子は非負cut次数の$B$に属する．有限個の$v,v'$に共通の閾値を取ると，
全高cut次数が$BpB$に入り(iii)を得る．
\end{proof}

\begin{corollary}[非空開集合と古典的六点モジュライ]\label{cor:coefficientopen}
十八係数の矢rescaling商は四交比トーラス$(\mathbb G_m)^4$であり，
その有限corner開集合$D(\Delta)$は，順序付き一般位置六点のmoduli $M_{\mathrm{cubic}}$と同型である．
全二十七係数非零のAOU係数商トーラス$(\mathbb G_m)^8$上で，有限corner locusは
この$D(\Delta)$の逆像である．その写像は全射のtorus準同型で切断を持ち，
従ってこの係数開集合は非標準的に$M_{\mathrm{cubic}}\times(\mathbb G_m)^4$と同型である．
\end{corollary}
\begin{proof}
$y$のrescalingで第一系列を1とし，残りの$x,z$のrescalingで$Q$を行・列正規化する．
$P=(Q^{\langle-1\rangle})^{\mathsf t}$による六点復元は，座標三点を固定した後の
行・列rescalingと一致する．定理\ref{thm:eleven}の証明が一般位置条件との同値を与える．
全係数商からの写像はこれら四つの独立な指標で与えられる．
$w_{0bc}=w_{1bc}=1$，$w_{2bc}=Q$と置けばtorus準同型としての切断を得る．
核の次元は$8-4=4$である．
\end{proof}
例えば$(a,b,c,d)=(-1/4,1/16,1/6,-1/4)$で全十一因子が非零であり，開集合は非空である．
AOU Theorem 1.1のbirational対応の定義域では，この係数開集合の逆像も非空稠密開である．
\emph{圏の結論は常に定理\ref{thm:AOUexactlocus}の適用域との共通部分で使う．}
とくに，無限位数translationの条件自体をZariski開と主張せず，全potentialのAS正則性や
全AS対象のnoetherian性もこの係数計算から導かない．

\begin{remark}[検算と適用範囲]
\texttt{verify\_four\_hearts\_corner.py}は，六点復元の六つの恒等式を12変数の整数係数多項式として
全係数で検算する．上の具体係数では，全Jacobi関係の全道による積を列挙して
有理数と$\mathbb F_{1000003}$でそれぞれ
$(\dim C_0,\ldots,\dim C_7)=(8,21,33,38,33,21,8,0)$を得る．
次数1生成ゆえ$C_7=0$は全高次数消滅を含意する．
しかし，定理\ref{thm:eleven}の全パラメータでの十分性は有限個の検算でなく六点幾何の証明による．
また，$\Delta=0$なら指定$\Psi$のt-exactnessは失敗するが，AOUの導来同値は適用条件内で残る．
その点で負指数のRains根が必ず存在することや，異なる構造対象を許す全abel同値の不存在は主張しない．
\end{remark}

\section{楕円decupleのpotentialと十一条件の幾何学的意味}\label{sec:ellipticcoefficients}
本節では，第\ref{sec:boundarydata}節で規約を固定したdecupleの関係式を直接計算する．
その結果，構造頂点を除くcornerは$L_0$に依存せず，
$L_1$による\emph{元の六点}の可換三次曲面cornerに一致することが分かる．
従って第\ref{sec:explicitcorner}節で係数空間に書いた十一条件を，
新しく作った補助的な六点でなく，入力の楕円データへ引き戻せる．
最後に，pointedなvery strong helixからの再構成を用いて，
この共通領域で残っていた右noetherian性を導出する．

\subsection{次数二の切断の関係式}
$E=\mathbb C/(\mathbb Z+\tau\mathbb Z)$とし，$\vartheta(z)$を奇Jacobi theta関数とする．
$\vartheta(z+1)=-\vartheta(z)$，
$\vartheta(z+\tau)=-e^{-\pi i\tau-2\pi iz}\vartheta(z)$と規格化する．
次数$d$，Abel和$\lambda$の線束の切断を，乗数
$(-1)^d$と$(-1)^de^{-\pi id\tau-2\pi idz+2\pi i\lambda}$を持つ関数で表す．
以下の数は$E$の点の複素数liftを選んだものであり，等式はこの選択と整合するliftで読む．
選択を変えると矢・関係式の基底が変わるが，零条件と後述の行列の行・列rescaling類は変わらない．

\begin{lemma}[三つのdegree-two sectionの関係]\label{lem:theta-relation}
$r_1,r_2,r_3$を$E$上で相異なる点とし，$d$を固定する．
$f_b(z)=\vartheta(z-r_b)\vartheta(z-d+r_b)$と置く．
添字を$3$で巡回させ，
\[
 c_b=\vartheta(r_{b+1}-r_{b+2})\,
       \vartheta(r_{b+1}+r_{b+2}-d)
\]
とすると，$\sum_{b=1}^3 c_bf_b=0$である．これが三切断間の唯一の関係を生成する．
特に$c_b=0$と$r_{b+1}+r_{b+2}=d$ in $E$は同値である．
\end{lemma}
\begin{proof}
三切断は同じ次数2の線束に属する．$z=r_1$へ代入すると第一項は零で，
残りの二項は$\vartheta$の奇性により相殺する．$z=r_2,r_3$でも同様である．
従ってその和は三つの異なる零点を持つ次数2の切断なので恒等的に零．
二切断が比例するのは有効零divisor
$\{r_b,d-r_b\}$が一致するときだけである．三点が相異なるので，
三切断すべてが比例することはなく，spanは次元2，関係空間は次元1である．
最後の主張もthetaの零点と相異性から従う．
\end{proof}

$D_x=p_1+p_2+p_3$，$D_y=p_4+p_5+p_6$とし，
$x_b=p_b$，$y_c=p_{3+c}$ $(1\le b,c\le3)$と書く．
$L_i$のAbel和を$\ell_i$，$X=\sum_bx_b$，$Y=\sum_cy_c$とする．
第\ref{eq:normalized-decuple}式の三blockを
\[
 A_0=\OO_E,\quad A_1=(L_0L_1)^{-1}(D_x+D_y),\quad
 A_2=L_0^{-1}(D_x),\quad
 B_b=\OO_E(x_b),\quad C_c=L_0(-y_c)
\]
とし，$u=X-\ell_0$，$t=\ell_1-Y$と置く．
条件$(U_+)$は$u,t,u-t$が$E$上で非零という条件である．
矢を$\mathsf x_{ab}:A_a\to B_b$，$\mathsf y_{bc}:B_b\to C_c$，
関係式に双対なrolled-up矢を$\mathsf z_{ca}$とする．

\begin{theorem}[正規化decupleの全potential]\label{thm:theta-potential}
相異なる六点と条件$(U_+)$の下で，矢の基底と関係式の基底を選べば，
正規化decupleから得るpotentialは
\begin{equation}\label{eq:theta-potential}
 W_+=\sum_{a=0}^2\sum_{b,c=1}^3
 \vartheta(y_c-x_b+t_a)
 [\mathsf z_{ca}\mathsf y_{bc}\mathsf x_{ab}],
 \qquad(t_0,t_1,t_2)=(u,t,0)
\end{equation}
と書ける．これは解析的な局所座標表示であり，矢rescalingを除く関係式の同型類を与える．
特に$w_{1bc},w_{2bc}\ne0$の範囲では，係数比行列は行・列rescalingを除き
\begin{equation}\label{eq:theta-Q}
 Q_{bc}=\frac{\vartheta(y_c-x_b)}{\vartheta(y_c-x_b+t)}
\end{equation}
である．そのrescaling類は$L_0$に依存しない．
\end{theorem}
\begin{proof}
$A_a$のAbel和を$a_a$，$C_c$のAbel和を$c_c=\ell_0-y_c$とする．
各一次元Homの基底を
$\vartheta(z-x_b+a_a)$と$\vartheta(z-c_c+x_b)$で表す．
その積に補題\ref{lem:theta-relation}を適用すると，関係係数は
\[
 \vartheta(x_{b+1}-x_{b+2})\,
 \vartheta(X-x_b-a_a-c_c)
 =\vartheta(x_{b+1}-x_{b+2})\,
   \vartheta(y_c-x_b+u-a_a).
\]
第一因子は非零で$b$だけに依存するから，矢$\mathsf x_{ab}$の基底変更で除ける．
各$(a,c)$の関係を規格化する自由度は$\mathsf z_{ca}$の基底変更である．
$a_0=0$，$a_1=u-t$，$a_2=u$を代入して\eqref{eq:theta-potential}を得る．
AOUのpotentialは各長さ二関係とその双対矢の積の和である
\cite[(4.7), \S7]{AOU}ため，これが全27係数を与える．
一般の矢基底の変更では$Q$は行・列rescalingだけ変わり，式\eqref{eq:theta-Q}の同型類は不変である．
\end{proof}

\subsection{係数比は入力六点の射影座標である}
$\OO_E(D_y)\not\simeq L_1$を仮定する．
$L_1$で埋め込んだ$p_4,p_5,p_6$は非共線なので，それらを座標三点へ送れる．
その座標で最初の三点を列に持つ行列を$P=(P_{cb})$とする．

\begin{theorem}[元の六点によるcornerの復元]\label{thm:original-six-points}
上の設定で，$w_{1bc}$がすべて非零なら
\begin{equation}\label{eq:original-P}
 P\sim (Q^{\langle-1\rangle})^{\mathsf t}
\end{equation}
が行・列rescalingを除いて成り立つ．従って
\[
 B(W_+)/B(W_+)e_{A_0}B(W_+)
\]
のgraded同型類は$L_0$に依存しない．六点が一般位置にある場合，
$L_1$埋込みでの同じ六点の滑らかな可換三次曲面から得るcornerと一致し，この商は有限次元である．
\end{theorem}
\begin{proof}
$\{1,2,3\}\setminus\{c\}=\{d,e\}$とする．$L_1$の切断
\[
 h_c(z)=\vartheta(z-y_d)\vartheta(z-y_e)\vartheta(z-y_c-t)
\]
の零divisorは$y_d+y_e+(y_c+t)$で，Abel和は$Y+t=\ell_1$である．
$h_c(y_c)\ne0$，$h_c(y_d)=h_c(y_e)=0$なので，三切断は座標三点を規格化する基底となる．
$G(z)=\prod_{j=1}^3\vartheta(z-y_j)$と置くと，$x_b\ne y_j$より
\[
 \frac{h_c(x_b)}{G(x_b)}
 =\frac{\vartheta(x_b-y_c-t)}{\vartheta(x_b-y_c)}
 =\frac{\vartheta(y_c-x_b+t)}{\vartheta(y_c-x_b)}=Q_{bc}^{-1}.
\]
切断基底と各点のfiber trivializationを変える自由度が行・列rescalingであるから
\eqref{eq:original-P}を得る．
構造頂点を除いた商は二系列$w_1,w_2$だけに依存し，それらが全て非零のとき
その矢rescaling類は$Q$の行・列rescaling類で決まる．
一般位置の場合には，第\ref{sec:explicitcorner}節の可換幾何の表示
$w_1^{\rm geo}=P^{\mathsf t},w_2^{\rm geo}=1$と比較すれば
主張した商の同型を得る．有限性は定理\ref{thm:eleven}による．
\end{proof}

\begin{corollary}[十一条件の楕円パラメータへの引戻し]\label{cor:elliptic-discriminant}
六点は相異なり，$\OO_E(D_y)\not\simeq L_1$とする．
このとき$w_{2bc}$は常に非零であり，
\[
 w_{1bc}=0\quad\Longleftrightarrow\quad
 L_1\simeq\OO_E(x_b+D_y-y_c).
\]
全$w_{1bc}$が非零となる範囲では，九つのminorと二つのdeterminantの零条件は
\begin{align*}
 Q_{bc}Q_{b'c'}-Q_{bc'}Q_{b'c}=0
 &\quad\Longleftrightarrow\quad
 L_1\simeq\OO_E(x_b+x_{b'}+y_e),
 \quad\{e\}=\{1,2,3\}\setminus\{c,c'\},\\
 \det Q^{\langle-1\rangle}=0
 &\quad\Longleftrightarrow\quad L_1\simeq\OO_E(D_x),\\
 \det Q=0
 &\quad\Longleftrightarrow\quad L_1^2\simeq\OO_E(D_x+D_y).
\end{align*}
従って十八係数非零と十一条件の同時成立は，
$L_1$による元の六点が一般位置にあることと同値である．
\end{corollary}
\begin{proof}
係数の条件は\eqref{eq:theta-potential}から従う．
$2\times2$ minorは非零成分の逆数を取っても零条件が同じであり，
\eqref{eq:original-P}を通じて，最初の二点と残り一つの座標点の共線性に等しい．
$\det Q^{\langle-1\rangle}$は最初の三点の共線性を表す．
六点がconic上にある条件は，座標三点を通るconicの残り三点への評価行列のdeterminantが零という条件である．
その評価行列は，各行の非零因子を除き$P$の成分逆数であり，零条件は$\det Q=0$となる．
楕円曲線上の三点共線条件と六点conic条件はそれぞれ
$L_1\simeq\OO_E(\text{三点和})$，$L_1^2\simeq\OO_E(\text{六点和})$である．
後者の十分性は，平面三次$E$に対する
$H^0(\mathbb P^2,\OO(2))\simeq H^0(E,L_1^2)$からも従う．
最後の三点の共線性は初めから排除されており，これで全20個の三点条件とconic条件を尽くす．
\end{proof}
\begin{remark}
式\eqref{eq:theta-Q}の逆数行列はelliptic Cauchy行列である．
Frobeniusのdeterminant公式\cite[(1.7), (2.1)]{EllipticCauchy}を使っても上の因子を得る．
本稿ではその恒等式自体を新結果とせず，元の六点の射影座標を直接同定する証明を与えた．
全27係数が非零という追加条件は不要である．実際，
$w_{0bc}=0$は$L_0\simeq\OO_E(D_x-x_b+y_c)$を意味し，
$L_1$一般位置でも起こり得るが構造頂点を除いたcornerには影響しない．
\end{remark}

\subsection{ampleな構造対象からのnoetherian再構成}
次の命題は，有限cornerだけから一般のJacobi環のnoetherian性を導く主張ではない．
\emph{既に存在する}pointed導来比較の中に，full very strong周期helixがあることを使う．

\begin{proposition}[ample構造対象を含むhelix]\label{prop:ample-helix-noetherian}
$\cH$はHom有限noetherian abel圏で，homological dimensionは有限$d$とする．
$(e,s)$はAZ-ampleで，全coherent対象への高次Ext tailが消えるとする．
$\Db(\cH)$にfull very strong helixがあり，その周期作用が$s$，
foundationが$(E_1=e,E_2,\ldots,E_l)$であるとする．
このとき全$E_i$は$\cH$に属し，$T=\bigoplus_iE_i$について$(T,s)$もAZ-ampleとなる．
従って$B_T=\bigoplus_{n\ge0}\Hom(T,s^nT)$は右noetherianで，
$\qgr B_T\simeq\cH$となる．
この同値は各$E_i$と対応する標準projectiveの像，およびそれらの射を保存する．
\end{proposition}
\begin{proof}
very strong性から$n\gg0$で$\Hom(e,s^nE_i[k])=0$ $(k\ne0)$となる．
命題\ref{prop:unique}のheart認識を標的$\cH$に適用し，$E_i\in\cH$を得る．
同様に$\Ext^{>0}(T,s^ne)=0$ $(n\gg0)$である．
任意の$M\in\cH$に対し，AZ生成性を$d$回用いて
\[
 0\to K_d\to P_{d-1}\to\cdots\to P_0\to M\to0
\]
を取り，各$P_j$を有限個の負twist $s^{-a}e$の直和とする．
有限個の$P_j$について高いtwistで$\Ext^{>0}(T,s^nP_j)=0$となる．
次元移動により
$\Ext^k(T,s^nM)\simeq\Ext^{k+d}(T,s^nK_d)=0$ $(k>0,n\gg0)$．
$e$は$T$の直和因子なので$(T,s)$は全対象を生成し，上の消滅はAZの全射条件を与える．

Artin--Zhangの再構成定理\cite[Theorem 4.5]{AZ}を適用する．
ample sequenceを用いた構成と逆関手は\cite[Theorem 2.4]{PolishchukProj}も参照されたい．
右noetherian性に必要な点を明示するため，graded右イデアル$I\subset B_T$を取る．
$I_n$を射$s^{-n}T\to T$と見た像の和は，$T$のnoetherian性から有限個の射の像で生成される部分対象$J$に安定する．
それらの射を$f_i:s^{-a_i}T\to J$とする．AZ全射条件を
$\bigoplus_i s^{-a_i}T\twoheadrightarrow J$へ適用すると，高次数の$I$は
$f_i$の右倍で生成される．低次数部分は有限次元なので，$I$全体は有限生成である．
一般の右イデアルにも次数filtrationの最高次項を用いて持ち上げられ，右noetherian性を得る．
再構成関手は$M\mapsto\pi\bigoplus_n\Hom(T,s^nM)$であり，最後の主張はこの定義から従う．
\end{proof}

\begin{theorem}[規格化共通領域でのAS--反標準一致]\label{thm:normalized-abelian}
第\ref{sec:boundarydata}節の規格化した構成を取り，VdBの導来比較条件，
六点の相異性，$L_1$一般位置，および$(U_+)$を満たすとする．
そこで用いた標準AS実現，すなわちfull very strong周期helixと
そのrolled-up Jacobi環$B(W_+)$の同定の下で，$B(W_+)$は右noetherianであり，
\[
 \boxed{\qgr\mathsf A_+\simeq\qgr F}
\]
となる．この同値は系\ref{cor:normalized-pointed}の指定dg同値$\Psi_+$のheartへの制限である．
従ってこの共通領域では，右noetherian性と十一条件の非消滅を独立に追加する必要はない．
\end{theorem}
\begin{proof}
$\cH=\qgr F$はnoetherian，Hom有限で，Serre関手$(-1)[2]$からhomological dimensionは2である．
標準の$(\pi F,(1))$はAZ-ampleかつ高次Extのtail消滅を満たす．
AS標準helixのvery strong性と周期性は\cite[Definition 3.1, (3.9)--(3.10)]{OU}の入力であり，
pointedな$\Psi_+^{-1}$で運べば，前命題の仮定を満たすhelixを$\Db(\qgr F)$に得る．
そのorbit代数は，指定したAS実現の下で$B(W_+)$である．
前命題から右noetherian性と$qgr$のabel同値を得る．周期coverによって$\qgr\mathsf A_+$と同定する．
この再構成はfoundationとそのEnd代数への作用を保存するため，
full strong generatorのdg Yoneda同定を通じて，誘導する導来同値は指定$\Psi_+$と同じである．
従って$\Psi_+$自体がt-exactとなる．
係数側でも，系\ref{cor:elliptic-discriminant}により$L_1$一般位置から十一条件が成立し，
第\ref{thm:AOUexactlocus}定理の判定と一致する．
\end{proof}
\begin{remark}[非循環性と範囲]
前命題は先にAS標準heartのnoetherian性やt-exactnessを仮定していない．
noetherianな標的$\qgr F$，そのample構造対象，および導来圏内のvery strong helixから導いている．
ただし任意のfull strong collectionをvery strongと読み替えてはならず，
本定理は第\ref{sec:boundarydata}節で用いたAS実現の範囲に限る．
任意のpotential，特異境界，あるいは未変更の印刷AOU座標の全領域への主張ではない．
\end{remark}

\begin{corollary}[整数倍点族の二つの比較]\label{cor:cyclic-AS-equality}
系\ref{cor:normalized-cyclic-family}の整数倍点族では，規格化AS heartは反標準heartと一致する．
一方，同じpointed導来圏内の順序付きRains heartは，
定理\ref{thm:sorted-spectrum}の整列手順によってそのheartへ移る．
特に$(q,3q,2q,5q,7q,11q)$では一回の単純tiltで移り，
反標準側の単純次元は$(6,1,1,2,2,4)$である．
\end{corollary}
\begin{proof}
正整数倍点は$L_1$一般位置と$(U_+)$を満たすので前定理を適用できる．
Rains heartの記述は既存の整列定理による．
ここでは同じ構造対象を保つdg比較によって共通の導来圏に運んでいる．
\end{proof}

\begin{remark}[検算]
\texttt{verify\_four\_hearts\_elliptic.py}は，Weierstrass曲線上のMiller関数を使い，
各一次元Homの切断とその積を有限体上の\emph{実際の有理関数}として構成する．
関係式は関数体で交差乗算して検査し，係数比の逆数と$L_1$射影座標の行・列rescaling類を照合する．
さらに二種類の混合共線壁，最初の三点の共線壁，conic壁を強制した例を別々に検査する．
これは有限体上の有限個の厳密検算であり，複素数上の全パラメータ定理や
noetherian性・heart同値の証明ではない．
\end{remark}

\section{section-model関手の内在的同定と比較図式}\label{sec:canonicalcomparison}
本節では，VdBのsection-model関手と反標準導来切断を，同じ既知のblow-up表示に対して
自然に同定する．単なる生成対象上の次元一致ではなく，次数付きの右作用を保つRHomの同定を用いる．
VdBの関手の存在と境界上の性質は既知の入力であり，その結果を別の新しい存在定理とは呼ばない．

\begin{lemma}[導来切断によるpointed関手の認識]\label{lem:pointedrecognition}
$R$は局所有限な連結次数付き右noetherian代数で，標準のtails構成が
$R_n\simeq\Hom(\pi R,\pi R(n))$ $(n\ge0)$を満たすとする．
$\cD$上の自己同値$s$と対象$e$に対し，dg同値
$G:\cD\simeq\Db(\qgr R)$，同型$G(e)\simeq\pi R$，自然同型$Gs\simeq(1)G$を取る．
これらを整合的に反復して得る次数付き環の同型を
$R(e,s)\simeq R$として固定すると，自然同型
\[
 G(X)\simeq\pi\bigoplus_{n\ge0}\RHom_{\cD}(e,s^nX)
\]
がある．ここでは右$R$作用とdg構造も同定する．
\end{lemma}
\begin{proof}
$\pi:\operatorname{Gr}R\to\QGr R$の右随伴を$\omega$とする．その次数$n$部分は
$\Hom_{\QGr R}(\pi R(-n),-)$である．$\pi$はexactで，counit
$\pi\omega I\to I$は各対象で同型である．従ってinjective resolutionに項別に適用して
$\pi R\omega Y\simeq Y$を得る．非負内部次数へ切断しても，捨てる負次数部分は
右有界なgraded加群としてtailsのtorsionとなるため$\pi$の像は変わらない．
一方，指定dg同値と$s$の整合性から
\[
 \RHom_{\cD}(e,s^nX)
 \simeq\RHom_{\QGr R}(\pi R,G(X)(n))
\]
となる．合成に関する自然性により，この同型は$R(e,s)\simeq R$の下で右$R$作用と両立する．
直和を取りcounitを適用すれば結論を得る．まず$D(\QGr R)$で構成し，
$G(X)$が有界coherentであることから所要の部分圏に入る．
\end{proof}
\begin{remark}
負次数切断に関する主張は，右有界なgraded加群の各元が十分高次数の作用で消えるというtorsionの意味である．
有限次元でない負次数全体を，一括して有限次元加群と呼んではいけない．
また環の同定を固定しなければ，次数付き環の自己同型によるtwistの曖昧さが残る．
$K_0$の一致やorbitのExtの次元だけでは本補題を適用できない．
\end{remark}

\begin{theorem}[VdBの最後の比較関手は反標準導来切断である]\label{thm:VdBcanonical}
VdBの三次曲面構成の最後の段階
$W=\widetilde X_6\xrightarrow{\delta}V=X_7=\operatorname{Proj}F$を取り，
Theorem 11.1.3の一般位置条件とcoherent有界部分への制限条件を満たすとする．
$e_W=\OO_W$，$s_W=\Sfun_W^{-1}[2]$とする．
標的のSerre関手が$(-1)[2]$である設定では，
\[
 R\delta_*e_W\simeq\pi F,\qquad
 R(e_W,s_W)\simeq F,\qquad
 \boxed{R\delta_*\simeq\pi\mathbf R\Gamma_{(e_W,s_W),*}}.
\]
これらは，$R\delta_*$が誘導する環の同定を使った自然同型である．
指定楕円境界へのrestrictionとも両立する．
\end{theorem}
\begin{proof}
VdB \cite[(9.5), Proposition 9.2.3(1)]{VdB}の記号で，blow-upを作るsheaf algebraを
$\mathcal D$，その大域切断を$D$と書くと
$R\Gamma(X,\mathcal D\OO_X)=D$である．従って(9.5)から
$R\delta_*\OO_W\simeq\pi D=\OO_V$となる．
Lemma 11.1.1によりこの入力は当該反復の各段階で適用できる．
Theorem 11.1.3は$R\delta_*$の導来同値を与える．標的のSerre同定
\cite[Lemma 9.2]{AOU}とSerre関手の一意性から
$R\delta_*s_W\simeq(1)R\delta_*$を得る．
$F_n\simeq\Hom(\pi F,\pi F(n))$の飽和は系\ref{cor:AOU-ring}の証明で確認した．
従ってRHomとその合成から$R(e_W,s_W)\simeq F$が従い，
補題\ref{lem:pointedrecognition}が最後の同定を与える．
境界についてはSerre関手の一意性だけを使わず，
VdB \cite[Proposition 9.1.8, (9.11)]{VdB}の
$Li_V^*R\delta_*\simeq R\delta_{Y*}Li_W^*$を用いる．この設定で$\delta_Y$は境界の同定である．
\end{proof}

\begin{corollary}[同じ表示での四者の接続]\label{cor:fourfunctors}
上の$W$を同じblow-upとgluingからRainsの言葉で実現し，
定理\ref{thm:RV}のpointed t-exact同値$J:\Db(\AV)\simeq\Db(\AR)$を取る．
さらに同じ$F$にAOU比較を適用し，実際のrestriction条件\eqref{eq:two-nontrivial}を確認する．このとき
\[
 R\delta_*\simeq\Phi_{\mathrm R}J,\qquad
 \Psi R\delta_*\simeq\Psi\Phi_{\mathrm R}J
\]
が構造対象と反Serre作用を保って成立する．
標準noetherian AS表示で十八係数非零かつ$\Delta\ne0$なら，
$\Psi R\delta_*$によるVdB--AS比較は，
$(\ker R^1\delta_*,\ker\delta_*)$に関する一回のHRS tiltである．
\end{corollary}
\begin{proof}
定理\ref{thm:VdBcanonical}を$J$で運び，Rains側のorbit環を同じ$F$に同定する．
その構造対象からの導来切断が$\Phi_{\mathrm R}$である．
$\Delta\ne0$では定理\ref{thm:AOUexactlocus}により$\Psi$はt-exactだから，
$R\delta_*$のHRS分解を変更せずAS側へ運べる．
\end{proof}
\begin{remark}[同じデータの意味と残る範囲]
この結果は，独立に選んだ比較同値を仮定する代わりに，VdBの指定した最後のraw blow-upと
その同じRains実現について自然関手を同定する．一方，VdBは反復の途中でsection modelを用いる．
任意の可換avatarから一度に作ったgluingや，途中でsection modelを使わない六回のraw blow-upが，
その同じheartになるとは結論しない．その区別を無視したパラメータ辞書は本稿の定理に含めない．
\end{remark}

\begin{proposition}[tiltの差は境界非交叉部分に集中する]\label{prop:boundarydefect}
定理\ref{thm:VdBcanonical}の範囲で，境界非交叉なcoherent対象のSerre部分圏を
$\operatorname{Iso}_Y(W)$，$\operatorname{Iso}_Y(V)$とする．
$\ker\delta_*\subset\operatorname{Iso}_Y(W)$であり，
$R^1\delta_*(M)\in\operatorname{Iso}_Y(V)$である．
また$\delta_*$は境界非交叉部分を除いたabel商圏の同値を誘導する．
\end{proposition}
\begin{proof}
VdB \cite[Theorem 9.1.9]{VdB}は，$\delta_*,\delta^*$が対応するIso部分圏を保ち，
高次導来関手の像がIsoに入り，Isoによる大きいabel商圏の逆同値になることを示す．
Lemma 9.1.5のcoherence保存とnoetherian性によってcoherent商圏へ制限する．
$\delta_*M=0$なら商圏で$M$も零なので，最初の包含も従う．
\end{proof}
\begin{remark}[既知のexceptional objectsとの関係]
VdB \cite[Theorem 11.2.1, Lemma 11.2.4]{VdB}は既に境界非交叉対象の有限長性と，
楕円点配置による単純対象数の公式を与えている．従って有限長性や個数の存在自体を本研究の新規性としない．
第\ref{sec:rootchambers}節では，反標準Euler標数からtorsion pairと単純対象の根類を選ぶ方法を記述する．
\end{remark}

\section{最大共鳴のZariski稠密性}\label{sec:maximal-density}
本節は次数$1$から$6$までを一様に扱う．単純根の共鳴軌跡の稠密性と
有効性の機構はRainsの既知結果\cite[\S4, pp.~34--35]{Filtered}である．
ここでは全根が\emph{同時に}共鳴する点を整数で表示し，
同じ楕円chartでheart一致・不一致の双方が稠密になることを明記する．
数値的な根の共鳴を，全ての根に対する既約sheafの実在と同一視しない．

\subsection{全根共鳴の整数表示}
滑らかな平面三次$E$の変曲点$o$を原点とし，$q\in E$を非torsionとする．
点$q$に対応する次数零線束も$q$と書く．$3\le r\le8$，$d=9-r$とし，
相異なる$p_1,\ldots,p_r\in E$のblow-up $S$を考える．
$h|_E=\OO_E(3o)$，$H=3h-\sum e_i$であり，
\[
 \rho\left(ah-\sum b_ie_i\right)=q^{-m(\alpha)}
 \quad\Longleftrightarrow\quad \sum b_ip_i=m(\alpha)q
\]
を共鳴の符号規約とする．$q$の非torsion性から指数は存在すれば一意である．
根格子の単純根は
\[
 \alpha_0=h-e_1-e_2-e_3,\qquad
 \alpha_i=e_i-e_{i+1}\quad(1\le i<r)
\]
である．

\begin{proposition}[最大共鳴の九価の整数パラメータ]\label{prop:maximal-integral-chart}
全根が$\langle q\rangle$にrestrictionを持つことは，全単純根についてそうであることと同値で，
整数$m_0,\ldots,m_{r-1}$による
\begin{equation}\label{eq:maximal-integral-chart}
 3p_1=(m_0-2m_1-m_2)q,\qquad
 p_i=p_1+\left(\sum_{j=1}^{i-1}m_j\right)q\quad(2\le i\le r)
\end{equation}
で記述される．指数は$m(\alpha_i)=m_i$であり，整数列を固定した式の解は，相異性・genuine性を課す前には
$E[3]$のtorsorで九通りである．相異なる点の範囲では，avatarがgenuineであることと
全根に対する$m(\alpha)\ne0$は同値である．
\end{proposition}
\begin{proof}
単純根が根格子の基底であることから最初の同値を得る．
$e_i-e_{i+1}$のrestrictionは$p_i-p_{i+1}$なので
$p_{i+1}=p_i+m_iq$であり，$p_1+p_2+p_3=m_0q$と合わせれば表示式になる．
複素楕円曲線の$[3]$は全射で核が九個だから，選択の主張も従う．

滑らかな反標準曲線の真の変換も$E$と書く．$H=[E]$はnefで$H^2=d>0$である．
$H$がampleでなければ，Nakai--MoishezonとHodge指数定理から
$HC=0$を満たす既約曲線$C$があり，adjunctionから$C^2=-2$となる．
$C\cap E=\varnothing$なので$\rho(C)=\OO_E$，従って$m(C)=0$である．
逆に$H$がampleなら，任意の根$\alpha$について$H^0(S,\pm\alpha)=0$であり，
$\chi(\alpha-H)=0$とSerre双対性から$H^1(S,\alpha-H)=0$となる．
restriction列より$H^0(E,\rho(\alpha))=0$，従って$\rho(\alpha)\ne\OO_E$である．
非torsion性によりこれは$m(\alpha)\ne0$と同値である．
\end{proof}

\begin{theorem}[全次数の最大共鳴は稠密]\label{thm:maximal-density-main}
相異なる正整数$N_1,\ldots,N_r$に対し$p_i=N_iq$と置く．
この配置のavatarはgenuine del Pezzoで，全ての根は非零整数指数で共鳴する．
さらに，任意の置換$\pi\in\mathfrak S_r$について
\[
 \{(N_1q,\ldots,N_rq):0<N_{\pi(1)}<\cdots<N_{\pi(r)}\}
\]
は$E^r$にZariski稠密である．従って最大共鳴のgenuine locusは稠密であり，
特に$N_1<N_3<N_2$という正負混合条件を課しても稠密である．
\end{theorem}
\begin{proof}
根を$ah-\sum b_ie_i$と書くと
\[
 \sum b_i=3a,\qquad \sum b_i^2=a^2+2.
\]
$a\ge0$の根は，差根$e_i-e_j$，三点を引く直線類，六点を引くconic類，
および$r=8$での$3h-2e_i-\sum_{j\ne i}e_j$に尽きる．
実際Cauchy--Schwarzから$(9-r)a^2\le2r$，
整数$b$に対する$b(b-1)\ge0$から
$\sum b_i(b_i-1)=a^2-3a+2$を得る．
$a=1,2$では全係数が$0,1$，$a=3$では$r=8$で一つが$2$，残りが$1$となる．
唯一残る$a=4,r=8$では，右辺$6$を一つの$3$または$-2$，あるいは
三つの$2$または$-1$で実現しても$\sum b_i=12$に届かず，不可能である．
$a=0$では一つずつの$1,-1$だけである．

従って$m(\alpha)=\sum b_iN_i$は，差根なら相異なる二整数の差，
$a>0$なら正整数の正係数和であり，零にならない．
点の相異性と前命題によりgenuine性を得る．

$D=\{nq:n\ge1\}$は曲線$E$の無限集合なので稠密で，$D^r$も$E^r$に稠密である．
正整数$u_j$を用いて$N_{\pi(j)}=u_1+\cdots+u_j$と書く．
この累積和と置換で定まる$E^r$の群自己同型は整数行列式$\pm1$を持ち，
$D^r$を表示した集合へ送る．よってその集合は稠密である．
$N_1<N_3<N_2$を満たす全順序を一つ選べば最後の主張も従う．
\end{proof}

\begin{center}
\begin{tabular}{c|rrrrrr}
$d$&6&5&4&3&2&1\\\hline
根系&$A_2+A_1$&$A_4$&$D_5$&$E_6$&$E_7$&$E_8$\\
全根数&8&20&40&72&126&240
\end{tabular}
\end{center}
この表は\emph{数値的}共鳴根の数である．同じheartの実在既約根sheafが
この個数あるという意味ではない．整数倍点族では，指定blowdownの$r$個の単純根への
Rainsの存在結果を用い，後述の定理\ref{thm:sorted-spectrum}により全単純とその像を決定する．

\subsection{指定heartの一致条件は構成可能でない}
$U_r\subset E^r$を，相異なる点をblow-upしたavatarがgenuineとなる開集合とする．
固定した$q$と指定順序のRains heartsを各幾何学的点で取る．

\begin{corollary}[一致・不一致の双方の稠密性]\label{cor:nonconstructible-texact}
$U_r$の複素点上で，指定反標準比較がt-exactとなる点集合を$U_r^{\rm ex}$とする．
$U_r^{\rm ex}$と$U_r\setminus U_r^{\rm ex}$はともにZariski稠密で，
いずれも$U_r$の構成可能部分集合の複素点集合ではない．
\end{corollary}
\begin{proof}
定理\ref{thm:sorted-spectrum}は，整数倍点族の指定Rains heartについて，
比較がt-exactであることと$N_1<\cdots<N_r$が同値であることを示す．
昇順の整数配置は前定理により稠密である．
一方$N_1<N_3<N_2$を課した稠密集合では$m(\alpha_2)=N_3-N_2<0$で，
同じ定理，または実在単純根による判定からt-exactでない．
既約noetherian空間の稠密な構成可能部分集合は非空開集合を含む．
一方の集合がそうであれば，他方の稠密性に反する．
\end{proof}

\begin{remark}[稠密性の範囲]
これは固定した楕円chart上の\emph{pointwiseな指定heart}の判定である．
相対的t構造や一様な比較kernelの構成を主張していないので，
そのような相対仮定を備えたt-exactnessの開性定理とは矛盾しない．
また全非可換del Pezzoモジュライの全既約成分に関する稠密性ではない．
共鳴で反標準環のHilbert関数が跳ぶという主張でもない．
$R/gR$が次数$d$の楕円座標環なら，中心元の完全列から依然として
$\dim R_n=1+dn(n+1)/2$である．
\end{remark}

\section{零指数根壁：古典的収縮と非可換gluingの核}\label{sec:zero-walls}
非零指数の共鳴と，$\rho(\alpha)=\OO_E$という零指数の壁は異なる．
前者はgenuineなavatarの内部にもあり，既存の導来同値の下でheartが変わる．
後者では反標準orbitから見えない非零対象が存在し得るため，
導来同値を前提とするHRSの議論をそのまま延長してはならない．
本節は古典的に知られた収縮を記載した後，標準的な弱楕円gluingについて
\emph{orbitから見えない全対象}を同定する．反標準$qgr$との商圏比較には
なお追加の有限性・局所化の議論が必要であり，未解決部分を末尾に分離する．

\subsection{十一因子と古典的な収縮}
$P=(Q^{\langle-1\rangle})^{\mathsf t}$の列を$p_1,p_2,p_3$，
座標点を$p_4,p_5,p_6$とする．十八係数非零のchartでは，
九つの$2\times2$ minorの消滅は，可動点二つと補集合の座標点一つの共線性を表す．
残る$\det Q^{\langle-1\rangle}=0$は$p_1,p_2,p_3$の共線性，
$\det Q=0$は六点conic条件である．対応する根はそれぞれ
\[
 h-e_i-e_j-e_k\quad\text{（十種類）},\qquad
 2h-\sum_{i=1}^6e_i
\]
であり，いずれも$\alpha^2=-2$，$H\alpha=0$を満たす．
これは第\ref{sec:explicitcorner}節の六点復元を壁で読むことで従う．
各divisorの一般点では対応曲線の真の変換$C$が唯一の$(-2)$曲線となる．

\begin{theorem}[古典的根壁の不可視対象]\label{thm:zero-classical-contraction}
滑らかな弱del Pezzo曲面$S$，$H=-K_S$，$H^2=d>0$について，
反標準収縮を$f:S\to\overline S$とする．$i:C\hookrightarrow S$が$(-2)$曲線なら
$K_C=i_*\OO_C(-1)$は非零の$2$-spherical対象で，
\[
 K_C(H)\simeq K_C,\qquad
 \RHom_S(\OO_S(nH),K_C)=0\quad(n\in\mathbb Z),\qquad
 Rf_*K_C=0.
\]
特に反標準pairは$\coh S$でAZ-ampleでない．$S$がgenuineでなければ
\[
 \Db(\coh S)\not\simeq
 \Db\left(\qgr\bigoplus_{n\ge0}H^0(S,nH)\right).
\]
三次曲面の一つの根壁の一般点では，$\overline S$は一つの$A_1$特異点を持つ．
\end{theorem}
\begin{proof}
弱del Pezzoの反標準モデルは$(-2)$曲線の各連結配置をDu Val特異点へ収縮し，
他の場所では同型である\cite[Theorem 8.3.2, \S8.3.2]{DolgachevWalls}．
$C\simeq\mathbb P^1$，$N_{C/S}=\OO_C(-2)$であり，局所Ext層は
$\mathcal Ext^0_S(K_C,K_C)=\OO_C$，$\mathcal Ext^1_S(K_C,K_C)=\OO_C(-2)$．
局所大域スペクトル列から自己Extは次数$0,2$で$k$，他は零である．
$K_S|_C=\OO_C$なのでSerre関手は$K_C$を$K_C[2]$へ送り，spherical性を得る．
$H|_C=\OO_C$と$R\Gamma(C,\OO_C(-1))=0$が二つの消滅を与える．
非零$K_C$への全反標準twistからのHomが零なのでAZ生成性は失敗する．

反標準環の$qgr$は$\coh\overline S$である（低次数では十分なVeroneseを用いる）．
滑らかな$S$の任意の有界coherent複体間では十分高次のHomが零である．
一方，$\overline S$の特異点$x$の剰余体は，非正則局所環上の最小自由解消が
無限に続くので，$\Ext^j(k(x),k(x))\ne0$となる$j$が無限にある．
支持が一点であるためこれは大域Extでも同じで，三角圏同値と矛盾する．
一壁の一般点では例外配置が一本なので型は$A_1$である．
\end{proof}

\begin{remark}[正負指数の完全相殺の古典模型]
$C\simeq\mathbb P^1$のEuler完全列を押し出すと
\[
0\longrightarrow i_*\OO_C(-2)\longrightarrow K_C^{\oplus2}
 \longrightarrow i_*\OO_C\longrightarrow0
\]
となる．両端のEuler値は$-1,+1$であり，connecting map
$H^0(C,\OO_C)\to H^1(C,\OO_C(-2))$は同型である．
中間項は非零だが全反標準orbitから見えない．これは零指数壁での完全相殺を示す．
両端は同じ曲線類を持ち，非torsionのgenuine gluingにおける異なる隣接根の
正負指数の組と同一視しない．
\end{remark}

\begin{example}[有理数上の一壁と反標準方程式]\label{ex:zero-rational-wall}
\[
 Q_0=\begin{pmatrix}1&1&1\\1&2&2\\1&3&4\end{pmatrix}
\]
では行$(0,1)$・列$(1,2)$のminorだけが零で，
$\det Q_0=1$，$\det Q_0^{\langle-1\rangle}=1/24$である．
六点は
\[
 (1:1:1),(2:1:1),(12:4:3),(1:0:0),(0:1:0),(0:0:1)
\]
となり，共線三点は$p_1,p_2,p_4$だけで，直線は$y=z$である．
反標準切断の基底を
\[
\begin{aligned}
 A&=\tfrac13x^2(y-z)-xy(y-z),&
 B&=\tfrac7{12}x^2(y-z)-x(y^2-z^2),\\
 C&=\tfrac16x^2(y+2z)-\tfrac32xy^2+y^2z,&
 D&=\tfrac14x^2(y+z)-\tfrac32xy^2+yz^2
\end{aligned}
\]
と取ると，像の三次方程式は
\[
\begin{aligned}
 0={}&23A^3-22A^2B+44A^2C-42A^2D+8AB^2-16ABC\\
 &-24AC^2+24AD^2+8B^2C+24BC^2-24BCD.
\end{aligned}
\]
$y=z$は$[0:0:1:1]$へ移り，$C=1$の局所座標$(A,B,D-1)$で
Hessian行列式は$-2304$である．従ってこの点は通常二重点である．
これらは有理数係数の恒等式で検査できる．ここまでは可換曲面の計算であり，
このpotentialの全非可換変形について同じ特異点型を主張していない．
\end{example}

\subsection{任意の実在零指数根sheafに対する障害}
\begin{proposition}[零指数の条件付き消失判定]\label{prop:zero-sheaf-obstruction}
Rainsの境界triangle，曲面のhomological dimension $2$，
line bundleから一次元sheafへの$\Ext^2$消滅と
構造line bundleの商に対するRiemann--Roch不等式が成立するheartを取る．
境界非交叉の既約純一次元sheaf $M$が実在し，
$c_1(M)^2=-2$，$Hc_1(M)=0$，$\chi(e,M)=0$とする．
すると全$n\in\mathbb Z$で$\RHom(e,s^nM)=0$である．
自然反標準導来切断が定義される範囲では$\Phi M=0$であり，保守的でない．
\end{proposition}
\begin{proof}
境界triangleから$sM\simeq M$である．非零$e\to M$の像を$J$とすると，
既約性から$c_1(J)=c_1(M)=\alpha$，余核は零次元である．
Rains\cite[Corollary 10.16]{Rains}の不等式は
$\chi(e,J)\ge-\alpha(\alpha+K)/2=1$を与え，
零次元余核のEuler値は非負なので$\chi(e,M)=0$に反する．
従って$\Hom(e,M)=0$．仮定した$\Ext^2$消滅とEuler値から$\Ext^1(e,M)=0$で，
他の次数も消える．全twistは同じであり，導来切断複体そのものが零となる．
\end{proof}
ここでは\emph{実在}が仮定である．任意の数値的零指数根をこの命題でsheafへ置き換えたり，
標準heartの変化を単なるtiltと呼んだりしない．次の結果ではsheafの実在を仮定せず，
古典的な収縮対象をgluingの共通部分圏に直接埋め込む．

\subsection{弱楕円gluingにおける全orbit核の不変性}
$S$を滑らかな射影有理弱del Pezzo曲面，$E\in|H|$を滑らかな楕円曲線とする．
$\cB=\OO_S^\perp\subset\Perf(S)$とし，$q\in\Pic^0(E)$について
\begin{equation}\label{eq:weak-wall-gluing}
 \RHom(e,M)=0,\qquad
 \RHom(M,e)=\RHom_E(Li^*M,q)\quad(M\in\cB)
\end{equation}
による標準Rains gluingを$\cD_q$，共通部分圏の包含を$\iota_q:\cB\hookrightarrow\cD_q$と書く．
この構成はgenuine性を必要としない\cite[\S4, (4.4)--(4.7)]{Filtered}．
境界関手$F_q$，右随伴$i_{q*}$，$s_q=\mathbb S_q^{-1}[2]$は
\begin{equation}\label{eq:weak-wall-boundary}
 F_qe=\OO_E,\quad F_q\iota_qG=Li^*G\otimes q^{-1},\quad
 s_q^{-1}X\longrightarrow X\longrightarrow i_{q*}F_qX\longrightarrow,
 \qquad F_qs_q=\Psi_qF_q
\end{equation}
を満たす．$\Psi_q$は平行移動と次数$d=H^2>0$の線束によるtensorからなる．
これらの\emph{実際の}dg gluingと境界関手を固定し，任意の数値的なavatarへ一般化しない．

反標準収縮$f:S\to\overline S$について，充満dg部分圏を
\[
 \mathcal N_f=\{G\in\Perf(S):Rf_*G=0\},\qquad
 \mathcal N_q=\{X\in\cD_q:\RHom(s_q^ne,X)=0\ (n\in\mathbb Z)\}
\]
と定める．$\mathcal N_f\subset\cB$である．
また$\mathcal N_q^{\rm tail}$を，各$j$について
$\Hom(e,s_q^nX[j])=0$が全ての十分大きい$n$で成り立つ対象の部分圏とする．

\begin{theorem}[全orbit核と古典的収縮核の同定]\label{thm:weak-orbit-kernel}
上の弱楕円gluingでは，$q$の位数によらず
\begin{equation}\label{eq:weak-orbit-kernel}
 \boxed{\mathcal N_q^{\rm tail}=\mathcal N_q=\iota_q(\mathcal N_f).}
\end{equation}
従ってorbitから見えない対象の\emph{充満dg圏}は，この標準gluingの$q$に依存しない．
特に$\mathcal N_q=0$であることと，$S$がgenuineであることは同値である．
\end{theorem}
\begin{proof}
\emph{tailから境界零へ．}
$X\in\mathcal N_q^{\rm tail}$なら$s_q^{-1}X$も同じ条件を満たす．
境界triangleへ$\RHom(s_q^{-n}e,-)$を施し，随伴を使うと
\[
 \Hom_E(\Psi_q^{-n}\OO_E,F_qX[j])=0\qquad(n\gg0)
\]
が各$j$について従う．$\Psi_q^n$は正次数のtwisted line bundleの作用である．
楕円曲線のcohomologyと十分高い次数の線束による消滅・生成性から，
この条件は有界対象$F_qX$が零であることを意味する．
より具体的には，その有限個のcohomology sheafへのhyper-Ext spectral sequenceを
$n\gg0$で退化させ，非零sheafの高い正twistは非零の大域切断を持つことを用いる．
ゆえに$s_q^{-1}X\to X$は同型．周期性からtail消滅は
$\RHom(e,X)=0$と全整数orbit消滅を与える．
$e^\perp=\iota_q(\cB)$なので$X=\iota_qG$と書け，
$R\Gamma(S,G)=0$かつ$Li^*G=0$となる．

\emph{古典的収縮の核へ．}
$Li^*G=0$ならCartier divisorのtriangleから
$G(-H)\to G$が同型で，全整数$n$について$R\Gamma(S,G(nH))=0$である．
十分可除な$a>0$と$\overline S$上のample line bundle$L$を取り，
$f^*L=\OO_S(aH)$とする．射影公式により
\[
 R\Gamma\bigl(\overline S,Rf_*G\otimes L^n\bigr)=0\quad(n\ge0).
\]
$Rf_*G$は有界coherentであり，Serre消滅と大域生成を各cohomologyへ適用すれば
$Rf_*G=0$を得る．従って$\mathcal N_q^{\rm tail}\subset\iota_q(\mathcal N_f)$．

\emph{逆包含．}
$G\in\mathcal N_f$なら$R\Gamma(S,G)=0$で，$G\in\cB$である．
$f$は例外曲線の外で同型であり，それらは$E$と交わらない．
$Rf_*G=0$をこの同型な開集合へ制限すると$G$はそこでも零なので，$Li^*G=0$である．
式\eqref{eq:weak-wall-boundary}と$\RHom(e,\iota_qG)=0$から
$s_q\iota_qG\simeq\iota_qG$と全orbit消滅を得る．
よって$\iota_q(\mathcal N_f)\subset\mathcal N_q\subset\mathcal N_q^{\rm tail}$．
$\iota_q$はdg完全忠実なので，等式は全RHomと合成を含む同定である．
最後にgenuineなら$f$は同型で核は零．そうでなければ
定理\ref{thm:zero-classical-contraction}の$K_C\ne0$が核にある．
\end{proof}

\begin{corollary}[実在する非可換零指数spherical対象]\label{cor:zero-spherical-persistence}
$S$の任意の$(-2)$曲線$C$について$K_{C,q}=\iota_q(i_*\OO_C(-1))$は
$\cD_q$の非零$2$-spherical対象で，$s_qK_{C,q}\simeq K_{C,q}$，
$\chi(e,K_{C,q})=0$，$K_{C,q}\in\mathcal N_q$である．
複数曲線の直和に対する導来End代数も$q$によらず可換側のものと同じである．
\end{corollary}
\begin{proof}
$K_C\in\mathcal N_f\subset\cB$と完全忠実性により自己Extと全合成は保存される．
境界triangleから$\mathbb S_qK_{C,q}\simeq K_{C,q}[2]$でありspherical性を得る．
Euler値と残りは前定理による．
\end{proof}
これは\emph{対象の構成}であって，任意のblowdown表示のheartでそれがsheaf，
あるいは単純sheafになることまで証明したものではない．
そのheart内の位置の同定を必要とせず，指定反標準orbitの生成障害が全$q$で残る．

\begin{proposition}[coherent導来切断に対する正確な核]\label{prop:weak-coherent-kernel}
さらに$R_q=\bigoplus_{n\ge0}\Hom(e,s_q^ne)$が局所有限右noetherianで，
自然dg関手$\Phi_q=\pi R\Gamma_*:\cD_q\to\Db(\qgr R_q)$が定義されるとする．
各$X$について
\[
 N^j(X)=\bigoplus_{n\ge0}\Hom(e,s_q^nX[j])
\]
が有限生成で，非零となる$j$は有限個，かつ$H^j\Phi_qX=\pi N^j(X)$と仮定する．
すると$\ker\Phi_q=\mathcal N_q=\iota_q(\mathcal N_f)$である．
\end{proposition}
\begin{proof}
有限生成・局所有限な次数付き加群では，$\pi N^j=0$は$N^j$が有限次元であること，
従って十分高い内部次数の成分が零であることと同値である．
$\Phi_qX=0$と$X\in\mathcal N_q^{\rm tail}$の同値を得て，前定理を使う．
\end{proof}
この命題のcoherent有限性は\emph{追加仮定}である．
全orbit核の同定から$qgr$値関手の有界性・有限生成性を逆に推論しない．
ただし$K_{C,q}$の生の導来切断は全次数で既に零なので，それが消えること自体に
全対象についての有限生成性は不要である．

\subsection{dg商による圏論的な収縮候補}
\begin{proposition}[orbit環を変えない圏論的収縮]\label{prop:zero-dg-quotient}
標準弱楕円gluingに対し，pretriangulated dg商の冪等完備化を
\[
 \mathcal Q_q=(\cD_q/\mathcal N_q)^{\natural},\qquad
 \ell_q:\cD_q\longrightarrow\mathcal Q_q
\]
とする．$s_q$は商上の自己同値$\overline s_q$へ降り，$\overline e=\ell_q e$はexceptionalである．
全整数$a,b,j$について自然な写像
\[
 \Hom_{\cD_q}(s_q^ae,s_q^be[j])
 \xrightarrow{\sim}
 \Hom_{\mathcal Q_q}(\overline s_q^a\overline e,
                     \overline s_q^b\overline e[j])
\]
は同型で，特に反標準orbit環は$R_q$のままである．
また$\mathcal Q_q$は，$(\cB/\mathcal N_f)^{\natural}$に$\overline e$を
$\RHom_E(Li^*(-),q)$の降下でglueした圏として記述できる．
\end{proposition}
\begin{proof}
$\mathcal N_q$はthickで$s_q$不変なので自己同値は降りる．
全orbit対象は$\mathcal N_q$に右直交し，Serre双対性により左にも直交する．
実際$N\in\mathcal N_q$に対して
$\RHom(N,s_q^be)=\RHom(s_q^be,s_q^{-1}N[2])^*=0$である．
従ってVerdier局所化はこれらの対象間の全Homを変えない．
dg商のHomを表すcone公式\cite[\S1.3.1]{DrinfeldWalls}でも，
$\mathcal N_q$を経由する導来tensor項が零となることから同じ結論を得る．
冪等完備化は既存の対象間のHomを変えない．合成との両立性が環の同定を与える．

さらに$\mathcal N_f\subset\cB$で，gluing moduleはその上で零である．
従ってこのmoduleはdg商$\cB/\mathcal N_f$へ降りる．
元の圏の二成分からなるgluingにdg商の普遍性を適用すると，
$\cB$側をこの商へ置き換え，exceptional成分をそのまま残したgluingを得る．
双方のpretriangulated hullと冪等完備化を取れば表示した圏となる．
\end{proof}
この構成は圏論的な収縮候補を与えるが，$\mathcal Q_q$に目的のnoetherian heartが
存在することや，それが$\Db(\qgr R_q)$であることは主張しない．
また$\overline s_q$が商のSerre関手の逆から生じることも，ここでは結論しない．
標的の特異性・proper性・coherent対象を別に扱う必要がある．

\subsection{反標準模型から見えないtwistとpointed比較の障害}
\begin{corollary}[零指数twistの不可視性]\label{cor:zero-twist-invisible}
$K=K_{C,q}$のspherical twistを$T_K$とする．
$T_K(e)\simeq e$で，構造対象と反Serre作用の整合する同定により
$R_q$上の作用は恒等となる．また$\Phi_q$が定義され$\Phi_qK=0$なら
\[
 \Phi_qT_K\simeq\Phi_q.
\]
ADE配置の隣接曲線に対応するtwistはbraid関係を，非隣接のtwistは交換関係を満たす．
\end{corollary}
\begin{proof}
$\RHom(K,e)=0$はgluing式と$Li^*K_C=0$から従い，
Serre双対性と全orbit直交性から$\RHom(K,s_q^ne)=0$も全$n$で成り立つ．
従ってtwistの評価triangle
\[
 \RHom(K,X)\otimes K\longrightarrow X\longrightarrow T_KX\longrightarrow
\]
の自然な第二射は，全orbit対象で同型であり，全ての射とその合成を保つ．
これが環の作用の主張を与える．$\Phi_qK=0$なら同じtriangleを$\Phi_q$で送ればよい．
異なる例外曲線のExtは交点一つに対し次数1で$k$，非交叉なら全て零である．
既知のspherical twist定理\cite[Propositions 2.10, 2.12--2.13]{SeidelThomasWalls}を適用する．
\end{proof}

\begin{proposition}[ampleな標準模型へのpointed同値はない]\label{prop:zero-pointed-obstruction}
$S$がgenuineでないとする．有界heart$\mathcal H\subset\mathcal D'$に
$(e',s')$がAZ-ampleで，$s'$はheartを両方向で保ち，
各$M\in\mathcal H$について全正次数Extが$n\gg0$で一様に消えるとする．
このとき$(\cD_q,e,s_q)$から$(\mathcal D',e',s')$へのpointed導来同値は存在しない．
\end{proposition}
\begin{proof}
非零$K_{C,q}$は全orbitに直交する．pointed同値があれば，その像も標的の全orbitに直交する．
一方，標的の任意の有界対象$Y$の有限個のcohomologyへ高次Ext消滅を同時に適用すると，
\[
 \Hom(e',s'^nY[j])=\Hom(e',s'^nH^j_{\mathcal H}Y)\quad(n\gg0).
\]
AZ生成性は，非零sheafへの全ての十分高いtwistのHomが零になることを排除する．
従って全orbit直交対象は零のみで，矛盾する．
\end{proof}
構造対象を忘れた任意の抽象的導来同値の不存在は，この命題からは従わない．
可換側の非同値は別途，定理\ref{thm:zero-classical-contraction}で証明した．

\subsection{非torsion gluing上の一つの零指数壁}
\begin{example}[単一の零根を持つ明示的な弱avatar]\label{ex:zero-elliptic-wall}
$E=\mathbb C/(\mathbb Z+i\mathbb Z)$，$L_1=\OO_E(3o)$とし，
\[
 \begin{gathered}
 q=\sqrt2,\quad x=\sqrt3,\quad y=\sqrt5,\quad z=\sqrt7,
 \quad w=\sqrt{11},\quad v=\sqrt{13},\\
 (p_1,\ldots,p_6)=(x,y,-x-y,z,w,v).
 \end{gathered}
\]
$L_0=L_1q$とする．$1,\sqrt2,\sqrt3,\sqrt5,\sqrt7,\sqrt{11},\sqrt{13}$の
$\mathbb Q$線型独立性から，点は相異なり，共鳴する根は
$\pm\alpha_0$，$\alpha_0=h-e_1-e_2-e_3$だけで，その指数はいずれも零である．
実際$\rho(ah-\sum b_ie_i)q^m=\OO_E$から
$b_1=b_2=b_3=b$，$b_4=b_5=b_6=0$，$m=0$となり，根の二方程式から$a=b=\pm1$を得る．
前半三点だけが共線で，対応する直線の真の変換が唯一の$(-2)$曲線$C$である．
$H$はnefで$H^2=3$だから$S$は弱del Pezzo．定理\ref{thm:weak-orbit-kernel}と
系\ref{cor:zero-spherical-persistence}が，この非torsion gluingで非零の零指数対象$K_{C,q}$を構成する．

楕円線束の条件$(U_+)$はなお成立する．実際$X=0$，$Y=z+w+v$なので
$u=-q$，$t=-Y$，$u-t=Y-q$は全て非零である．
実際の楕円decupleのtheta表示では27係数が全て非零である．
実際$a=0,2$では引数に$z,w,v$の一つの係数が残り，$a=1$では残り二つの係数が残る．
入力六点からの$Q$は
九minorと$\det Q$が非零，$\det Q^{\langle-1\rangle}=0$となる．
従って純代数的な十一条件の定理はこのcornerの無限次元性を与える．
ここで使用するtheta式と六点復元は楕円線束の計算であり，genuine性を使わない．
一方，このpotentialのAS模型と$\cD_q$との指定導来比較はここでは構成していない．
一般位置外のVdB比較を適用したり，corner無限性だけからその模型の特異点型を決めたりしない．
\end{example}

\subsection{証明した範囲と将来の課題}\label{subsec:zero-wall-future}
最大共鳴の稠密性はgenuineな点の定理であり，零指数壁の消失と矛盾しない．
後者について本節で得たのは，任意次数の古典的な収縮障害，標準弱楕円gluingでの
全orbit核の$q$不変性，実在spherical対象，orbit環を保存するdg商，
twistの不可視性とpointed同値の障害である．
coherent反標準導来切断の全核との一致は命題\ref{prop:weak-coherent-kernel}の追加仮定付きである．
既知のgluing・収縮・spherical twist理論を入力とするこれらの組合せの網羅的新規性は未確認である．

\paragraph{coherent収縮と商圏．}
任意の弱楕円gluingについて$N^j(X)$の有限生成性と有界性を証明し，
自然関手が全核を除いた商から標的へ完全忠実・本質的全射になるかを調べる．
\emph{核が同定できたことだけではVerdier商と$\Db(\qgr R_q)$の同値は従わない．}
perfect対象の商，coherent有界導来圏，冪等完備化を区別する必要がある．
一つの$A_1$壁から始め，特異な標的の非perfect対象がどのように現れるかを点検する．

\paragraph{heartと局所特異点．}
$K_{C,q}$の指定Rains heartでの位置，零指数根配置全体に適合する
相対的perverse heart，および収縮後の局所・完成代数を構成する．
既存の非可換収縮環\cite{DonovanWemyssWalls}は\emph{与えられた幾何学的収縮}から出発するため，
その定理だけで本節の非可換曲面に収縮射が存在するとは結論しない．
ADEの古典配置が残ることから標的の特異点型や全有限次元表現を直ちに推論しない．

\paragraph{AS模型・壁の交叉・族．}
例\ref{ex:zero-elliptic-wall}の実際の楕円potentialについて，AS実現，noetherian性，
標準helix，pointed比較の各入力を壁上で確認する．
その後に複数零指数壁の交叉，特異楕円境界，曲面パラメータに関する底変更と降下へ進む．
本節の$q$不変性は，\emph{固定した$(S,E)$のgluing}におけるdg部分圏の同定であり，
任意の非可換del Pezzoの全モジュライ上の相対収縮定理ではない．
\endinput

\section{混合\texorpdfstring{$A_2$}{A2}によるheartのずれの具体化}\label{sec:A2}
この節は定理\ref{thm:derived}のgenuine楕円gluingで，$q$非torsionとする．
$H=-K$はnef，$H^2>0$．楕円境界と交わらない一次元以下の対象を
$\cZ=\{M\in\AR:\dim M\le1,\ F_qM=0\}$と書く．

\subsection{単純根sheafと有限長部分圏}
\begin{lemma}\label{lem:offboundary}
$\cZ$は有限長Serre部分圏であり，その既約純一次元root sheafは$\AR$の単純対象である．
\end{lemma}
\begin{proof}
境界非交叉な対象の零次元商$Z$には$i^*Z=0$が従い，$\theta Z\to Z$は全射となる．
両者の数値類は等しいので\cite[Corollary 7.7]{Rains}から同型である．
$[Z]=\ell[\mathrm{pt}]$，$\ell\ge0$に対して境界determinantは$q^{\ell}$であり，
非torsion性から$\ell=0$，従って$Z=0$．既約純一次元sheafの非零部分対象の商は
零次元なので，この場合には単純性が得られる．

$0\to N\to M\to Q\to0$，$M\in\cZ$とすると，nef性と加法性から
$Hc_1(N)=Hc_1(Q)=0$．$i^*Q=0$なので$\theta Q\to Q$は全射であり，
rank零，$Kc_1(Q)=0$と\cite[Corollary 7.15]{Rains}から数値類は変わらない．
Corollary 7.7と自然性により$\theta Q\to Q$，$\theta N\to N$はいずれも同型で，
$N,Q\in\cZ$となる．拡大閉性は境界triangleから従う．
真の降鎖の逐次商は零次元になれないため，\cite[Proposition 7.32]{Rains}の
finitely partitive性が降鎖を止める．coherent heartのnoetherian性と合わせて有限長となる．
\end{proof}

\begin{corollary}[全境界外圏のtorsion pair]\label{cor:Ztilt}
定理\ref{thm:globaltilt}のtorsion pairは$\cZ$へ制限され，
\[
\cT\cap\cZ=\ker\Ext^1(e,-),\qquad
\cF\cap\cZ=\ker\Hom(e,-)
\]
となる．これは二根部分圏だけでなく全$\cZ$についての結論である．
\end{corollary}
\begin{proof}
$\cZ$がSerre部分圏なので，その対象のtorsion分解の部分対象と商も$\cZ$に属する．
$M\in\cZ$では$g$が可逆に作用し，各$\Ext^j(e,s^nM)$の移行写像が同型となる．
その局所化を$V^j(M)$とすれば$\dim V^j(M)=\dim\Ext^j(e,M)$であり，
$H^j\Phi(M)=0$と$\Ext^j(e,M)=0$は同値である．
\end{proof}

異なる二根sheaf$S_\alpha,S_\beta\in\cZ$には
\begin{equation}\label{eq:rootext}
\Hom(S_\alpha,S_\beta)=\Ext^2(S_\alpha,S_\beta)=0,
\qquad \dim\Ext^1(S_\alpha,S_\beta)=\alpha\cdot\beta
\end{equation}
が成り立つ．実際，単純性からHomが消え，境界非交叉とSerre双対性からExt$^2$も消える．
rank-zero Euler形式$\chi(S_\alpha,S_\beta)=-\alpha\beta$が残る次元を与える．

隣接根$\alpha\beta=1$，$\chi(e,S_+)=p>0$，$\chi(e,S_-)=-n<0$を固定し，
両向きの非分裂拡大を
\begin{equation}\label{eq:PQ}
0\to S_-\to P\to S_+\to0,
\qquad 0\to S_+\to Q\to S_-\to0
\end{equation}
とする．存在は\eqref{eq:rootext}から従い，各中間項は同型を除き一意である．

\begin{proposition}[connecting mapの最大rank]\label{prop:rank}
$P$のconnecting map$\delta:\Hom(e,S_+)\to\Ext^1(e,S_-)$は
$\rk\delta=\min(p,n)$を満たす．従って
\[
(h^0(P),h^1(P))=((p-n)_+,(n-p)_+),\qquad (h^0(Q),h^1(Q))=(p,n),
\]
ただし$h^j(M)=\dim\Ext^j(e,M)$，$x_+=\max(x,0)$とする．
このrankの証明には定理\ref{thm:derived}を必要としない．
\end{proposition}
\begin{proof}
符号別Extから$\Hom(e,P)=\ker\delta$，$\Ext^1(e,P)=\operatorname{coker}\delta$．
非零$f:e\to P$があれば，$P\to S_+$との合成は非零で全射となる．
像と$S_-$との交わりは零か$S_-$であり，前者なら拡大が分裂するので後者である．
従って$f$は全射．核$L$はrank-one torsion-freeで$c_1(L)=-(\alpha+\beta)$．
Corollary 10.12と$H(K+\alpha+\beta)=-H^2<0$により
$\Ext^2(e,L)=\Hom(L,\theta e)^*=0$．$e$のexceptionalityから$h^1(P)=0$となる．
$p>n$なら$\ker\delta\ne0$なので$\delta$は全射．$p\le n$で$h^0(P)>0$とすると
$\chi(e,P)=h^0(P)>0$に矛盾するので，この場合は単射である．
$Q$には逆向き短完全列の六項完全列を使う．
\end{proof}

\begin{proposition}[二根の全有限長対象]\label{prop:A2category}
$\cC=\langle S_+,S_-\rangle_{\mathrm{Serre}}$は
$\operatorname{mod}_{\mathrm{fd}}\Pi(A_2)$と同値であり，直既約対象は$S_+,S_-,P,Q$の四つである．
ここで$\Pi(A_2)$は二頂点・両向き二矢と，両方の長さ二の道を零とする関係を持つ代数である．
$p>n$では$(\cT\cap\cC,\cF\cap\cC)=(\add(S_+,P),\add(S_-))$，
$p<n$では$(\add(S_+),\add(S_-,P))$である．
\end{proposition}
\begin{proof}
非分裂性と単純性から，順序$(S_+,S_-,P,Q)$でHom次元は
\[
\begin{pmatrix}
1&0&0&1\\0&1&1&0\\1&0&1&1\\0&1&1&1
\end{pmatrix}.
\]
例えば$\Hom(P,Q)$は$P\twoheadrightarrow S_+\hookrightarrow Q$で生成され，
逆方向も同様で，両合成は零である．$\theta$はこれらの対象で自然に恒等なので
\[
\dim\Ext^1(U,V)=\dim\Hom(U,V)+\dim\Hom(V,U)+c_1(U)c_1(V).
\]
上の表を代入すると，$P,Q$から二単純へのExt$^1$はすべて零である．
長さによる帰納法から$P\oplus Q$は$\cC$のprojective generatorであり，
そのEnd代数は表示した四次元代数となる．
その表現$x:U_+\to U_-$，$y:U_-\to U_+$，$xy=yx=0$は，
$\operatorname{im}y\subset\ker x$，$\operatorname{im}x\subset\ker y$に適合する基底により，
一方の矢だけが非零の一次元対と孤立一次元空間へ分解される．従って直既約は四つである．
最後の主張は命題\ref{prop:rank}と符号別Extによる．
\end{proof}

$\cC$の内在的ExtはambientのExtと異なる．実際$\Pi(A_2)$の単純加群には
二周期的なprojective resolutionがあり，$\Ext^{2j}_{\cC}(S_+,S_+)=k$ $(j\ge0)$である．
従って$\Db(\cC)\to\Db(\AR)$は完全忠実でない．
また$p=n$なら$P$は自然な導来切断で消えるが，genuine avatarでは
$\alpha+\beta$が零指数根となるためこの同時共鳴自体が排除される．

\subsection{反標準側の二つのcharacter}
\begin{example}[実在する指数$(2,-1)$]
滑らかな平面三次$E$の変曲点を原点に取り，非torsion $q$について
\[
(p_1,\ldots,p_6)=(q,3q,2q,5q,7q,11q)
\]
とする．三点共線なら対応する正整数の和の$q$倍が零となり，六点conicも同様なので，
六点は一般位置である．$\alpha=e_1-e_2$，$\beta=e_2-e_3$の指数は$2,-1$．
Rainsのsimple-root effectivityと既約sheaf存在結果
\cite[\S4]{Filtered}\cite[\S11, p.~141]{Rains}の適用条件を
補題\ref{lem:Rains-simple-existence}で確認する．このchartでは上記の$S_+,S_-$が同じheartに実在する．
数値計算だけをsheaf存在の代わりには用いない．
\end{example}

\begin{lemma}[境界外でのExtの比較]\label{lem:localization}
有限次元$A=R/(g-1)$加群$U,V$を境界外の$\qgr R$の対象とみなすと，
全$j\ge0$で$\Ext_A^j(U,V)\simeq\Ext_{\qgr R}^j(U,V)$である．
\end{lemma}
\begin{proof}
$R[g^{-1}]$は強次数付き環で，その次数付き加群圏は$\operatorname{Mod}A$と同値である．
exactな中心局所化の右随伴はrestrictionで，これもexactかつ完全忠実である．
局所化加群は内部次数torsionに導来的に直交するため，右随伴は$\QGr R$にも降りる．
よって導来随伴も完全忠実となる．有限次元加群の非負次数格子は$k[g]$上有限自由で，
有限生成$R$加群を与える．locally noetherian圏でのcoherent Extとの一致により主張を得る．
\end{proof}

\begin{theorem}[三角表現と四次元商]\label{thm:affinequotient}
上の指数$(2,-1)$の例で$V=\Phi(S_+)$，$U=\Phi(S_-)[1]$，$W=\Phi(P)$とする．
これらはそれぞれ次元$2,1,1$の$A$加群であり，
\[
0\longrightarrow W\longrightarrow V\longrightarrow U\longrightarrow0
\]
は非分裂で，$W\not\simeq U$である．$A$の$V$への作用像は全三角行列代数となる．
逆向き非分裂拡大$0\to U\to V'\to W\to0$も存在し，
\[
A/\Ann_A(V\oplus V')\simeq\Pi(A_2).
\]
新しい二単純の根類は$\alpha+\beta,-\beta$であり，指数は$1,1$である．
\end{theorem}
\begin{proof}
最大rankから短完全列を得る．定理\ref{thm:derived}，補題\ref{lem:localization}と
$\Ext_X^1(P,S_-)=0$から$\Hom_A(W,U)=0$となり，二つの一次元加群は異なる．
非零包含$S_-\to P$は同値によって非零の$\Ext_A^1(U,W)$へ移るので，拡大は非分裂である．
異なるcharacterを持つ二次元非分裂表現の像には，二つの対角冪等元と非零の非対角cornerがあり，
像は全三角行列代数である．

同様に$\Ext_A^1(W,U)=\Ext_X^2(P,S_-)=\Hom_X(S_-,P)^*=k$から$V'$が存在する．
$V\oplus V'$への像は，共通の二つの対角成分と両向きの非対角成分からなる
四次元代数$C_0\simeq\Pi(A_2)$の部分代数$C$である．characterの分離から
$C$の半単純像は$k\times k$である．$J=\operatorname{rad}C_0$は$J^2=0$なので，
冪等元を持ち上げ，$1+J$による共役で標準の対角冪等元に合わせる．
両拡大の非分裂性により二つの一次元cornerがともに像へ入り，$C=C_0$となる．
最後に$[P]=[S_+]+[S_-]$，$[S_-[1]]=-[S_-]$を使う．
\end{proof}

反射$s_\beta(v)=v+(v\beta)\beta$は$(\alpha,\beta)$を$(\alpha+\beta,-\beta)$へ送る．
従ってこの局所模型では，正負の共鳴から反標準heartの正指数二単純への組替えがWeyl反射に対応する．
ここまでの主張は二根部分系に限る．Rainsの単純根存在結果を適用した$E_6$への拡張は，
後の系\ref{cor:E6spectrum}で別に与える．
商$A\to\Pi(A_2)$もhomological epimorphismではない．商の単純加群の高次Extは二周期的だが，
ambientのExtは次数2を超えて消えるからである．

\begin{remark}[同型型と生成元の係数]
この定理は実際の有限次元商の同型型を与える．しかし，楕円データからあらかじめ固定した
$A$の生成元に対するcharacterの値や拡大cocycleの係数を数値式で与えてはいない．
その明示計算と全$E_6$のアフィン作用行列は別の課題として残る．
単純加群の根類と次元の決定とは区別する．
\end{remark}

\section{反標準Euler標数によるheartの選択}\label{sec:rootchambers}
ADE型のspherical対象とWeyl反射，および単純tiltによる隣接heartの記述は既知の機構である．
Kleinian特異点についてはBridgeland \cite[\S2, Lemma 3.5]{BridgelandKleinian}を参照されたい．
本節の目的は，その機構を反標準比較へ適用し，\emph{指定構造対象のEuler標数}が
実際の標的heartをどのように選ぶかを明らかにすることである．
内部の$\Db(\cZ)$をambient導来部分圏と同一視する必要はなく，zigzagの形式性も仮定しない．
一般のlength heartの単純tiltが再びlengthとなる条件は\cite[Theorem 2.7]{BCPW}で扱われている．
ここでは2-spherical・単純辺という仮定から必要な有限長性を直接証明する．
幾何学的適用には第\ref{sec:A2}節と同じgenuine楕円gluing・非torsionの仮定を用いる．

\subsection{反標準側の有限次元heart}
$\cD_{\cZ}\subset\Db(\AR)$を，旧heartでの全cohomologyが$\cZ$に属する対象の充満三角部分圏とする．
定理\ref{thm:derived}の自然な$\Phi$を用い，$A=R/(g-1)$とする．
以下の$\operatorname{mod}_{\rm fd}A$は有限次元加群の\emph{abel圏}であり，
その内在的有界導来圏を$\cD_{\cZ}$と同定しているのではない．

\begin{lemma}\label{lem:targetfiniteheart}
$\Phi$は$\cD_{\cZ}$を，境界非交叉かつcohomologyが有限次元$A$加群に対応する
標的の充満三角部分圏へ同値に写す．従って$\cD_{\cZ}$には二つの有界heart
\[
 \cZ,\qquad \cK=\Phi^{-1}(\operatorname{mod}_{\rm fd}A)
\]
があり，$\cZ\subset D^{[0,1]}_{\cK}$である．
この部分圏でSerre関手は自然に$[2]$となり，
$m(X)=\chi(e,X)$は$0\ne X\in\cK$で正整数となる．
\end{lemma}
\begin{proof}
境界triangleの自然な射$\theta\to\mathrm{id}$は$\cD_{\cZ}$上で同型となる．
$\Phi$はこの射を中心元$g$に送るので，像のcohomologyは$g$可逆である．
各$M\in\cZ$について$\Ext^j(e,s^nM)$は$n$に関して周期的で有限次元だから，
像の二つのcohomologyは有限次元$A$加群となる．有界拡大により全$\cD_{\cZ}$で同様である．

逆に有限次元$A$加群は$k[g]$上有限自由な非負次数格子から$\qgr R$の対象を与える．
その逆像$X$には$\theta X\simeq X$が自然な境界射から従う．$\theta$は旧heartにexactなので
全$H^i(X)$も境界非交叉である．$\theta$の数値作用で$c_1$はrank倍の$K_X$だけ変わるから，
$K_X^2>0$の下でそのrankは零．Rainsの次元関数のrank-zero部分は次元高々1なので
$H^i(X)\in\cZ$となる．これは有限次元cohomologyを持つ標的の有界対象にも拡張される．
従って両t構造はこの充満部分圏へ制限される．
$\theta\simeq\mathrm{id}$からSerre関手は$[2]$．
中心局所化の導来随伴を$\pi R$と有限次元$U$に適用すると
$\RHom(\pi R,U)=\RHom_A(A,U)=U$であり，$m(\Phi^{-1}U)=\dim_k U>0$を得る．
\end{proof}

\subsection{全heartのずれは境界非交叉部分で検出される}
\begin{lemma}[中心元の全射性と有限次元局所化]\label{lem:g-surjection}
$U\in\qgr R$について，境界への通常の制限$i_R^*U$が零なら，
$g:U(-1)\to U$は同型であり，$U$は有限次元$A=R/(g-1)$加群から来る．
\end{lemma}
\begin{proof}
有限生成の次数付き代表$M$を取る．$i_R^*U=0$なら$M/gM$は有限次元なので，
$n\gg0$で$g:M_{n-1}\to M_n$は全射である．有限非負整数$\dim M_n$はやがて単調非増加となり，
従って定数になる．十分高い次数では$g$は同型となり，その核も有限次元である．
ゆえに$g$は$\qgr R$で同型で，$\varinjlim(M_n,g)$は有限次元．
この極限が$(M[g^{-1}])_0$であり，強次数付き環による局所化が所要の$A$加群を与える．
\end{proof}

\begin{theorem}[全heartの差を検出する負指数単純根]\label{thm:globaldefect}
定理\ref{thm:derived}，\ref{thm:globaltilt}のgenuine楕円gluingで$q$は非torsionとする．
全Rains heartのtorsionfree部分$\cF=\ker H^0\Phi$は$\cZ$に含まれる．
従って次の三条件は同値である：
\[
 \Phi\text{ はt-exact}
 \quad\Longleftrightarrow\quad(e,s)\text{ はAZ-ample}
 \quad\Longleftrightarrow\quad
 \chi(e,S)>0\text{ for every simple }S\in\cZ.
\]
特にこの範囲では，t-exactnessの失敗は実在する負指数の境界非交叉単純根によって必ず検出される．
\end{theorem}
\begin{proof}
定理\ref{thm:derived}の構成は$\Phi i_{q*}\simeq i_{R*}$を与えるので，左随伴の一意性から
$Li_R^*\Phi\simeq Li_q^*$となる．$M\in\cF$を$\Phi M=U[-1]$と書く．
旧heartでは境界の押出しがexactかつfaithfulであり\cite[\S5冒頭]{Filtered}，
境界triangleと$\theta$のexactnessから$Li_q^*$は振幅$[-1,0]$を持つ．
標的の$Li_R^*$も中心正則元の長さ1の解消から同じ振幅を持つため，
$Li_q^*M\simeq Li_R^*U[-1]$は次数0だけにcohomologyを持つ．
特に$H^0(Li_R^*U)=0$である．補題\ref{lem:g-surjection}から$Li_R^*U=0$となり，
$U$は有限次元$A$加群となる．補題\ref{lem:targetfiniteheart}より$M\in\cZ$である．

$\cZ$の単純対象はtorsion分解によって$\cT$か$\cF$の一方に入る．
そのEuler値はそれぞれ正，負であって零ではない．
$\cF\ne0$なら，その非零対象の$\cZ$での単純部分対象も$\cF$に入り，負指数を持つ．
逆に負指数単純があれば$\cF\ne0$である．t-exactnessと$\cF=0$の同値を用いる．
$t$-exactならpointedな同値が標準ample pairを運ぶのでAZ-ampleである．
逆に非零$M\in\cF\subset\cZ$では$\Hom(e,s^nM)=0$が全$n$で成り立ち，AZ生成性に反する．
\end{proof}
\begin{remark}
ここで得た逆向きの含意は，指定genuine gluing，非torsion，既知の導来比較と振幅を使う．
数値的共鳴の有無だけを判定しているのではなく，\emph{実在する向き付き単純根}を全て調べる条件である．
対象の構成やこれらの仮定を外した一般の非可換曲面への主張ではない．
\end{remark}

\subsection{全ての部分対象・商によるtorsion pairの認識}
\begin{proposition}[Euler標数によるtorsion pairの完全な記述]\label{prop:numericaltorsion}
$\cZ$上の反標準torsion pairは
\[
\begin{split}
 \cT\cap\cZ&=\{M:\ m(Q)>0\text{ for every nonzero quotient }Q\text{ of }M\},\\
 \cF\cap\cZ&=\{M:\ m(N)<0\text{ for every nonzero subobject }N\text{ of }M\}
\end{split}
\]
である．零対象に対する条件は空条件とする．
\end{proposition}
\begin{proof}
非零$T\in\cT\cap\cZ$では$\Phi T$が非零有限次元加群なので$m(T)>0$．
非零$F\in\cF\cap\cZ$では$\Phi F=U[-1]$で$m(F)=-\dim U<0$．
$\cT$は商に，$\cF$は部分対象に閉じるため，それぞれ表示した条件を満たす．
逆に$0\to T\to M\to F\to0$をtorsion分解とすれば，商の条件は$F\ne0$を排除し，
部分対象の条件は$T\ne0$を排除する．
\end{proof}
$m(M)$の符号だけでは不十分である．混合$A_2$の$P,Q$が同じEuler標数を持つ例は，
上の全商・全部分対象の条件を全対象の符号へ置き換えられないことを示す．

\subsection{有限ADE型における単純tiltの有限手順}
まず一般のExt有限な三角圏$\cD_0$でSerre関手が$[2]$である場合を考える．
その有界有限長heart $\cH$の単純対象を$S_1,\ldots,S_r$とし，
全て2-sphericalで，異なる単純の$\Ext^1$が有限ADE図式の隣接数に等しいと仮定する．
$\alpha_i=[S_i]$，$(\alpha_i,\alpha_j)=-\chi(S_i,S_j)$を負定値根形式とする．

\begin{lemma}[一つの単純対象でのtilt]\label{lem:simpletilt}
頂点$j$について
$\cF_j=\add(S_j)$，$\cT_j=\{M:\Hom(M,S_j)=0\}$とすると，$(\cT_j,\cF_j)$はtorsion pairである．
$\cH'=\langle S_j[1],\cT_j\rangle$は有限長heartで，その単純対象は
\[
 S_j[1],\qquad
 P_i\ (i\text{ が }j\text{ に隣接}),\qquad
 S_i\ (i\ne j\text{ が }j\text{ に非隣接})
\]
である．ここで$P_i$は非分裂拡大$0\to S_j\to P_i\to S_i\to0$．
新しい単純根は$s_j(\alpha_i)$であり，$s_j(v)=v+(v,\alpha_j)\alpha_j$とする．
新しい単純も2-sphericalで，同じADE図式を持つ．
\end{lemma}
\begin{proof}
$\Ext^1(S_j,S_j)=0$なので，最大の$S_j$-半単純商への評価射
$M\twoheadrightarrow S_j^{\oplus\dim\Hom(M,S_j)}$の核は$\cT_j$に入る．
これがtorsion分解を与える．新heartの候補は表示した通りで，隣接$i$については
新heartで$0\to P_i\to S_i\to S_j[1]\to0$となる．

有限長性を確かめる．$M\in\cT_j$の旧heartでの単純商を$S_i$とすると$i\ne j$である．
核$K$のtorsion部分を$T_j(K)$とし，$E=M/T_j(K)$と置く．
$0\to S_j^{\oplus b}\to E\to S_i\to0$であり，$\Hom(E,S_j)=0$から
$k^b\hookrightarrow\Ext^1(S_i,S_j)$を得る．従って$b=0$，または$i$が隣接して$b=1$の非分裂拡大だけである．
旧lengthによる帰納法と上の新heartの短完全列により，$M$は表示した候補で有限filtrationを持つ．
全新heartもそのようなfiltrationを持つ．

候補間のHomは異なる二対象で零，同じ対象で$k$である．
例えば$\Hom(P_i,S_j)=0$は非零拡大類によるconnecting mapから従い，
$\Hom(P_i,S_j[1])=0$はSerre双対性とEuler形式から得る．
異なる$P_i,P_k$間では，$P_k$の部分対象$S_j$と商$S_k$へのHomを用いれば零となる．
従って候補は本当に全単純対象である．その根類は表示した反射で，反射はEuler形式を保つ．
Serre双対性とHomの計算から各自己Ext$^1$は零，cross Ext$^1$は元の隣接数となる．
\end{proof}

\begin{theorem}[反標準heartのEuler chamber認識]\label{thm:Eulerchamber}
$\cH\subset\cD_0$を上の有限ADE型heartとし，もう一つの有界heart$\cK$について
$\cH\subset D^{[0,1]}_{\cK}$とする．加法的関数$m:K_0(\cD_0)\to\mathbb Z$が
全ての非零$U\in\cK$で$m(U)>0$を満たすと仮定する．
このとき次の有限手順は$\cK$を回復する．
\begin{enumerate}
\item 現在の単純のうち$m(S_j)<0$のものを一つ選ぶ．
\item 補題\ref{lem:simpletilt}の$\langle S_j[1],\cT_j\rangle$へ移る．
\item 負の単純がなくなるまで繰り返す．
\end{enumerate}
手順は初期正根系$\mathcal R_+$について
\[
 N(m)=\#\{\alpha\in\mathcal R_+:m(\alpha)<0\}
\]
回で終了する．終端heartは選択順序に依存せず$\cK$であり，その単純根系は
正根系$\{\alpha:m(\alpha)>0\}$の単純根系に等しい．
特に$m$は全ての根で非零であることが従う．
\end{theorem}
\begin{proof}
現在のheartの単純$S$には，振幅1の比較からtorsion分解がある．単純性により
$S\in\cK$または$S\in\cK[-1]$のいずれかであり，その$m$はそれぞれ正または負である．
従って$m(S)=0$は起こらない．負の$S_j$を選ぶと$S_j[1]\in\cK$であるから，
新heartも$D^{[0,1]}_{\cK}$に入る．補題により有限長性とADE条件も保たれる．

現在の正根系を単純反射$s_j$で反射すると，正根の集合は$\alpha_j$だけを
$-\alpha_j$に入れ替えたものとなる．その$m$は負から正へ変わるので，負の正根の数は正確に1減る．
この数は非負整数であるため手順は停止する．全単純の$m$が正となったとき，
それらは$\cK$に入り，新heartの有限長性から新heart全体が$\cK$に含まれる．
同じ三角圏の有界heartの包含は一致を意味するので，終端は$\cK$である．
終端の各正根は単純根の非負整数結合で$m>0$となり，全根での非零性も従う．
最後に，正根系を指定するregular線形汎関数のchamberは一意なので単純根系も一意である．
対象としての終端heartの独立性は，根類の一致だけでなく，既に示した$\cK$との一致による．
\end{proof}
\begin{remark}
これは一般のWeyl反射や単純tiltを新しく導入する定理ではない．
本稿では特に$\cH=\cZ$，$m=\chi(e,-)$，$\cK=\Phi^{-1}(\operatorname{mod}_{\rm fd}A)$に適用し，
\emph{反標準比較が指定する}heartと，そこでの有限次元単純加群の次元を決める．
得られた根類$\beta_i$に対応する単純加群の次元は$m(\beta_i)$である．
異なる経路が同じheartに到達することと，それらの合成自己同値が自然に同型であることは別であり，後者は主張しない．
またこれは固定圏におけるEuler汎関数によるheartの選択であり，
楕円パラメータの族での全ての壁越えや安定性空間全体の記述ではない．
\end{remark}

\subsection{局所的な反射手順の全heartへの延長}
\begin{lemma}[境界非交叉単純の有限ADE性]\label{lem:finiteADEautomatic}
定理\ref{thm:globaldefect}の設定では，$\cZ$の単純対象は有限個で，
その根類は線形独立な有限ADE単純根系をなす．空集合も許す．
\end{lemma}
\begin{proof}
単純$S$に対して$\theta S=S$であり，$\End(S)=\Ext^2(S,S)=k$．
$c_1(S)=\alpha\ne0$，$H\alpha=0$と負定値性，Euler形式の偶奇から
$0<-\alpha^2=2-\dim\Ext^1(S,S)\le2$，従って$\alpha^2=-2$である．
同根の異なる単純はHomとExt$^2$が零なのにEuler値が2となって矛盾する．
負定値格子には長さ2の根が有限個しかないので，単純も有限個である．
異なる二根の交点はExt$^1$の次元に等しく非負．
反対根は有効錐のpointednessで排除されるので，交点は0か1である．

根類が有理係数の線形関係を持つとする．正負の係数を分けて$x=y$と書けば，
支持が異なるため$x\cdot y\ge0$，従って$x^2\ge0$となる．
負定値性から$x=0$であるが，非自明な正係数の有効類の和が零であることは
有効錐のpointednessに反する．片側しかない関係も同様に排除される．
従って根類は独立であり，対角2，非対角0または$-1$の正定値Cartan行列を持つ．
有限型の分類からADE型である．
\end{proof}

\begin{theorem}[全Rains heartの有限単純tilt分解]\label{thm:globalADEtilts}
定理\ref{thm:globaldefect}の範囲で，定理\ref{thm:Eulerchamber}の負指数単純tiltは
全Rains heart上に延長できる．初期$\cZ$の正根系を$\mathcal R_+$とすると，
$\#\{\alpha\in\mathcal R_+:m(\alpha)<0\}$回で全反標準heartへ到達する．
単純対象がない場合は初めから両heartは一致する．
\end{theorem}
\begin{proof}
$\cZ$がSerre部分圏なので，その単純対象は全heartでも単純である．
$\Hom$の有限性と$\Ext^1(S_j,S_j)=0$から，全対象に対する最大$S_j$-半単純商を用いて
$(\ker\Hom(-,S_j),\add S_j)$というtorsion pairを全heart上で構成できる．
負指数なら$S_j[1]$は標的heartに属すから，tilt後も振幅$[0,1]$を保つ．

現在の境界非交叉Serre部分圏を$\cZ'$とする．新heartの短完全列の旧cohomologyを取ると，
次数$-1$の項は全て$\add S_j\subset\cZ'$に入る．従って旧$\cZ'$のSerre性から，
$\cD_{\cZ}$と新heartとの交わりもSerre部分圏である．その制限は補題\ref{lem:simpletilt}のheartである．
新しいtorsionfree部分の対象$X$を，旧cohomologyのtriangle
$H^{-1}(X)[1]\to X\to H^0(X)$で記述する．
$H^{-1}(X)[1]$は標的heartに属し，$H^0(X)$の標的cohomologyは負次数にないため，
$H^{-1}(X)[1]\hookrightarrow H^0_{\rm target}(X)=0$となる．
よって$X$は旧heartにも属し，旧torsionfree部分の対象である．
従って各段階の残るtorsionfree部分は$\cD_{\cZ}$の中にある．

有限ADE性は補題\ref{lem:finiteADEautomatic}で得られる．局所的な手順が標的の有限次元heartに
到達すると，残るtorsionfree部分は零になる．命題\ref{prop:hrs}により全heartも一致する．
回数は定理\ref{thm:Eulerchamber}から従う．
\end{proof}
この有限分解は一回のHRS tiltという全体記述を，有限個のspherical単純による操作へ細分する．
中間heartを別に指定したblowdown表示のRains heartと同一視する必要はない．

\subsection{最大共鳴の単純対象}

\begin{proposition}[全単純根が実在する場合の網羅性]\label{prop:maximalrootexhaustion}
$\cZ$の根格子が有限ADE型で，選んだ全単純根$\alpha_i$に対する
既約境界非交叉sheafが実在すると仮定する．このときそれらが$\cZ$の全単純対象であり，
同根の単純は同型を除き一意である．
\end{proposition}
\begin{proof}
任意の単純$M\in\cZ$は根類を持つ．実際，$D=c_1(M)\ne0$は$H^\perp$にあり，
負定値性，Riemann--Rochの偶奇，Serre双対性から
$0<-D^2=2-\dim\Ext^1(M,M)\le2$となるので$D^2=-2$である．
負根が有効なら，その反数が実在単純根の非負和として有効になり，
有効錐のpointedness \cite[Proposition 7.19]{Rains}に反する．
非単純正根$\gamma=\sum c_i\alpha_i$では
$-2=\gamma^2=\sum c_i(\gamma,\alpha_i)$だから，ある$i$で$(\gamma,\alpha_i)<0$．
しかし異なる単純間のExt$^1$の次元はこの交点に等しいため不可能である．
同根の異なる単純も，HomとExt$^2$が零なのにEuler値が2となって矛盾する．
\end{proof}

\subsection{単純根の実在入力と整数倍点族の完全な次元表}
\begin{lemma}[Rainsの単純根存在結果の適用]\label{lem:Rains-simple-existence}
滑らかな楕円境界を持つ指定Rains blowdown表示で，$\alpha$をその向き付き単純根とする．
$q$が非torsionで，ある整数$m$について$\rho(\alpha)q^m=\OO_E$なら，
境界非交叉の既約純一次元sheaf $S_\alpha$で
$c_1(S_\alpha)=\alpha$，$\chi(e,S_\alpha)=m$となるものが存在する．
本稿のnef・正次数の範囲では，これは単純であり同型を除き一意である．
\end{lemma}
\begin{proof}
存在はRainsの既知結果を適用する．\cite[\S4, pp.~34--35]{Filtered}は，
指定blowdownの単純根が$q$の冪へのrestrictionを持つとき有効となることと，
対応する$(-2)$類のpure sheafに対するdeterminant条件を記述している．
\cite[\S11, p.~141, Proposition 11.23の直前]{Rains}では，
\emph{有効単純根そのもの}の類に対し，そのdeterminant条件を満たすEuler標数で
既約sheafが存在し，境界と交わらないことが明記されている．
Proposition 11.23のnef・fundamental chamberの主張を単純根に誤用しているのではない．
本稿の規約でrank-zero類のrestriction determinantは$\rho(\alpha)q^m$である．
従って入力の仮定を満たす．単純性は補題\ref{lem:offboundary}，一意性は
同根の異なる単純間でEuler値2となる矛盾による．
\end{proof}
\begin{remark}
これは任意の数値的共鳴根に無条件にsheafを割り当てる主張ではない．
\emph{固定したRains heartのblowdown単純根}であることを追加して既知の存在結果を使っている．
以下の整数倍点族では，各単純根についてこの条件を個別に確認する．
\end{remark}

\begin{theorem}[整数倍点族：整列した差による全単純の次元]\label{thm:sorted-spectrum}
$E\subset\mathbb P^2$のflexを原点$o$とし，$q\in E$を非torsionとする．
$3\le r\le8$，相異なる正整数$N_1,\ldots,N_r$に対して$p_i=N_iq$と置く．
$S=\operatorname{Bl}_{p_1,\ldots,p_r}\mathbb P^2$と，この順序のblowdownによるRains heartを取る．
$q$のPicard類は$\OO_E(q-o)$であり，共鳴の符号は\eqref{eq:resonance}に従う．
このとき$S$はgenuine del Pezzoであり，境界非交叉heartの全単純根は
\[
 \alpha_0=h-e_1-e_2-e_3,\qquad \alpha_i=e_i-e_{i+1}\quad(1\le i<r)
\]
である．対応するsheafは全て実在し，その指数は
$(N_1+N_2+N_3,N_2-N_1,\ldots,N_r-N_{r-1})$である．

$N_{\pi(1)}<\cdots<N_{\pi(r)}$となる並べ替えを$\pi$とする．
$A=R/(g-1)$の有限次元単純加群は同型を除き正確に$r$個で，その根類と次元は
\begin{align*}
 \beta_0&=h-e_{\pi(1)}-e_{\pi(2)}-e_{\pi(3)},
 &\dim V_0&=N_{\pi(1)}+N_{\pi(2)}+N_{\pi(3)},\\
 \beta_i&=e_{\pi(i)}-e_{\pi(i+1)},
 &\dim V_i&=N_{\pi(i+1)}-N_{\pi(i)}\quad(1\le i<r).
\end{align*}
全Rains heartから反標準heartへの負単純tiltの回数は
\[
 \operatorname{inv}(N)=\#\{(i,j):i<j,\ N_i>N_j\}
\]
である．特に両heartが一致することと$N_1<\cdots<N_r$は同値である．
\end{theorem}
\begin{proof}
$q$が非torsionなので点は相異なる．blow-up後の滑らかな反標準曲線の類$H$はnefで，
$H^2=9-r>0$である．$H^\perp$の根を$ah-\sum b_ie_i$と書くと
$\sum b_i=3a$，$\sum b_i^2=a^2+2$．$a\ge0$の根は差根，三点の和を引く直線類，
六点の和を引くconic類，および$r=8$で一つの係数が2・残りが1の三次曲線類だけである．
実際Cauchy--Schwarzで$(9-r)a^2\le2r$，かつ
$\sum b_i(b_i-1)=a^2-3a+2$を使えば$a\le3$とこの列挙を得る．
可能性$a=4,r=8$では右辺6を一つの3または三つの2で分配しても$\sum b_i=12$に届かず排除される．

restrictionは$\rho(ah-\sum b_ie_i)=q^{-\sum b_iN_i}$である．
差根の指数は相異なる二整数の差，$a>0$の根の指数は正整数の正係数和なので零にならない．
もし$H$がampleでなければ，nef・bigとNakai--Moishezonから$HC=0$の既約曲線がある．
Hodge指数定理とadjunctionより$C^2=-2$で，境界と交わらずrestrictionが自明となり矛盾する．
よって$S$はgenuineである．全ての表示した単純根は$q$の整数冪にrestrictionするため，
補題\ref{lem:Rains-simple-existence}を\emph{同じblowdown表示}で各根に適用する．
命題\ref{prop:maximalrootexhaustion}が全単純の網羅性を与える．

正根のうち負指数を持つものは差根$e_i-e_j$ $(i<j,N_i>N_j)$だけである．
負の隣接差を選ぶ単純反射は隣接二整数の交換であり，$\alpha_0$の指数は常に
現在の先頭三整数の和だから正である．従って隣接逆転を解消する整列が
定理\ref{thm:Eulerchamber}の手順となり，終端の単純根と正指数は表示したものになる．
補題\ref{lem:targetfiniteheart}がこれらを全有限次元単純加群に対応させる．
定理\ref{thm:globalADEtilts}から同じ回数で全heartも一致する．
初期が整列済みのときだけ負の単純がないので，最後の同値も従う．
\end{proof}

\begin{corollary}[同じpointed導来圏内の異なるRains hearts]\label{cor:permuted-hearts}
上の無順序点集合を固定すると，その$r!$通りの順序から得るRains heartsは，
同じ$\cD_q$と同じ$e$の中で相異なる．反標準環$R(e,s)$は共通であり，
この$r!$個のうち反標準heartと一致するのは整数を昇順に並べたものだけである．
\end{corollary}
\begin{proof}
可換blow-up $S$は点の順序によらず，gluing\eqref{eq:glue}は$(S,E,q)$と$e$から定まる．
従って導来圏と構造対象，Serre関手，反標準環は共通に固定できる．
一方，異なる二順序には前後関係が逆になる点の対がある．一方のheartでは$e_i-e_j$が
隣接単純根の正和として有効で，他方ではその反数が有効である．
もし両heartが同じなら有効錐のpointednessに反する．最後は定理\ref{thm:sorted-spectrum}による．
\end{proof}

\begin{corollary}[六点の具体例]\label{cor:E6spectrum}
$(N_1,\ldots,N_6)=(1,3,2,5,7,11)$では，全六単純根sheafが実在し，
一回の$\alpha_2$-tiltで全反標準heartへ移る．全有限次元単純の根類と次元は
\[
\begin{array}{c|rrrrrr}
\text{根類}&\alpha_0&\alpha_1+\alpha_2&-\alpha_2&\alpha_3+\alpha_2&\alpha_4&\alpha_5\\\hline
\text{次元}&6&1&1&2&2&4
\end{array}
\]
であり，一次元characterは正確に二つである．同じ無順序六点集合は同じpointed導来圏内に
$6!=720$個の相異なるRains heartsを与える．
\end{corollary}
\begin{proof}
整列は$(1,2,3,5,7,11)$で，逆転は一つだけである．前二結果を適用する．
\end{proof}
\begin{remark}[共鳴とheartの不一致は同義ではない]
この族では全根が数値的に共鳴しているが，整列した順序のRains heartは反標準heartと一致し，
$(e,s)$はAZ-ampleである．非整列順序では少なくとも一つの負指数単純があり，AZ-ampleでない．
従って最大の数値的共鳴とheartの一致は両立する．この主張はheartの\emph{同じ導来圏内での位置}についてであり，
構造対象を忘れた抽象的abel同値の全分類ではない．
単純加群の根類・次元は決定したが，元の楕円生成元の作用行列や全アフィン環の表示は未計算である．
既知のRains存在結果とWeyl反射を適用した帰結として扱い，有限次元表現論一般の新発見とは主張しない．
\end{remark}


\section{四者を同じ図式に置くための条件}\label{sec:alignment}
比較は二圏ずつの同値の存在だけでは完了しない．同じ$e$と$s$を使って反標準環を作り，
比較の合成がその構造を保つことを確認する必要がある．
\begin{proposition}[pointedな比較の接続]\label{prop:alignment}
次のデータが与えられたとする．
同じblow-up表示によるpointedな$J:\Db(\AV)\simeq\Db(\AR)$，
定理\ref{thm:derived}の$\Phi$，およびpointedな
$E:\Db(\AAS)\simeq\Db(\AR)$がある．さらにAS側で系\ref{cor:ASac}の条件と
命題\ref{prop:unique}の高次消滅が成立するとする．
このとき$\Phi E$は標準AS heartから反標準heartへのt-exact同値であり，
$\Phi J$はVdB heartに対して$J^{-1}(\cT),J^{-1}(\cF)$によるHRS比較となる．
従って$\cF=0$なら四者はこれらの指定関手の下で一致する．
\end{proposition}
\begin{proof}
corner同値からAS側の指定反標準orbitは豊富である．pointed性により環とorbitを同定でき，
両heartは命題\ref{prop:unique}の認識式を満たすので$\Phi E$はt-exact．
$J$のt-exactnessはtorsion pairを$\AV$へ運び，定理\ref{thm:globaltilt}の式を保つ．
$\cF=0$のときの結論は命題\ref{prop:hrs}から従う．
\end{proof}

実際のrestriction条件\eqref{eq:two-nontrivial}の下では，定理\ref{thm:AOU-pointed}により
section model--AS比較のpointed性が得られる．
VdBの最後のraw blow-upと同じRains表示については，系\ref{cor:fourfunctors}で
指定関手を実際に接続した．異なるavatarとの比較にはなおデータの同定を要する．
確認済みの範囲と残る同定は，次の表のとおりである．
\begin{center}
\small
\begin{tabular}{p{.18\textwidth}p{.36\textwidth}p{.35\textwidth}}
\toprule
比較 & 結論と適用範囲 & 自動的には従わないこと\\
\midrule
VdB--Rains & 同じblow-up・gluing・blow-down表示でabel同値 & 独立に選んだavatarのデータ対応\\
Rains--反標準 & genuine smooth楕円gluingで導来同値，一回のHRS tilt & 任意nc曲面・特異境界への拡張\\
VdB--section model & VdBの一般位置条件下で振幅1の導来同値 & raw blow-upとのheart一致\\
AS--反標準 & 規格化共通領域でnoetherian性と標準heartの一致を導出 & 任意potential・全AS locusへの拡張\\
section model--AS & 実際のrestriction条件の下でpointed導来同値，反標準環も一致 & 全域での標準heartの一致\\
四者の接続 & 同じ最終VdB表示では自然関手を同定，$\Delta\ne0$で単一tilt & 独立avatarの同定，相対降下\\
\bottomrule
\end{tabular}
\end{center}

\subsection{可換基準点と非可換共鳴}
滑らかな可換三次曲面$S$では，実際のblow-up heartとRainsの可換heartは$\coh S$である．
$-K_S$がampleなので反標準環の$\qgr$も$\coh S$を回復する．
標準的なline bundle helixの各周期部分列$\OO(D_i-nH)$もSerreの生成・消滅を満たし，
ASの$\mathbb Z$-代数から同じheartを再構成する．この場合には四者の一致を実際に確認できる．
しかし，一つの可換点での一致から，非可換パラメータ空間での一致locusの開性を推論しない．
命題\ref{prop:density}は，楕円chartでは負指数共鳴が可算個の軌跡の合併として稠密となることを示す．
「稠密な共鳴」と「一つの閉じた判別因子」，あるいは「導来特異化」は異なる概念である．

\subsection{族に移すときの追加課題}
本稿の主要な比較は幾何学的点での結果である．基礎$T$上の族へ移すには，
少なくともgluing bimoduleとkernelの相対perfect性，orbitのHomの基底変更，
有限性と消滅次数の一様性，t構造の基底変更適合性を確認する必要がある．
構造対象の同型と反Serre作用への適合を降下データとして保つことも必要である．
構造対象の同定については命題\ref{prop:AOU-relative-object}が相対Homのperfectnessと
基底変更の下で評価triangleから自然な同型を与える．その前提となる相対圏の構成と
全heartの降下を，この対象単独の同定から推論しない．

有限cornerの範囲だけなら次の初等的な相対判定が使える．
\begin{proposition}\label{prop:relativecorner}
有限型基礎$T$上に次数1生成のgraded代数族$C$があり，各$C_N$が有限表示で
fiberの形成と両立するとする．このときfiberが有限次元となるlocusは
$\bigcup_{N\ge1}\{t:(C_N)_t=0\}$であり開である．
$T$がnoetherianなら，この開locus上である共通の$N$を選べる．
\end{proposition}
\begin{proof}
各集合はcoherent module $C_N$のsupportの補集合であり開である．
補題\ref{lem:cutoff}によってこれらは$N$とともに増加し，合併が求めるlocusとなる．
noetherian空間の開集合は準コンパクトなので有限部分被覆を選び，最大の$N$を取る．
\end{proof}
三次曲面型では第\ref{sec:explicitcorner}節の十一条件がこの開集合の明示式を与える．
この命題は有限cornerのfiber条件を述べるだけであり，heartの相対同値・flatness・stack降下を
証明するものではない．本稿のt構造比較と，内在的非可換del Pezzoモジュライの
全成分分類は分離して扱う．

\appendix
\section{三次曲面の数値データと再現手順}\label{app:numerics}
可換三次曲面の標準格子を
$\Pic(S)=\mathbb Zl\oplus\bigoplus_{i=1}^6\mathbb Ze_i$，
$l^2=1$，$e_i^2=-1$，$H=3l-\sum e_i$とする．
AOU \cite[(4.1)]{AOU}の九線束のChern類は
\[
0,\quad -2l+\sum e_i,\quad -l+e_1+e_2+e_3,
\quad e_1,e_2,e_3,\quad l-e_4,l-e_5,l-e_6.
\]
各三つが一blockをなし，$H$次数は$0,1,2$である．
$\chi(\OO(D),\OO(E))=1+\tfrac12(E-D)(E-D+H)$からEuler行列は
\[
C=\begin{pmatrix}I_3&J_3&2J_3\\0&I_3&J_3\\0&0&I_3\end{pmatrix},
\qquad \det C=1,
\]
ただし$J_3$は全成分1の行列．成分総和は45である．
これは可換点の既知のfull strong性と合わせてEnd代数の次元を確認するが，
Euler行列だけからstrong性を証明してはいない．
$K_0$でのSerre作用は$C^{-1}C^{\mathsf t}$，反Serre作用はその逆であり，
後者は$H$によるtensorと一致する．

根は$e_i-e_j$ $(i\ne j)$，$\pm(l-e_i-e_j-e_k)$ $(i<j<k)$，
$\pm(2l-\sum e_i)$の72個である．各根について$\alpha^2=-2$，$H\alpha=0$．
反標準Euler値は$n=0,\ldots,8$で
\[
1,4,10,19,31,46,64,85,109
\]
となる．これらを添付コードで有理数上の厳密計算として確認する．

混合$A_2$では，四直既約のHom表をquiver射の線形方程式から求め，
Euler形式とSerre双対からambient Ext$^1$を検算する．
指数$p,n=1,\ldots,8$，四直既約の重複度$0,\ldots,4$について
40,000通りのcohomology公式を検査する．$p=n$の検算は仮想的線形代数の確認であって，
genuine avatar上にそのsheafが存在することの証拠ではない．

構造対象のmutationについては，AOUの操作の合成順序を実装し，
標識対象の位置が$9,9,9,6,6,3,1$と移ることと，最後の二箇所だけに
Hom消滅の入力が必要であることを検査する．
実際の$\mathcal U_2,\mathcal U_1$を使った\eqref{eq:two-degree-zero}の次数零性も検査する．
印刷された旧指数の検算は境界規約を保証するものではなく，現在は
第\ref{subsec:boundaryaudit}節の不整合を再現する検査と区別する．
境界検算ではさらに，逆の右mutationの$[-1]$と，末尾三対象への$[1]$を明示して
各段階のrankとdeterminantを計算し，実際の正規化したdecupleの指数と一致することを検査する．
これは非自明な楕円線束のcohomology消滅やmutation対象の同型を計算で証明するものではなく，
証明に使う順序と数値の独立な確認である．

Euler chamberの検算は有限根系$A_2+A_1,A_4,D_5,E_6,E_7,E_8$について行う．
各型で32個のregular汎関数を作り，負の単純を選ぶ三通りの規則を比較する．
停止回数が初期負正根数に一致し，終端根集合が正値根の単純根集合に一致することを整数上で検査する．
これは有限個の根格子計算であって，単純sheafの実在やheartの同定を数値実験に置き換えるものではない．

整数倍点族では$3\le r\le8$の各型で40組の相異なる正整数を選び，
Picard格子での根の明示列挙とWeyl反射による列挙を照合する．
具体的六点の全720順列について，停止回数が逆転数に等しく，全単純の次元が同じ整列公式となることを検査する．
これらは格子の検算であり，Rainsの既知のsheaf存在結果を計算で証明するものではない．

本稿は主TeXファイルと九つのsectionソースからなり，\texttt{latexmk -lualatex}で再生成できる．
日本語LuaLaTeX，LuaTeX-ja，\texttt{tikz-cd}，Python~3を用いる．
リポジトリrootでの主要コマンドは
\begin{verbatim}
python3 computations/verify_cubic_heart_numerics.py
python3 computations/verify_four_hearts.py
python3 computations/verify_four_hearts_corner.py
python3 computations/verify_four_hearts_chambers.py
python3 computations/verify_four_hearts_sorted.py
python3 computations/verify_four_hearts_boundary.py
python3 computations/verify_four_hearts_elliptic.py
latexmk -lualatex papers/nc_del_pezzo_hearts.tex
\end{verbatim}
である．数値検算は，sheafの存在，derived type，導来完全忠実性，
または新規性を計算によって証明するものではない．




\begin{thebibliography}{99}
\bibitem{VdB}
M.~Van den Bergh, \emph{Blowing up of non-commutative smooth surfaces},
Mem. Amer. Math. Soc. \textbf{154} (2001), no.~734;
\href{https://arxiv.org/abs/math/9809116}{arXiv:math/9809116v1}.
使用箇所：(5.4)，Theorem 6.3.1，\S10冒頭，(9.5)，Lemma 9.1.5，Proposition 9.1.8，Theorem 9.1.9，Proposition 9.2.3，
Lemma 11.1.1，Definition 11.1.2，Theorems 11.1.3, 11.2.1，Lemma 11.2.4，Remark 12.0.5．
\bibitem{Rains}
E.~M.~Rains, \emph{The birational geometry of noncommutative surfaces},
\href{https://arxiv.org/abs/1907.11301}{arXiv:1907.11301v1} (2019).
使用箇所：Lemma 4.9，Theorem 4.10，Propositions 4.13, 4.17，Theorem 4.30，
\S5，\S7，Theorem 9.10，Corollaries 10.12, 10.16，Proposition 10.13，\S11（特にp.~141，Proposition 11.23の直前の有効単純根の存在についての記述）．
\bibitem{Filtered}
E.~M.~Rains, \emph{Filtered deformations of elliptic algebras},
\href{https://arxiv.org/abs/2107.13540}{arXiv:2107.13540v2} (2022).
使用箇所：Corollary 3.8とtraceの証明，\S4（(4.4)--(4.7)の一般の有理曲面のgluingと境界，
特にpp.~34--35の単純根の有効性・存在），Proposition 5.1．
\bibitem{RainsBook}
E.~M.~Rains, \emph{The (noncommutative) geometry of difference equations},
\href{https://arxiv.org/abs/2504.16187}{arXiv:2504.16187} (2025).
v1のLemma 11.10，Proposition 11.19，Theorem 11.31，Theorem 16.10，
Corollary 17.17，およびProposition 18.27直前の有効単純根の存在記述を参照．
旧稿の対応番号は序論に示す．
\bibitem{AOU}
T.~Abdelgadir, S.~Okawa, K.~Ueda,
\emph{Compact moduli of marked noncommutative cubic surfaces},
\href{https://arxiv.org/abs/2404.00175}{arXiv:2404.00175v1} (2024).
使用箇所：\S3--4，(4.1)，\S9，Lemma 9.2，Theorem 9.4，(9.2)--(9.3)，(9.25)，(9.28)--(9.57)，Definition 7.1，Proposition 9.3．
これらの境界表示の同時使用に関する検査は第\ref{subsec:boundaryaudit}節を参照．
\bibitem{OU}
S.~Okawa, K.~Ueda, \emph{Marked Artin--Schelter surfaces of del Pezzo types},
\href{https://arxiv.org/abs/2508.04986}{arXiv:2508.04986v1} (2025).
使用箇所：Definitions 1.1, 3.1，\S6．
\bibitem{OUHelices}
S.~Okawa, K.~Ueda, \emph{AS-regular algebras from acyclic spherical helices},
\href{https://arxiv.org/abs/2007.07620v1}{arXiv:2007.07620v1} (2020).
使用箇所：Definitions 3, 4，Theorem 5，(18)--(22)．
\bibitem{AZ}
M.~Artin, J.~J.~Zhang, \emph{Noncommutative projective schemes},
Adv. Math. \textbf{109} (1994), 228--287.
\href{https://doi.org/10.1006/aima.1994.1087}{doi:10.1006/aima.1994.1087}.
\bibitem{CHZ}
X.-W.~Chen, Z.~Han, Y.~Zhou, \emph{Derived equivalences via HRS-tilting},
Adv. Math. \textbf{354} (2019), 106749;
\href{https://arxiv.org/abs/1804.05629}{arXiv:1804.05629v2}.
使用箇所：Theorems A, B，\S2--3．
\bibitem{RSS}
D.~Rogalski, S.~J.~Sierra, J.~T.~Stafford,
\emph{Noncommutative Blowups of Elliptic Algebras},
\href{https://arxiv.org/abs/1308.2216}{arXiv:1308.2216v2} (2015).
使用箇所：Proposition 2.4（正次数，無限位数の境界自己同型）．
\bibitem{YZ}
A.~Yekutieli, J.~J.~Zhang, \emph{Rings with Auslander Dualizing Complexes},
J. Algebra \textbf{213} (1999), 1--51;
\href{https://arxiv.org/abs/math/9804005}{arXiv:math/9804005v3}.
使用箇所：Theorems 4.5, 5.1(1), 5.13．
\bibitem{Orlov}
D.~Orlov, \emph{Smooth and proper noncommutative schemes and gluing of DG categories},
\href{https://arxiv.org/abs/1402.7364}{arXiv:1402.7364v5} (2015).
使用箇所：Theorems 3.18, 3.25，Corollary 3.19．
\bibitem{BridgelandKleinian}
T.~Bridgeland, \emph{Stability conditions and Kleinian singularities},
Int. Math. Res. Not. \textbf{2009} (2009), no.~21, 4142--4157;
\href{https://arxiv.org/abs/math/0508257}{arXiv:math/0508257v3}.
使用箇所：\S2（ambient圏と内在的導来圏の区別），Lemma 3.5と\S3（単純tiltとWeyl反射）．
\bibitem{BCPW}
N.~Broomhead, R.~Coelho Sim\~oes, D.~Pauksztello, J.~Woolf,
\emph{Simple tilts of length hearts and simple-minded mutation},
Adv. Math. \textbf{503} (2026), Part B, 111197;
\href{https://arxiv.org/abs/2401.02947}{arXiv:2401.02947v3}.
使用箇所：Theorem 2.7およびRemark 2.8（left/right tiltの有限長性）．
\bibitem{EllipticCauchy}
V.~Prokofev, A.~Zabrodin, \emph{Elliptic Cauchy matrices},
\href{https://arxiv.org/abs/2305.02837}{arXiv:2305.02837v1} (2023).
使用箇所：(1.7)，(2.1)（elliptic Cauchy行列と既知のFrobenius determinant）．
\bibitem{PolishchukProj}
A.~Polishchuk, \emph{Noncommutative Proj and coherent algebras},
\href{https://arxiv.org/abs/math/0212182}{arXiv:math/0212182v4} (2004).
使用箇所：Theorem 2.4とその証明（ample sequenceからのabel圏の再構成）．
\bibitem{LiuWangZigzag}
J.~Liu, Z.~Wang, \emph{$A_\infty$-deformations of zigzag algebras via Ginzburg dg algebras},
\href{https://arxiv.org/abs/2301.06757v3}{arXiv:2301.06757v3} (2025).
使用箇所：Theorem 1.1，\S4.3．形式性の先行結果は同論文のEtg\"u--Lekili，Lekili--Uedaへの帰属による．
\bibitem{KellerAinf}
B.~Keller, \emph{Introduction to $A$-infinity algebras and modules},
Homology Homotopy Appl. \textbf{3} (2001), no.~1, 1--35;
\href{https://arxiv.org/abs/math/9910179}{arXiv:math/9910179}.
使用箇所：\S\S6.5, 7.5, 7.7（twisted objectsと拡大閉包）．
\bibitem{Mizuno}
Y.~Mizuno, \emph{Classifying $\tau$-tilting modules over preprojective algebras of Dynkin type},
\href{https://arxiv.org/abs/1304.0667v3}{arXiv:1304.0667v3} (2013).
使用箇所：Dynkin preprojectiveの有限次元性とtorsion classの既知分類．
\bibitem{dTdVV}
L.~de Thanhoffer de V\"olcsey, M.~Van den Bergh,
\emph{Some new examples of non-degenerate quiver potentials},
\href{https://arxiv.org/abs/1005.0808v3}{arXiv:1005.0808v3}.
使用箇所：\S3，Theorem 4.2.1（geometric helixのrolled-up代数に関する既知一般論）．
\bibitem{EOR}
P.~Etingof, A.~Oblomkov, E.~Rains,
\emph{Generalized double affine Hecke algebras of rank 1 and quantized Del Pezzo surfaces},
\href{https://arxiv.org/abs/math/0406480v3}{arXiv:math/0406480v3} (2006).
比較箇所：\S1（quantized del Pezzoと境界），\S11（preprojectiveとの接続）．
\bibitem{EtingofRains}
P.~Etingof, E.~Rains, \emph{Central extensions of preprojective algebras, the quantum Heisenberg algebra, and 2-dimensional complex reflection groups},
\href{https://arxiv.org/abs/math/0503393v3}{arXiv:math/0503393v3} (2006).
使用箇所：\S1, (1)に記載された既知のDynkin preprojectiveのmatrix Hilbert公式．
\bibitem{AiharaMizuno}
T.~Aihara, Y.~Mizuno, \emph{Classifying tilting complexes over preprojective algebras of Dynkin type},
Algebra Number Theory \textbf{11} (2017), no.~6, 1287--1315;
\href{https://doi.org/10.2140/ant.2017.11.1287}{doi:10.2140/ant.2017.11.1287}.
使用箇所：\S3A（自己入射性と既知のNakayama置換）．
\bibitem{Stacks}
The Stacks Project Authors, \emph{The Stacks Project},
\href{https://stacks.math.columbia.edu/tag/00FV}{Tag 00FV}（Nullstellensatz），
\href{https://stacks.math.columbia.edu/tag/00IP}{Tag 00IP}（Krullの共通部分定理）．
\bibitem{DolgachevWalls}
I.~V.~Dolgachev, \emph{Classical Algebraic Geometry: a modern view. II},
author's draft, August 10, 2025,
\href{https://sites.lsa.umich.edu/idolga/wp-content/uploads/sites/1467/2025/10/CAG-2.pdf}{著者公開稿}.
使用箇所：Chapter 8, Theorem 8.3.2, \S8.3.2（弱del Pezzoの反標準収縮）．
\bibitem{SeidelThomasWalls}
P.~Seidel, R.~Thomas, \emph{Braid group actions on derived categories of coherent sheaves},
\href{https://arxiv.org/abs/math/0001043}{arXiv:math/0001043v2} (2000).
使用箇所：Propositions 2.10, 2.12--2.13（spherical twistとbraid関係）．
\bibitem{DrinfeldWalls}
V.~Drinfeld, \emph{DG quotients of DG categories},
\href{https://arxiv.org/abs/math/0210114}{arXiv:math/0210114v7} (2008).
使用箇所：\S\S1.2--1.3, 1.6（dg商，Homのcone公式と普遍性）．
\bibitem{DonovanWemyssWalls}
W.~Donovan, M.~Wemyss, \emph{Noncommutative enhancements of contractions},
Adv. Math. \textbf{344} (2019), 99--136;
\href{https://arxiv.org/abs/1612.01687}{arXiv:1612.01687}.
比較箇所：Introduction, \S\S1.1, 6．与えられた幾何学的収縮に対する理論であり，
本稿の任意の非可換曲面上に収縮射を構成する定理としては用いない．
\end{thebibliography}
\end{document}
